\section[General monodromy data: Painleve-III(D\_6)]{General monodromy data: Painlev\'e-III($\boldsymbol{D_6}$)}
\label{sec:monodromy-rep-$D_6$}

\subsection[Lax pair for Painleve-III(D\_6)]{Lax pair for Painlev\'e-III($\boldsymbol{D_6}$)}
\label{sec:lax-pair-$D_6$}

Following Jimbo and Miwa \cite{MR625446}, we use the fact that each Painlev\'e equation can be recast as an isomonodromic deformation condition for a $2\times 2$ system of linear ODEs with rational coefficients. The case of Painlev\'e-III ($D_6$) corresponds to the situation where the coefficient matrix for the equation in the spectral variable, $\lambda$, has exactly two poles on the Riemann sphere leading to irregular singularities at $\lambda = 0$ and $\lambda = \infty$, at each of which the leading term is diagonalizable. After some normalization, the differential equation can be written in the form\footnote{Henceforth, we use bold capital letters to denote matrices, with the only exceptions being the identity matrix, denoted $\mathbb{I}$ and the Pauli matrices, denoted $\sigma_k$, $k = 1, 2, 3$.}
\begin{equation}
\label{eq:generic-system}
\frac{\partial\mathbf{\Psi}}{\partial\lambda} = \mathbf{\Lambda}^{(6)}(\lambda, x) \mathbf{\Psi}(\lambda, x),
\end{equation}
where
\begin{equation}
\label{eq:lambda-coefficient-D6}
\mathbf{\Lambda}^{(6)}(\lambda, x) = \frac{\ii x}{2}\sigma_3 + \dfrac{1}{2\lambda} \begin{bmatrix} -\Theta_\infty & 2y \\ 2v & \Theta_\infty \end{bmatrix} + \dfrac{1}{2\lambda^2} \begin{bmatrix} \ii x - 2\ii st & 2\ii s \\ -2\ii t(st - x) & -\ii x + 2\ii st\end{bmatrix},
\end{equation}
In this case, the deformation equation is
\begin{equation}
 \frac{\partial\mathbf{\Psi}}{\partial x} = \mathbf{X}(\lambda, x) \mathbf{\Psi}(\lambda, x),
 \label{eq:generic-deformation}
\end{equation}
where
\begin{equation*}
\mathbf{X}(\lambda, x) = \dfrac{\ii \lambda}{2} \sigma_3 + \dfrac{1}{x} \begin{bmatrix} 0 & y \\ v & 0 \end{bmatrix} - \dfrac{1}{2\lambda x} \begin{bmatrix} \ii x - 2\ii st & 2\ii s \\ -2\ii t(st - x) & -\ii x + 2\ii st\end{bmatrix}.
\end{equation*}
In the expressions for $\mathbf{\Lambda}^{(6)}(\lambda, x)$ and $\mathbf{X}(\lambda, x)$, $\Theta_\infty$ is a complex parameter and $s = s(x)$, $t = t(x)$, $ v = v(x)$, $y = y(x)$. The equations \eqref{eq:generic-system} and \eqref{eq:generic-deformation} constitute an over-determined system with compatibility condition
\[
\dpd{\mathbf{\Lambda}^{(6)}}{x}(\lambda, x) - \dpd{\mathbf X}{\lambda}(\lambda, x) + \big[\mathbf{\Lambda}^{(6)}(\lambda, x), \mathbf X(\lambda, x)\big] = \mathbf{0},
\]
where $\big[\mathbf{\Lambda}^{(6)}, \mathbf{X}\big]$ is the commutator. This boils down to the scalar equations
\begin{alignat}{3}
& x \dod{y}{x} = -2xs + \Theta_\infty y, \qquad&&
x \dod{v}{x} = -2xt(st - x) - \Theta_\infty,&\nonumber \\
&x \dod{s}{x} = (1 - \Theta_\infty)s - 2xy + 4yst, \qquad&&
x \dod{t}{x} = \Theta_\infty t - 2yt^2 + 2v.&
\label{eq:scalar-equations}
\end{alignat}
If we let
\begin{equation}
u(x) := - \dfrac{y(x)}{s(x)},
\label{eq:u-from-y-s}
\end{equation}
then it follows from \eqref{eq:scalar-equations} that
\[
x \dod{u}{x} = 2x - (1 - 2\Theta_\infty) u + 4stu^2 - 2xu^2,
\]
which can be seen to be equivalent to \eqref{eq:PIII-$D_6$} by taking another $x$-derivative and using \eqref{eq:scalar-equations} again, after which the quantity
\[
I(x) := \dfrac{2\Theta_\infty }{x} st - \Theta_\infty - \dfrac{2yt}{x}(st - x) + \dfrac{2sv}{x},
\]
appears. However, from \eqref{eq:scalar-equations} it follows that $I'(x) = 0$, so denoting the constant value of $I$ by~$\Theta_0$, we arrive at \eqref{eq:PIII-$D_6$} with parameters
\begin{equation}
 \Theta_0=\frac{\alpha}{4}, \qquad \Theta_\infty=1-\frac{\beta}{4}.
 \label{eq:thetas-alpha-beta}
\end{equation}

The constants $\Theta_0$, $\Theta_\infty$ can be naturally interpreted on the level of the $2 \times 2$ system \eqref{eq:generic-system}, which we now explore. For all the calculations that follow, we assume for simplicity that $x > 0$. The system \eqref{eq:generic-system} admits formal solutions near the singular points\footnote{Here, we use the standard notation $f^{\sigma_3} := \mathrm{diag}\big(f, f^{-1}\big)$.}
\begin{equation}
\mathbf{\Psi}_{\text{formal}}^{(\infty)}(\lambda, x) = \big( \mathbb{I} +\mathbf{\Xi}^{(6)}(x)\lambda^{-1}+ \mathcal{O}\big(\lambda^{-2}\big) \big) \ee^{\ii x \lambda \sigma_3 /2} \lambda^{-\Theta_\infty \sigma_3/2} \qquad \text{as} \quad \lambda \to \infty,
\label{eq:formal-infty}
\end{equation}
and
\begin{equation}
\mathbf{\Psi}^{(0)}_{\text{formal}}(\lambda, x) = \big( \mathbf{\Delta}^{(6)}(x) + \mathcal{O}(\lambda) \big) \ee^{-\ii x \lambda^{-1} \sigma_3 /2} \lambda^{\Theta_0 \sigma_3/2} \qquad \text{as} \quad \lambda \to 0.
\label{eq:formal-zero}
\end{equation}
Here $\mathbf{\Delta}^{(6)}(x)$ is an (invertible) eigenvector matrix of the coefficient of $\lambda^{-2}$ in \eqref{eq:lambda-coefficient-D6}, so the leading term of $\mathbf{\Delta}^{(6)}(x)^{-1}\mathbf{\Lambda}^{(6)}(\lambda,x)\mathbf{\Delta}^{(6)}(x)$ at $\lambda=0$ is diagonal.

For $k = 1, 2,3$, we define the Stokes sectors,
\begin{gather*}
S_k^{(\infty)} = \bigl\{ \lambda \in \C \colon |\lambda| > R, \ k\pi - 2\pi < \arg (\lambda) < k\pi \bigr \},\\
S_k^{(0)} = \bigl \{ \lambda \in \C \colon |\lambda| < r, \ k\pi - 2\pi < \arg (\lambda) < k\pi \bigr\}.
\end{gather*}
 It follows from the classical theory of linear systems that there exist \emph{canonical solutions} \smash{$\mathbf{\Psi}_k^{(\infty)}$} and~\smash{$\mathbf{\Psi}_k^{(0)}$} analytic for \smash{$\lambda\in S_k^{(\infty)}$} and \smash{$\lambda\in S_k^{(0)}$} respectively and determined uniquely by the asymptotic condition
\begin{equation}
\mathbf{\Psi}_k^{(\nu)}(\lambda, x) = \mathbf{\Psi}^{(\nu)}_{\text{formal}}(\lambda, x), \qquad \lambda \in S_k^{(\nu)}, \qquad \nu \in \{0, \infty\}, \qquad k=1,2,3.
\label{eq:psi-asymptotic-condition}
\end{equation}
In these asymptotic conditions, the meaning of the power functions in \eqref{eq:formal-infty} and \eqref{eq:formal-zero} is determined from the range of $\arg(\lambda)$ in the definition of \smash{$S_k^{(\nu)}$}.
The canonical solutions in consecutive Stokes sectors are related to one another by multiplication on the right with \emph{Stokes matrices}, i.e.,
\begin{align}%\label{eq:stokes-condition-2}
\label{eq:stokes-condition-1}\mathbf{\Psi}_2^{(\infty, 0)}(\lambda, x) &= \mathbf{\Psi}_1^{(\infty, 0)}(\lambda, x) \mathbf{S}^{\infty, 0}_{1}, \qquad
\mathbf{\Psi}_3^{(\infty, 0)}(\lambda, x) = \mathbf{\Psi}_2^{(\infty, 0)}( \lambda, x) \mathbf{S}^{\infty, 0}_{2},
\end{align}
where for some \emph{Stokes multipliers} $s_j^{\infty,0}\in\mathbb{C}$, $j=1,2$,
\begin{equation}
\mathbf{S}^{\infty, 0}_{1}= \begin{bmatrix}1&s^{\infty, 0}_1\\0&1\end{bmatrix},\qquad \mathbf{S}^{\infty, 0}_{{2}}=
\begin{bmatrix}1&0\\{s^{\infty, 0}_2}&1\end{bmatrix}.
 \label{eq:Stokes-infty}
\end{equation}
Likewise, by uniqueness and the different interpretation of the multi-valued powers in the formal solutions on the otherwise identical sectors $S^{(\infty,0)}_1$ and $S^{(\infty,0)}_3$, we have the identities
\begin{align}%\label{eq:stokes-condition-3}
\mathbf{\Psi}_3^{(\infty)}(\lambda, x) &= \mathbf{\Psi}_1^{(\infty)}\big(\ee^{-2\pi \ii} \lambda, x\big) e_\infty^{-2\sigma_3},
\qquad
\mathbf{\Psi}_3^{(0)}(\lambda, x) = \mathbf{\Psi}_1^{(0)}\big(\ee^{-2\pi \ii} \lambda, x\big) e_0^{2\sigma_3},\label{eq:stokes-condition-4}
\end{align}
where, combining \eqref{eq:e0-einfty-alpha-beta} with \eqref{eq:thetas-alpha-beta} gives
\begin{equation}
 e_0 = \ee^{\ii \pi \Theta_0/2},\qquad e_\infty = \ee^{\ii \pi \Theta_\infty/2}.
 \label{eq:e0-einfty}
\end{equation}

Canonical solutions in, say $S_k^{(\infty)}$ admit analytic continuation into $S_k^{(0)}$ and since both canonical solutions solve \eqref{eq:generic-system} in the same domain, they must be related by multiplication on the right by a constant \emph{connection matrix}, which we define using
\begin{align}
\label{eq:connection-matrix-1}\mathbf{\Psi}_1^{(0)}(\lambda, x) &= \mathbf{\Psi}_1^{(\infty)}(\lambda, x) \mathbf{C}_{0\infty}^+, \\
\label{eq:connection-matrix-2}\mathbf{\Psi}_2^{(0)}(\lambda, x) &= \mathbf{\Psi}_2^{(\infty)}(\lambda, x) \mathbf{C}_{0\infty}^-.
\end{align}

The condition that the coefficients $y$, $v$, $s$, $t$ in the matrix $\mathbf{\Lambda}^{(6)}(\lambda,x)$ depend on $x$ as a solution of \eqref{eq:scalar-equations} implies simultaneous solvability of \eqref{eq:generic-system} and \eqref{eq:generic-deformation}, and the latter system implies that the Stokes matrices and connection matrices are, like $\Theta_0$ and $\Theta_\infty$, independent of $x$. We show below in Sections~\ref{subsec:cyclic} and \ref{subsec:monodromy-D6} that the four Stokes multipliers and the elements of the two connection matrices are determined from just two essential \emph{monodromy parameters} that we denote by $e_1$ and $e_2$.


\subsection[Riemann--Hilbert problem for Painlev\'e-III(D\_6)]{Riemann--Hilbert problem for Painlev\'e-III($\boldsymbol{D_6}$)} \label{sec:initial-rhp}
Using the canonical solutions, we define the following sectionally-analytic function
\[
\mathbf{\Psi}(\lambda, x) = \begin{cases}
\mathbf \Psi ^{(\infty)}_1(\lambda, x), & |\lambda| >1 \ \text{and} \ \re (\lambda) > 0, \\
\mathbf \Psi ^{(\infty)}_2(\lambda, x), & |\lambda| >1 \ \text{and} \ \re (\lambda) < 0, \\
\mathbf \Psi ^{(0)}_1(\lambda, x), & |\lambda| <1 \ \text{and} \ \re (\lambda) > 0, \\
\mathbf \Psi ^{(0)}_2(\lambda, x), & |\lambda| <1 \ \text{and} \ \re (\lambda) < 0. \\
\end{cases}
\]
Then, it follows from the asymptotic conditions \eqref{eq:psi-asymptotic-condition} and the relations \eqref{eq:stokes-condition-1}--\eqref{eq:stokes-condition-4} and \eqref{eq:connection-matrix-1}--\eqref{eq:connection-matrix-2} that $\mathbf{\Psi}$ solves the following $2 \times 2$ Riemann--Hilbert problem.
Let $\lambda_{\lw}^p$ denote the branch of the power function analytic in $\C\setminus \ii \R_-$ with argument chosen so that
\begin{equation}
-\dfrac{\pi}{2} < \mathrm{arg}_{\lw}(\lambda) < \dfrac{3\pi}{2}.
\label{eq:lw-branch-x-real}
\end{equation}
The notation reminds us that the branch cut of these functions is the contour carrying lower triangular Stokes matrices.
\begin{rhp} \label{rhp:initial}
Fix generic monodromy parameters $( e_1, e_2 )$ determining the Stokes and connection matrices, and $x>0$. We seek a $2 \times 2$ matrix function $\lambda\mapsto\mathbf \Psi(\lambda, x)$ satisfying:
\begin{itemize}\itemsep=0pt
\item Analyticity: $\mathbf \Psi (\lambda, x)$ is analytic in $\C \setminus L^{(6)}$, where $L^{(6)} = \{ | \lambda| = 1\} \cup \ii \R$ is the jump contour shown in Figure {\rm\ref{fig:1}}.
\item  Jump condition: $\mathbf \Psi (\lambda, x)$ has continuous boundary values on $L^{(6)} \setminus \{0\}$ from each component of $\mathbb{C}\setminus L^{(6)}$, which satisfy
$\mathbf \Psi_+ (\lambda, x) = \mathbf \Psi_- (\lambda, x) \mathbf J_{\mathbf \Psi}(\lambda)$, where $\mathbf J_{\mathbf \Psi}(\lambda)$ is as shown in Figure {\rm\ref{fig:1}} and where the $+$ $($resp., $-)$ subscript denotes a boundary value taken from the left $($resp., right$)$ of an arc of $L^{(6)}$.
\item  Normalization: $\mathbf \Psi$ satisfies the asymptotic conditions
\begin{equation}
\label{eq:Psi-asymptotic-infinity}
\boldsymbol \Psi(\lambda, x) = \big( \mathbb{I} +\mathbf{\Xi}^{(6)}(x)\lambda^{-1}+ \mathcal{O}\big(\lambda^{-2}\big) \big) \ee^{\ii x \lambda \sigma_3 /2} \lambda_{\lw}^{-\Theta_\infty \sigma_3/2} \qquad \text{as} \quad \lambda \to \infty,
\end{equation}
and
\begin{equation}
\label{eq:Psi-asymptotic-zero}
\boldsymbol \Psi(\lambda, x) = \big( \mathbf{\Delta}^{(6)}(x) + \mathcal{O}(\lambda) \big) \ee^{-\ii x \lambda^{-1} \sigma_3 /2} \lambda_{\lw}^{\Theta_0 \sigma_3/2} \qquad \text{as} \quad \lambda \to 0,
\end{equation}
where $\mathbf{\Delta}^{(6)}(x)$ is a matrix determined from $\boldsymbol\Psi(\lambda,x)$ having unit determinant.
\end{itemize}
\end{rhp}
\begin{figure}[t]\centering
\includegraphics[scale = 1]{Figures/1c.pdf}
\caption{The jump contour $L^{(6)}$ for $\mathbf{\Psi}(\lambda, x)$ and definition of $\mathbf J_{\mathbf \Psi}(\lambda)$ when $x>0$.}
\label{fig:1}
\end{figure}

Observe that if $\mathbf \Psi$ solves Riemann--Hilbert Problem~\ref{rhp:initial}, then the following limit exists:
\begin{equation}
\label{eq:T-definition}
 \mathbf{\Xi}^{(6)}(x):=\lim_{\lambda\to\infty}\lambda\big[\mathbf{\Psi}(\lambda,x)\ee^{-\ii x\lambda\sigma_3/2}\lambda_{\lw}^{\Theta_\infty\sigma_3/2}-\mathbb{I}\big].
\end{equation}
 Existence of a solution $\mathbf{\Psi}(\lambda, x)$ to Riemann--Hilbert Problem \ref{rhp:initial} which is meromorphic in $x$ on a covering of the plane is well established; see, e.g., \cite[Theorem 5.4]{FIKN}. Furthermore, it follows from the Riemann--Hilbert problem that $\det (\mathbf{\Psi}(\lambda; x)) = 1$; hence $\mathbf{\Xi}^{(6)}(x)$ defined by \eqref{eq:T-definition} has zero trace. The solution of the Painlev\'e-III($D_6$) equation for the initial data that generated the matrices for the inverse monodromy problem is given by
 \begin{equation}
 u(x)=\frac{-\ii \Xi^{(6)}_{12}(x)}{{\Delta}^{(6)}_{11}(x){\Delta}^{(6)}_{12}(x)},
 \label{eq:u-recover-general}
 \end{equation}
where $\mathbf{\Delta}^{(6)}$, $\mathbf{\Xi}^{(6)}$ are as in \eqref{eq:Psi-asymptotic-zero}, \eqref{eq:T-definition}, respectively.


To study the direct monodromy problem and obtain the jump matrices given just the values of $u$ and $u'$ at an initial point $x_0$, it is necessary to introduce artificial initial values of the auxiliary functions $s$, $t$, $v$, $y$ at $x_0$ in way consistent with the definition \eqref{eq:u-from-y-s} of $u(x)$. Different consistent choices lead to different jump matrices, but the jump matrices determine the same function $u(x)$ via \eqref{eq:u-recover-general}. This symmetry is reflected at the level of $\mathbf{\Psi}(\lambda,x)$ by the conjugation
 $\mathbf{\Psi}(\lambda,x)\mapsto \delta^{-\sigma_3}\mathbf{\Psi}(\lambda,x)\delta^{\sigma_3}$
 for any $\delta\neq 0$. Another symmetry that also leaves $u(x)$ invariant but changes the jump matrices $\mathbf{C}_{0\infty}^\pm$ is multiplication of $\mathbf{\Psi}(\lambda,x)$ on the right for $|\lambda|<1$ only by a~unit-determinant diagonal matrix. Therefore, having obtained the jump matrices for the inverse monodromy problem via a direct monodromy calculation, after the fact we may introduce an arbitrary transformation of $\mathbf{\Psi}(\lambda,x)$ of the form
\begin{equation}\label{eq:transf}
 \boldsymbol \Psi(\lambda,x)\mapsto \widetilde{\mathbf{\Psi}}(\lambda,x):=\begin{cases}
 \delta^{-\sigma_3} \boldsymbol \Psi(\lambda,x)\gamma^{\sigma_3},&|\lambda|<1,\\
 \delta^{-\sigma_3}\boldsymbol \Psi(\lambda,x)\delta^{\sigma_3},&|\lambda|>1
 \end{cases}
\end{equation}
without changing $u(x)$. This transformation modifies the Stokes matrices as follows:
 \begin{equation}\label{eq:stokes-transformation}
 \mathbf{S}^\infty_{1,2}\mapsto \widetilde{\mathbf{S}}^\infty_{1,2}:=\delta^{-\sigma_3}\mathbf{S}^\infty_{1,2}\delta^{\sigma_3}\qquad\text{and}\qquad
 \mathbf{S}^0_{1,2}\mapsto \widetilde{\mathbf{S}}^0_{1,2}:=\gamma^{-\sigma_3}\mathbf{S}^0_{1,2}\gamma^{\sigma_3}
 \end{equation}
 and it modifies the connection matrices as
 \begin{equation}
 \mathbf{C}_{0\infty}^\pm\mapsto\widetilde{\mathbf{C}}_{0\infty}^\pm:=\delta^{-\sigma_3}\mathbf{C}_{0\infty}^\pm\gamma^{\sigma_3}.
 \label{eq:connection-modify}
 \end{equation}

\subsection[Monodromy parameters (e\_1, e\_2)]{Monodromy parameters $\boldsymbol{(e_1, e_2)}$}\label{subsec:cyclic}
The cyclic products of the jump matrices for the inverse monodromy problem about the two non-singular self-intersection points of the jump contour $\lambda=\pm\ii$ read
 \begin{gather}
 \text{about } \lambda=+\ii \colon\ (\mathbf{C}_{0\infty}^-)^{-1}(\mathbf{S}^\infty_{1})^{-1}\mathbf{C}_{0\infty}^+\mathbf{S}^0_{1}=\mathbb{I},\nonumber\\
 \text{about } \lambda=-\ii\colon\
 \mathbf{S}_{2}^\infty e_\infty^{2\sigma_3}\mathbf{C}_{0\infty}^+e_0^{2\sigma_3}\big(\mathbf{S}^0_{2}\big)^{-1}(\mathbf{C}_{0\infty}^-)^{-1}=\mathbb{I}.
 \label{eq:cyclic-condition}
 \end{gather}
 We can use the second relation to explicitly write $\mathbf{C}_{0\infty}^+$ in terms of two Stokes matrices and the other connection matrix:
 \begin{equation} \mathbf{C}_{0\infty}^+=e_\infty^{-2\sigma_3}(\mathbf{S}_{2}^\infty )^{-1}\mathbf{C}_{0\infty}^-\mathbf{S}^0_{2}e_0^{-2\sigma_3}.
 \label{eq:right-connection-formula}
 \end{equation}
This identity is an analog of \cite[equation~(3.17)]{Jimbo}.
 Under the condition that $\det \mathbf{C}_{0\infty}^- = 1$, we immediately get that $\det \mathbf{C}_{0\infty}^+ = 1$. Furthermore, using \eqref{eq:right-connection-formula} we eliminate $\mathbf{C}_{0\infty}^+$ from the first equation of \eqref{eq:cyclic-condition} to obtain the identity
 \begin{equation}
 \big(\mathbf{S}^0_{1}\big)^{-1}e_0^{2\sigma_3}\big(\mathbf{S}_{2}^0 \big)^{-1} = \big(\mathbf{C}_{0\infty}^-\big)^{-1}\big(\mathbf{S}^\infty_{1}\big)^{-1}e_\infty^{-2\sigma_3}(\mathbf{S}_{2}^\infty )^{-1}\mathbf{C}_{0\infty}^-.
 \label{eq:left-connection-identity}
 \end{equation}

 In other words, $ (\mathbf{S}^\infty_{1})^{-1}e_\infty^{-2\sigma_3}(\mathbf{S}_{2}^\infty )^{-1}$ and $\big(\mathbf{S}^0_{1}\big)^{-1}e_0^{2\sigma_3}\big(\mathbf{S}_{2}^0 \big)^{-1}$ are similar unit-determinant matrices. Note that this is merely reflective of the fact that both products are monodromy matrices, possibly expressed in terms of different bases of fundamental solutions, for a simple circuit about the origin for solutions of the system \eqref{eq:generic-system}. Let us assume that they have distinct eigenvalues that we will denote $e_1^{\pm 2}$. Then, both products are diagonalizable, so there exist unit-determinant eigenvector matrices $\mathbf{E}^\infty$ and $\mathbf{E}^0$ such that
 \begin{equation}
 (\mathbf{S}^\infty_{1})^{-1}e_\infty^{-2\sigma_3}(\mathbf{S}_{2}^\infty )^{-1}\mathbf{E}^\infty=\mathbf{E}^\infty e_1^{2\sigma_3}\qquad\text{and}\qquad
 \big(\mathbf{S}^0_{1}\big)^{-1}e_0^{2\sigma_3}\big(\mathbf{S}_{2}^0 \big)^{-1}\mathbf{E}^0=\mathbf{E}^0 e_1^{2\sigma_3}.
 \label{eq:Stokes-products-eigenvectors}
 \end{equation}
To specify the eigenvector matrices $\mathbf{E}^\infty$, $\mathbf{E}^0$ uniquely, we agree that their (2,2) entries are both equal to $1$.


 Using \eqref{eq:Stokes-products-eigenvectors} in \eqref{eq:left-connection-identity} gives a homogeneous linear equation on $\mathbf{C}_{0\infty}^-$ that can be written in commutator form as
 \begin{equation*}
 \big[e_1^{2\sigma_3},(\mathbf{E}^\infty)^{-1}\mathbf{C}_{0\infty}^-\mathbf{E}^0\big]=\mathbf{0}.
 \end{equation*}
 The diagonal matrix $e_1^{2\sigma_3}$ can be written in the form
 \begin{equation*}
 e_1^{2\sigma_3}=f\mathbb{I}+g\sigma_3,\qquad f:= \frac{1}{2}\big(e_1^2+e_1^{-2}\big),\qquad g:=\frac{1}{2}\big(e_1^2-e_1^{-2}\big).
 \end{equation*}
 Under the assumption $e_1^4\neq 1$ we already invoked to obtain diagonalizability, $g\neq 0$ so the commutator equation implies that $(\mathbf{E}^\infty)^{-1}\mathbf{C}_{0\infty}^-\mathbf{E}^0$ is a diagonal unit-determinant matrix that we may write in the form $e_2^{\sigma_3}$. Thus we have the identity
 \begin{equation}
 \label{eq:C-zero-infty-diagonalization}
 \mathbf{C}_{0\infty}^-=\mathbf{E}^\infty e_2^{\sigma_3}\big(\mathbf{E}^0\big)^{-1}.
 \end{equation}
 \begin{Remark}\label{remark:change-sign-e2}
 Changing the sign of $e_2$ changes the sign of the connection matrix. This corresponds to multiplication of the solution of Riemann--Hilbert Problem \ref{rhp:initial} by $-1$ inside of unit disc. Looking at formula \eqref{eq:u-recover-general}, we see that solution $u(x)$ does not change after such transformation. Therefore, we can assume $-\frac{\pi}{2}<\arg(e_2)\leq\frac{\pi}{2}$.
 \end{Remark}

 Using \eqref{eq:C-zero-infty-diagonalization} in \eqref{eq:right-connection-formula} then gives the equivalent representations
 \begin{equation}
 \mathbf{C}_{0\infty}^+
 =e_\infty^{-2\sigma_3}(\mathbf{S}_{2}^\infty )^{-1}\mathbf{E}^\infty e_2^{\sigma_3}\big(\mathbf{E}^0\big)^{-1}\mathbf{S}^0_{2}e_0^{-2\sigma_3}
 =\mathbf{S}_{1}^\infty \mathbf{E}^\infty e_2^{\sigma_3}\big(\mathbf{E}^0\big)^{-1}\big(\mathbf{S}^0_{1}\big)^{-1}.
 \label{eq:C-zero-infty-plus-diagonalization}
 \end{equation}

 \subsection{Parametrization of Stokes multipliers and connection matrix}
%\subsection{Parametrization of monodromy manifold}
 Taking the trace of \eqref{eq:left-connection-identity}, we get
\begin{equation*}
%\label{eq:stokes-products}
e_1^2+\frac{1}{e_1^2}=e_\infty^2+\frac1{e_\infty^2}+{s^\infty_1 s^\infty_2}{e_\infty^2}=e_0^2+\frac1{e_0^2}+\frac{s^0_1s^0_2}{e_0^2}.
\end{equation*}
It is clear that one can solve for the products $s^\infty_1s^\infty_2$ and $s_1^0s_2^0$ in terms of $\big(e_1^2,e_\infty^2\big)$ and $\big(e_1^2,e_0^2\big)$ respectively. Using the transformation \eqref{eq:stokes-transformation}, we can take a particular solution of this relation and hence obtain the Stokes multipliers:
\begin{equation}
\label{eq:stokes-parameters}
 s^\infty_1=\frac{e_\infty^2-e_1^2}{e_1^2e_\infty^4},\qquad s^\infty_2=1-e_1^2e_\infty^2,\qquad s^0_1=\frac{e_1^2-e_0^2}{e_1^2},\qquad s^0_2={e_1^2e_0^2 - 1}.
\end{equation}
With the Stokes matrices specified in this way, the eigenvector matrices $\mathbf{E}^\infty$ and $\mathbf{E}^0$ are uniquely specified as mentioned earlier by taking the $(2,2)$ entry to be $1$ in each case, which yields
\begin{gather}
\mathbf{E}^{\infty}=\left[
\begin{matrix}
 \dfrac{e_1^2 \big(e_1^2-e_\infty^2\big)}{e_1^4-1} & -\dfrac{1}{e_1^2e_\infty^2} \vspace{1mm} \\
 \dfrac{e_1^2 e_\infty^2\big(e_1^2 e_\infty^2-1\big)}{e_1^4-1} & 1
\end{matrix}
\right],\label{eq:E-infinity-formula}\\
\mathbf{E}^{0}=\left[
\begin{matrix}
 \dfrac{e_1^2 \big(e_0^2 e_1^2-1\big)}{e_0^2 \big(e_1^4-1\big)} & \dfrac{e_1^2-e_0^2}{e_1^2 \big(e_0^2 e_1^2-1\big)} \vspace{1mm}\\
 \dfrac{e_1^2 \big(1 - e_0^2 e_1^2\big)}{e_0^2 \big(e_1^4-1\big)} & 1
\end{matrix}
\right].\label{eq:E-zero-formula}
\end{gather}
After making such choices, we obtain the formul\ae\ \eqref{eq:x1}--\eqref{eq:x3}. At this point, it can be directly checked that our choices are consistent with the full equation \eqref{eq:left-connection-identity} with $\mathbf C_{0\infty}^-$ given by \eqref{eq:C-zero-infty-diagonalization}.

One can think of fixing the $(2,2)$ entry in the following way: the eigenvector matrices $\mathbf{E}^\infty$ and~$\mathbf{E}^0$ represent ``internal degrees of freedom'' that have an additional symmetry, namely, arbitrary scalings of the eigenvectors that preserve determinants. In other words, while \eqref{eq:transf} induces a conjugation symmetry on the eigenvector matrices, there is an additional symmetry for each involving multiplication on the right by an arbitrary unit-determinant diagonal matrix. Thus, the matrices $\mathbf{E}^\infty$ and $\mathbf{E}^0$ undergo the transformations
 \begin{equation*}
 \mathbf{E}^{\infty}\mapsto \widetilde{\mathbf{E}}^\infty:=\delta^{-\sigma_3}\mathbf{E}^\infty\delta^{\sigma_3}\epsilon_\infty^{\sigma_3}\qquad\text{and}\qquad
 \mathbf{E}^{0}\mapsto \widetilde{\mathbf{E}}^0:=\gamma^{-\sigma_3}\mathbf{E}^0\gamma^{\sigma_3}\epsilon_0^{\sigma_3}
 \end{equation*}
 for some arbitrary nonzero quantities $\epsilon_\infty$, $\epsilon_0$. Note that these transformations along with~\eqref{eq:connection-modify} and $e_2^{\sigma_3}=(\mathbf{E}^\infty)^{-1}\mathbf{C}_{0\infty}^-\mathbf{E}^0$ imply that
 \begin{equation*}
 %\label{eq:e2-gauge-change}
 e_2\mapsto\widetilde{e}_2:=\frac{\gamma\epsilon_0}{\delta\epsilon_\infty}e_2.
 \end{equation*}
 By contrast, $e_1^2\mapsto\widetilde{e}_1^2:=e_1^2$ is a symmetry invariant.

\begin{Remark} \label{rmk:general-arg-x}
 In the case where one is interested in values of $x \in \C$ with $|{\arg} (x)| < \pi$, the analogue of Figure \ref{fig:1} is shown in Figure \ref{fig:rotated-contour}, where the nonsingular self-intersection points are at $\lambda=\pm \ii \ee^{\pm \ii \arg(x)}$ (independent $\pm$ signs).
 \begin{figure}
 \centering
 \includegraphics{Figures/rotated.pdf}
 \caption{The analogue of the contour $L^{(6)}$ in Figure \ref{fig:1} when $|{\arg}(x)| \neq 0$.}
 \label{fig:rotated-contour}
 \end{figure}
 The angles of the rays in the contour $L^{(6)}$ are chosen so that $\ii \lambda x \in \R$ on the rays extending to $\lambda = \infty$, and $\ii \lambda^{-1} x \in \R$ on the rays extending to $\lambda = 0$. Similar to Section~\ref{sec:initial-rhp}, one can formulate a Riemann--Hilbert problem for a sectionally analytic function~$\lambda\mapsto\mathbf \Psi(\lambda, x)$ off of the contour $L^{(6)}$ and one finds \emph{four} connection matrices instead of two, denoted $\mathbf{C}_1$ through~$\mathbf{C}_4$, defined on corresponding arcs of the unit circle as shown in Figure~\ref{fig:rotated-contour}. These satisfy cyclic conditions similar to~\eqref{eq:cyclic-condition}, namely,
 \begin{alignat*}{3}
 & \text{about}\ \lambda=\ii\ee^{-\ii \arg(x)}\colon\ &&\mathbf{C}_1^{-1}\mathbf{S}_1^\infty \mathbf{C}_2= \mathbb{I}, &\\
 & \text{about} \ \lambda=-\ii\ee^{\ii \arg(x)}\colon\ &&\mathbf{C}_1\big(\mathbf{S}_2^0e_0^{-2\sigma_3}\big)^{-1} \mathbf{C}_4^{-1} = \mathbb{I}, &\\
 & \text{about} \ \lambda=-\ii\ee^{-\ii \arg(x)}\colon\ &&\mathbf{C}_3^{-1} \mathbf{S}_2^\infty e_\infty^{2\sigma_3} \mathbf{C}_4 = \mathbb{I},&\\
 & \text{about}\ \lambda=\ii\ee^{\ii \arg(x)}\colon\ &&\mathbf{C}_3 \big(\mathbf S_1^0\big)^{-1} \mathbf C_2^{-1} = \mathbb{I}.&
 \end{alignat*}
 Eliminating all but $\mathbf{C}_3$ from the above identities yields the analog of \eqref{eq:left-connection-identity}, namely,
 \[
 \big(\mathbf S_1^0\big)^{-1} e_0^{2\sigma_3} \big(\mathbf S_2^0\big)^{-1} = \mathbf C_3^{-1} (\mathbf S_1^\infty)^{-1} e_\infty^{-2\sigma_3} (\mathbf S_2^\infty)^{-1} \mathbf C_3.
 \]
 Reasoning similar to that of Section \ref{subsec:cyclic} yields
 \begin{equation}
 \mathbf{C}_3 = \mathbf{E}^\infty e_2^{\sigma_3} \big(\mathbf{E}^0\big)^{-1},
 \label{eq:general-connection-factorization-1}
 \end{equation}
 which, in turn, yields
 \begin{gather}
 \mathbf{C}_1 = \mathbf S_1^\infty \mathbf{E}^\infty e_2^{\sigma_3} \big(\mathbf{E}^0\big)^{-1} \big(\mathbf S_1^0\big)^{-1}, \qquad
 \mathbf{C}_2 = \mathbf{E}^\infty e_2^{\sigma_3} \big(\mathbf{E}^0\big)^{-1} \big(\mathbf S_1^0\big)^{-1},\nonumber \\
 \mathbf{C}_4 = e_\infty^{-2\sigma_3}\big(\mathbf{S}_2^\infty\big)^{-1}\mathbf{E}^\infty e_2^{\sigma_3} \big(\mathbf{E}^0\big)^{-1} = \mathbf S_1^\infty \mathbf E^\infty e_1^{2\sigma_3} e_2^{\sigma_3} (\mathbf{E}^0)^{-1}.
 \label{eq:general-connection-factorization-2}
 \end{gather}
 In this setting, we must adjust our choice of the branch $\lambda\mapsto \mathrm{arg}_{\lw}(\lambda)$, and we choose a branch which satisfies (cf.\ \eqref{eq:lw-branch-x-real} when $\arg(x) = 0$)
 \begin{equation*} % \label{eq:lw-branch-near-infty}
 -\dfrac{\pi}{2} - \arg(x) < \mathrm{arg}_{\lw} (\lambda) < \dfrac{3\pi}{2} - \arg(x), \qquad |\lambda| \to \infty,
 \end{equation*}
 and
 \begin{equation*} % \label{eq:lw-branch-near-zero}
 -\dfrac{\pi}{2} + \arg(x) < \mathrm{arg}_{\lw} (\lambda) < \dfrac{3\pi}{2} + \arg(x), \qquad |\lambda| \to 0.
 \end{equation*}
 A concrete branch cut is chosen later, see Remark \ref{rmk:rotated-lenses} below.
\end{Remark}

\subsection[Example: rational solutions of Painlev\'e-III(D\_6)]{Example: rational solutions of Painlev\'e-III($\boldsymbol{D_6}$)}\label{sec:rational-solutions-parameters}
One can check that the Painlev\'e-III($D_6$) equation with parameters related by $\Theta_0=\Theta_\infty-1$ admits the constant solution $u(x)\equiv 1$. Its monodromy data was calculated in \cite[Section 4]{BMS18} by taking advantage of the fact that the compatible $x$-equation \eqref{eq:generic-deformation} in the Lax pair has simple coefficients. Denoting $m=\Theta_0=\Theta_\infty-1$ gives $e_0^{-2}=\ee^{-\ii\pi m}$ and $e_\infty^2=-\ee^{\ii\pi m}$. Choosing $\gamma$, $\delta$ in~\eqref{eq:transf} satisfying
\[ %\label{eq:rational-gamma-delta}
\delta^2 = \ee^{-2\pi \ii m} \big(1 - \ii \ee^{\pi \ii m}\big)\dfrac{\Gamma\big(\frac{1}{2} - m\big)}{\sqrt{2\pi}} \qquad \text{and} \qquad \gamma^2 = \big(1 + \ii \ee^{\pi \ii m}\big)\dfrac{\Gamma\big(\frac{1}{2} - m\big)}{\sqrt{2\pi}},
\]
one obtains
 \begin{alignat*}{3}
& s^0_1=\frac{\sqrt{2\pi}}{\Gamma\big(\tfrac{1}{2}-m\big)},\qquad&&
 s^0_2=-\ee^{\ii\pi m}\frac{\sqrt{2\pi}}{\Gamma\big(\tfrac{1}{2}+m\big)},&\\
& s^\infty_1=-\frac{\sqrt{2\pi}}{\Gamma\big(\tfrac{1}{2}-m\big)},\qquad&&
 s^\infty_2=\ee^{-\ii\pi m}\frac{\sqrt{2\pi}}{\Gamma\big(\tfrac{1}{2}+m\big)}.&
 \end{alignat*}
 With this choice of $\gamma$, $\delta$ and $\mathbf{E}^\infty$, $\mathbf{E}^0$ chosen as in \eqref{eq:E-infinity-formula}, \eqref{eq:E-zero-formula} (that is, we insist that the $(2,2)$ entry of $\mathbf{E}^\infty$, $\mathbf{E}^0$ is 1 by setting $\epsilon_0 = \epsilon_\infty = 1$), the connection matrices are
 \begin{equation*}
 \mathbf{C}_{0\infty}^+=\begin{bmatrix}1&-\dfrac{\sqrt{2\pi}}{\Gamma\big(\tfrac{1}{2}-m\big)} \\0&1\end{bmatrix},\qquad\mathbf{C}_{0\infty}^-=\begin{bmatrix}1&\dfrac{\sqrt{2\pi}}{\Gamma\big(\tfrac{1}{2}-m\big)} \\0&1\end{bmatrix},
 \end{equation*}
 and
 \begin{equation*}
 e_1^2=\ii \qquad \text{and} \qquad e_2 = 1.
 \end{equation*}
\begin{Remark}
 The above gauge is only needed to match our setup with that of \cite{BMS18}; in the sequel we will be working with $\gamma = \delta = 1$. Formula \eqref{eq:C-zero-infty-diagonalization} then implies
 \[
 e_2^2 = \ee^{-2\pi \ii m}\dfrac{1 - \ii \ee^{\pi \ii m}}{1 + \ii \ee^{\pi \ii m}}.
 \]
 This is important to note when, for example, one tries to verify that \eqref{eq:u-n-leading} below reduces to~\eqref{eq:u-n-rat-leading}.
\end{Remark}

Before beginning to study the large $n$ behavior of $u_n$, we must first establish a similar monodromy representation of the limiting solution of Painlev\'e-III($D_8$), which we do in Section~\ref{sec:monodromy-rep-$D_8$} below.

\subsection{Monodromy manifold}
\label{subsec:monodromy-D6}
It is known that the monodromy manifold for Painlev\'e-III($D_6$) can be given by a cubic equation (see, e.g., \cite{PS}), which can be recovered from our point of view as follows. Denote
\begin{equation*}
%\label{eq:connection-l}
 \mathbf{C}_{0\infty}^-=\begin{bmatrix}
 \ell_1&\ell_2\\\ell_3&\ell_4
 \end{bmatrix}
\end{equation*}
and
\begin{equation*}
%\label{eq:stokes-m}
 \mathbf{S}^0_{2}e_0^{-2\sigma_3}\mathbf{S}_{1}^0 =e_0^{-2} \begin{bmatrix}
 1& s^0_1\\ s^0_2& \left(e_0^4+s^0_1s^0_2 \right)
 \end{bmatrix}=\begin{bmatrix}
 m_1&m_2\\m_3&m_4
 \end{bmatrix}.
\end{equation*}
Then, the inverse of the cyclic relation \eqref{eq:left-connection-identity} allows us to solve for $s_1^\infty$, $s_2^\infty$ in terms of parameters $m_i$, $\ell_i$, and imposes the constraint
\begin{equation}\label{eq:beta}
 e_\infty^2=\ell_1\ell_4m_1-\ell_1\ell_3m_2+\ell_2\ell_4m_3-\ell_2\ell_3m_4.
\end{equation}
Hence, we are left with these eight parameters subject to the constraint \eqref{eq:beta} and the unit-determinant conditions
\begin{equation}\label{eq:det}
 \ell_1\ell_4-\ell_2\ell_3=1,\qquad m_1m_4-m_2m_3=1.
\end{equation}
We may define coordinates which are invariant under the transformation \eqref{eq:transf}:
\begin{equation*}
 I_1:= \ell_1\ell_4,\qquad I_2 := m_2\ell_1\ell_3,\qquad I_3 :=m_3\ell_2\ell_4,\qquad I_4:= m_4, \qquad I_5:= m_1.
 \end{equation*}
Equations \eqref{eq:beta}, \eqref{eq:det} imply
\begin{equation*}
 e_\infty^2 =e_0^{-2} I_1-I_2+I_3-I_4(I_1-1),\qquad I_2I_3-I_1\big(e_0^{-2} I_4-1\big)(I_1-1)=0.
\end{equation*}
We eliminate $I_3$ and get
\begin{gather*}
 -I_1+e_0^{-2} I_1 I_4+I_1^2-e_0^{-2} I_4 I_1^2+e_\infty^2 I_2- I_2I_4 -e_0^{-2}I_1I_2+I_1I_2I_4+I_2^2=0.
\end{gather*}
Introducing new variables
\begin{equation}
x_1:=I_1-1,\qquad x_2:=-e_0^{-2} I_1+I_2,\qquad x_3:=I_4+e_0^{-2}
\label{eq:x_i-def}
\end{equation}
yields the following equation, which defines the \emph{monodromy manifold} for the problem
\begin{equation}
 \label{eq:cubic} x_1x_2x_3+x_1^2+x_2^2+x_2\big(e_0^{-2}+e_\infty^2\big)+x_1\big(1+e^{-2}_0 e^{2}_\infty\big)+e^{-2}_0e^{2}_\infty=0.
\end{equation}
This matches \eqref{eq:cubic-D6} upon using \eqref{eq:thetas-alpha-beta} and \eqref{eq:e0-einfty}. Using \eqref{eq:x_i-def}, we obtain formul\ae\ \eqref{eq:x1}--\eqref{eq:x3}.

To find the singularities of \eqref{eq:cubic}, we adjoin to \eqref{eq:cubic} the three equations obtained by setting to zero the components of the gradient vector of the left-hand side of \eqref{eq:cubic} with respect to $(x_1,x_2,x_3)$. There is therefore at most one singularity:
\begin{alignat}{3}\label{eq:crit1}
 &\text{for}\ e^{-2}_0=e^2_\infty\colon\ &&(x_1,x_2,x_3)=\big(0,-e^{-2}_0,e^{2}_0+e_0^{-2}\big),&\\
 &\text{for}\ e^{2}_0=e_\infty^{2}\colon\ &&(x_1,x_2,x_3)=\bigl(-1,0,e^{2}_0+e_0^{-2}\bigr).&\label{eq:crit2}
\end{alignat}
In particular, if neither $e^{-2}_0=e^2_\infty$ nor $e^{2}_0=e^2_\infty$, then the monodromy manifold is a smooth curve with no singular points.
Notice that we can use $(x_1,x_2)$ as parameters for the generic collection of points on monodromy manifold \eqref{eq:cubic} for which $x_1x_2\neq 0$, because $x_3$ can be explicitly expressed in terms of the other coordinates. The points satisfying \eqref{eq:cubic} with $x_1=0$ form a 1-dimensional variety consisting in general of two distinct lines:
\begin{equation*}
 (x_1,x_2,x_3)=\big(0,-e^{-2}_0,x_3\big)\qquad\text{or}\qquad (x_1,x_2,x_3)=\big(0,-e^2_\infty,x_3\big)
\end{equation*}
each parametrized by $x_3\in\mathbb{C}$. If $e_0^{-2}=e_\infty^2$, the two lines coincide and pass through the critical point \eqref{eq:crit1} of \eqref{eq:cubic}. Likewise there are generally two lines on \eqref{eq:cubic} along which $x_2=0$ each parametrized by $x_3\in\mathbb{C}$:
\begin{equation*}
 (x_1,x_2,x_3)=(-1,0,x_3)\qquad\text{or}\qquad (x_1,x_2,x_3)=\bigl(-e^{-2}_0e^2_\infty,0,x_3\bigr)
\end{equation*}
and if $e^{2}_0=e_\infty^{2}$, the two lines again coincide and pass through the critical point \eqref{eq:crit2} of \eqref{eq:cubic}.

\section[General monodromy data: Painlev\'e-III(D\_8)]{General monodromy data: Painlev\'e-III($\boldsymbol{D_8}$)} \label{sec:monodromy-rep-$D_8$}

\subsection[Lax pair for Painlev\'e-III(D\_8)]{Lax pair for Painlev\'e-III($\boldsymbol{D_8}$)}
\label{sec:lax-pair-$D_8$}

The Painlev\'e-III($D_8$) equation \eqref{eq:PIII-$D_8$} can also be formulated as an isomonodromic deformation of a linear system. In this case we need two ramified irregular singularities at $\lambda = 0$ and $\lambda = \infty$, i.e., we consider the system
\begin{align}\label{eq:lax_pair_D8}
\frac{\partial\mathbf{\Omega}}{\partial\lambda} (\lambda, z) &= {\mathbf{\Lambda}}^{(8)}(\lambda, z) \mathbf{\Omega}(\lambda, z), \\
\frac{\partial\mathbf{\Omega}}{\partial z} (\lambda, z) &= {\mathbf{Z}}(\lambda, z) \mathbf{\Omega}(\lambda, z),
\label{eq:lax_pair_D8_2}
\end{align}
where
\begin{gather*}
{\mathbf{\Lambda}}^{(8)}(\lambda, z) = \begin{bmatrix}0&{\ii z}{} \\0&0\end{bmatrix} +\dfrac{1}{4\lambda}\begin{bmatrix}V(z) & {W}(z)\\2 & -V(z)\end{bmatrix} + \dfrac{1}{\lambda^2}\begin{bmatrix}X(z) & -2\ii X(z)^2U(z)\\-\ii/(2U(z)) & -X(z)\end{bmatrix},
\\
{\mathbf{Z}}(\lambda, z) = \lambda\begin{bmatrix}0& {\ii} \\0&0\end{bmatrix} + \frac{1}{4z}\begin{bmatrix}V(z) & W(z)\\2 & -V(z)\end{bmatrix} - \frac{1}{z\lambda}\begin{bmatrix}X(z) & -2\ii X(z)^2U(z)\\-\ii/(2U(z)) & -X(z)\end{bmatrix},
\end{gather*}
and functions $U(z)$, $V(z)$, $W(z)$, $X(z)$ satisfy the identities
\begin{gather}
 W(z) +4zU(z) +4\ii U(z)V(z)X(z) + 8 U(z )^2X(z)^2 =0,
 \label{eq:UVWXidentity1}\\
 U(z) V(z)^2-4 U(z) V(z)+2 U(z) W(z)+3 U(z)+8 z=0.\label{eq:UVWidentity1}
\end{gather}
Note the characteristic feature that the leading terms of $\mathbf{\Lambda}^{(8)}(\lambda,z)$ and of $\mathbf{Z}(\lambda,x)$ at the singular points $\lambda=0,\infty$ are singular and nondiagonalizable matrices.

Since $\mathbf{\Omega}(\lambda,z)$ is a simultaneous fundamental solution matrix of the Lax system \eqref{eq:lax_pair_D8}--\eqref{eq:lax_pair_D8_2}, the zero-curvature compatibility condition for that system is therefore satisfied:
\begin{gather*}
 \frac{\partial\mathbf{\Lambda}^{(8)}}{\partial z}(\lambda,z)-\frac{\partial\mathbf{Z}}{\partial\lambda}(\lambda,z)+\big[\mathbf{\Lambda}^{(8)}(\lambda,z),\mathbf{Z}(\lambda,z)\big]=\mathbf{0}.
\end{gather*}
Equating to zero the coefficients of different powers of $\lambda$ on the left-hand side gives a first-order system of four differential equations on the four functions $U(z)$, $V(z)$, $W(z)$, and $X(z)$:
\begin{gather}
 zU'(z)=V(z)U(z)-U(z)-4\ii X(z)U(z)^2,\qquad
 V'(z)=\frac{4}{U(z)},\nonumber\\
 W'(z)=-16\ii X(z),\qquad
 zX'(z)=X(z)+2\ii U(z)X(z)^2-\frac{\ii W(z)}{4U(z)}.
 \label{eq:PossiblyD8system}
\end{gather}
It is possible to express the functions $W(z)$, $X(z)$, and $V(z)$ in terms of $U(z)$ and $U'(z)$ using \eqref{eq:UVWXidentity1}, \eqref{eq:UVWidentity1}, and \eqref{eq:PossiblyD8system}, but since we do not use these formul\ae, we do not present them here. Using \eqref{eq:PossiblyD8system} to repeatedly eliminate all derivatives, it is straightforward to obtain the following identity:
\begin{gather*}
 U''(z)-\frac{U'(z)^2}{U(z)}+\frac{U'(z)}{z} -\frac{4U(z)^2+4}{z}\\
 \qquad=-\frac{U(z)}{z^2}\big[ W(z)+4 z U(z) +4 \ii U(z) V(z) X(z) + 8 U(z)^2 X(z)^2\big].
\end{gather*}
Of course the right-hand side vanishes as a result of the identity \eqref{eq:UVWXidentity1}. Hence $U(z)$ is a solution of \eqref{eq:PIII-$D_8$}, the Painlev\'e-III($D_8$) equation.

For all the calculations that follow, we assume for simplicity that $z > 0$. The system \eqref{eq:lax_pair_D8} admits formal solutions near the singular points
\begin{equation}\label{eq:Omega-expand-infty-early}
\mathbf{\Omega}_{\text{formal}}^{(\infty)}(\lambda,z) = \bigg( \mathbb{I} + \frac{\mathbf{\Xi}^{(8)}(z)}{\lambda}+\mathcal{O}\big(\lambda^{-2}\big) \bigg) \rho_\infty^{\sigma_3/2}
 {\dfrac{1}{\sqrt{2}}} \begin{bmatrix} \ii & -1\\ 1 &-\ii \end{bmatrix} \ee^{\ii \rho_\infty \sigma_3} \qquad \text{as} \quad \lambda \to \infty,
\end{equation}
and
\begin{align}
\mathbf{\Omega}^{(0)}_{\text{formal}}(\lambda, z) ={}& {\mathbf{\Delta}}^{(8)}(z)\big( \mathbb{I} +\mathbf{\Pi}(z)\lambda+ \mathcal{O}\big(\lambda^2\big) \big)\nonumber\\
&\times\rho_0^{\sigma_3/2} \ee^{-\pi \ii \sigma_3/4}{\dfrac{1}{\sqrt{2}}} {}{} \begin{bmatrix} \ii & -1 \\ 1 &-\ii \end{bmatrix} \ee^{ \rho_0 \sigma_3} \qquad \text{as} \quad \lambda \to 0,\label{eq:Omega-expand-zero-early}
\end{align}
where
\begin{equation*}
\rho_\infty=\sqrt{-2 \ii z\lambda},\qquad \rho_0= \sqrt{2 \ii z\lambda^{-1}}
\end{equation*}
 and the square roots denote principal branches.
The function ${\mathbf{\Delta}}^{(8)}(z)$ satisfies the identity
\begin{equation}
{\mathbf{\Delta}}^{(8)}(z)\begin{bmatrix}0&-\ii z\\0&0\end{bmatrix}{\mathbf{\Delta}}^{(8)}(z)^{-1}=\begin{bmatrix}X(z) & -2\ii X(z)^2U(z)\\-\ii/(2U(z)) & -X(z)\end{bmatrix},
\label{eq:Delta-U-X-identity}
\end{equation}
and hence the solution $U(z)$ can be expressed as
\begin{equation}\label{eq:un-U-early}
U(z):=-\frac{1}{2z\Delta^{(8)}_{21}(z)^2}.
\end{equation}
For $k = 1, 2$, we define the Stokes sectors,
\begin{align*}
&\mathcal S_k^{(\infty)}= \left \{ \lambda \in \C \colon  |\lambda| > R, \, 2\pi k - \frac{7\pi}{2} < \arg (\lambda) < 2\pi k+\frac{\pi}{2} \right \},\\
&\mathcal S_k^{(0)}= \left \{ \lambda \in \C \colon  |\lambda| < r, \, 2\pi k - \frac{5\pi}{2} < \arg (\lambda) < 2\pi k+\frac{3\pi}{2}\right \}.
\end{align*}

It follows from the classical theory of linear systems that there exist canonical solutions \smash{$\mathbf{\Omega}_k^{(\infty)}$},~\smash{$\mathbf{\Omega}_k^{(0)}$} determined uniquely by the asymptotic condition
\begin{equation}
\mathbf{\Omega}_k^{(\nu)}(\lambda, z) = \mathbf{\Omega}^{(\nu)}_{\text{formal}}(\lambda, z), \qquad \lambda \in \mathcal S_k^{(\nu)}, \quad \nu \in \{0, \infty\}.
\label{eq:psi-asymptotic-condition-D8}
\end{equation}
The canonical solutions in consecutive Stokes sectors at $\lambda=0,\infty$ are related to one another by multiplications on the right with Stokes matrices, i.e.,
\begin{gather}%\label{eq:stokes-condition-D8-2}\label{eq:stokes-condition-D8-3}
\label{eq:stokes-condition-D8-1}\mathbf{\Omega}_2^{(\infty)}(\lambda, z) = \mathbf{\Omega}_1^{(\infty)}(\lambda,z) \mathbf{S}^{\infty}_{1},\qquad \lambda\in \mathcal S_1^{(\infty)}\cap \mathcal S_2^{(\infty)},\\
\mathbf{\Omega}_1^{(0)}(\lambda, z) = \mathbf{\Omega}_0^{(0)}(\lambda, z) \mathbf{S}^{0}_{0}, \qquad \lambda\in \mathcal S_0^{(0)}\cap \mathcal S_1^{(0)}, \\
\mathbf{\Omega}_{2}^{(\infty)}(\lambda, z) = \mathbf{\Omega}_1^{(\infty)}\big(\ee^{-2\pi \ii} \lambda, z\big) (-\ii\sigma_2), \\
\label{eq:stokes-condition-D8-4}\mathbf{\Omega}_{1}^{(0)}(\lambda, z) = \mathbf{\Omega}_0^{(0)}\big(\ee^{-2\pi \ii} \lambda, z\big) (\ii\sigma_2),
\end{gather}
where
\begin{equation}
\mathbf{S}^\infty_{1}= \begin{bmatrix}1&t^\infty_1\\0&1\end{bmatrix},\qquad \mathbf{S}^0_{{0}}=
 \begin{bmatrix}1&0\\{t^0_0}&1\end{bmatrix} .
 \label{eq:Stokes-zero}
\end{equation}
Canonical solutions in, say, $\mathcal S_k^{(\infty)}$ admit analytic continuation into $\mathcal S_k^{(0)}$ and since both canonical solutions solve \eqref{eq:generic-system} in the same domain, they must be related by multiplication on the right by a constant connection matrix, which we define using
\begin{align}
\mathbf{\Omega}_0^{(0)}(\lambda, z) &= \mathbf{\Omega}_1^{(\infty)}(\lambda, z) \mathbf{C}_{0\infty}.
\label{eq:connection-matrix-D8}
\end{align}

\subsection[Riemann--Hilbert problem for Painlev\'e-III(D\_8)]{Riemann--Hilbert problem for Painlev\'e-III($\boldsymbol{D_8}$)} In a fashion similar to Section \ref{sec:initial-rhp}, we now formulate a $2\times 2$ Riemann--Hilbert problem for a~sectionally-analytic function ${\mathbf{\Omega}}$ defined by
\[
{\mathbf{\Omega}}(\lambda, z) =
 \begin{cases}
\mathbf \Omega_1^{(\infty)}(\lambda, z) , & |\lambda| >1 \ \text{and} \ -\dfrac{\pi}{2} < \arg (\lambda) < \dfrac{3\pi}{2}, \vspace{1mm}\\
\mathbf \Omega_0^{(0)}(\lambda, z) , & |\lambda| <1 \ \text{and} \ -\dfrac{\pi}{2} < \arg (\lambda) < \dfrac{3\pi}{2}.
\end{cases}
\]
Then, it follows from the asymptotic conditions \eqref{eq:psi-asymptotic-condition-D8} and the relations \eqref{eq:stokes-condition-D8-1}--\eqref{eq:stokes-condition-D8-4} and \eqref{eq:connection-matrix-D8} that $\mathbf{\Omega}$ solves the following $2 \times 2$ Riemann--Hilbert problem.
\begin{rhp} \label{rhp:D8}
Fix monodromy data $\big( t_0^0, t_1^\infty \big)$ and $z>0$. We seek a $2 \times 2$ matrix function $\lambda\mapsto{\mathbf{\Omega}}(\lambda, z)$ satisfying:
\begin{itemize}\itemsep=0pt
\item Analyticity: ${\mathbf{\Omega}} (\lambda, z)$ is analytic in $\C \setminus L^{(8)}$, where $L^{(8)} = \{ | \lambda| = 1\} \cup \ii \R_-$ is the jump contour shown in Figure {\rm\ref{fig:D8}}.
\item Jump condition: ${\mathbf{\Omega}} (\lambda, z)$ has continuous boundary values on $L^{(8)} \setminus \{0\}$ from each component of $\mathbb{C}\setminus L^{(8)}$, which satisfy
\[{\mathbf{\Omega}}_+ (\lambda, z) = {\mathbf{\Omega}}_- (\lambda, z) \mathbf J_{{\mathbf{\Omega}}}(\lambda),\] where $\mathbf J_{{\mathbf{\Omega}}}(\lambda)$ is as shown in Figure {\rm \ref{fig:D8}}.
\item Normalization: ${\mathbf{\Omega}}(\lambda,z)$ satisfies the asymptotic conditions
\begin{equation*}
%\label{eq:Omega-asymptotic-infinity}\nonumber
{\mathbf{\Omega}}(\lambda, z) = \big( \mathbb{I} + \mathcal{O}\big(\lambda^{-1}\big) \big) \rho_\infty^{\sigma_3/2}
 {\dfrac{1}{\sqrt{2}}} \begin{bmatrix} \ii & -1\\ 1 &-\ii \end{bmatrix} \ee^{\ii \rho_\infty \sigma_3} \qquad \text{as} \quad \lambda \to \infty,
\end{equation*}
and
\begin{equation*}
%\label{eq:Omega-asymptotic-zero}
{\mathbf{\Omega}}(\lambda, z) = \big( {\mathbf{\Delta}}^{(8)}(z) + \mathcal{O}(\lambda) \big)\rho_0^{\sigma_3/2} {\dfrac{\ee^{-\frac{\ii\pi \sigma_3} {4}}}{\sqrt{2}}} {}{} \begin{bmatrix} \ii & -1\\ 1 &-\ii \end{bmatrix} \ee^{ \rho_0 \sigma_3} \qquad \text{as} \quad \lambda \to 0,
\end{equation*}
where ${\mathbf{\Delta}}^{(8)}(z)$ is a matrix determined from $\mathbf{\Omega}(\lambda,z)$ having unit determinant.
\end{itemize}
\end{rhp}
Solvability of Riemann--Hilbert Problem \ref{rhp:D8} is discussed in Section~\ref{sec:suleimanov-solution-connection}.

\begin{figure}[t]\centering
\includegraphics[scale = 1]{Figures/1d.pdf}
\caption{The jump contour $L^{(8)}$ and definition of $\mathbf J_{{\mathbf{\Omega}}}(\lambda)$ when $z>0$.}\label{fig:D8}
\end{figure}

 \subsection[Lax pair equations for Omega (lambda,z)]{Lax pair equations for $\boldsymbol{\Omega(\lambda,z)}$}\label{sec:Lax-pair}

Since the jump matrices depend on neither $\lambda$ nor $z$, the matrices
\begin{equation}
 \mathbf{\Lambda}^{(8)}(\lambda,z):=\frac{\partial\mathbf{\Omega}}{\partial\lambda}(\lambda,z)\mathbf{\Omega}(\lambda,z)^{-1}\qquad\text{and}\qquad\mathbf{Z}(\lambda,z):=\frac{\partial\mathbf{\Omega}}{\partial z}(\lambda,z)\mathbf{\Omega}(\lambda,z)^{-1}
 \label{eq:Lambda-Z-def}
\end{equation}
are both analytic functions of $\lambda$ in the domain $\mathbb{C}\setminus\{0\}$.
We determine these analytic functions by computing sufficiently many terms in their asymptotic expansions as $\lambda\to\infty$ and $\lambda\to 0$ using \eqref{eq:Omega-expand-infty-early}--\eqref{eq:Omega-expand-zero-early}.
We will use the identities
\begin{equation}
 \frac{\partial\rho_\infty}{\partial\lambda} = -\ii z\rho_\infty^{-1}\qquad\text{and}\qquad\frac{\partial\rho_\infty}{\partial z} = -\ii\lambda\rho_\infty^{-1}
 \label{eq:rhoinfty-derivs}
\end{equation}
and
\begin{equation}
 \frac{\partial\rho_0}{\partial\lambda}=-\ii z\lambda^{-2}\rho_0^{-1}\qquad\text{and}\qquad\frac{\partial\rho_0}{\partial z} = \ii\lambda^{-1}\rho_0^{-1}.
 \label{eq:rhozero-derivs}
\end{equation}
Using \eqref{eq:rhoinfty-derivs} and \eqref{eq:Omega-expand-infty-early} gives, in the limit $\lambda\to\infty$, the expansions
\begin{gather}
 \mathbf{\Lambda}^{(8)}(\lambda,z)= \begin{bmatrix}0&{\ii z}{} \\0&0\end{bmatrix}
 +\frac{1}{4\lambda}\begin{bmatrix}1-{4\ii z}\Xi_{21}^{(8)}(z) & {4\ii z}{}(\Xi_{11}^{(8)}(z)-\Xi_{22}^{(8)}(z)) \vspace{1mm}\\ 2 &-1 +{4\ii z}\Xi_{21}^{(8)}(z)\end{bmatrix}+\mathcal{O}\big(\lambda^{-2}\big),\\
 \mathbf{Z}(\lambda,z)=
 \lambda\begin{bmatrix}0 & {\ii}{} \\0&0\end{bmatrix}
 +\frac{1}{4z}\begin{bmatrix}
 1-{4\ii z}\Xi_{21}^{(8)}(z) &{4\ii z}{}(\Xi_{11}^{(8)}(z)-\Xi_{22}^{(8)}(z)) \vspace{1mm}\\2 & -1+{4\ii z}\Xi_{21}^{(8)}(z)
 \end{bmatrix}
 +\mathcal{O}\big(\lambda^{-1}\big).
\label{eq:Lax-coefficients-infty}
\end{gather}
Actually, we can also go to higher order and compute the coefficient of $\lambda^{-2}$ in the matrix element~$\Lambda_{21}(\lambda,z)$,
in the limit $\lambda\to\infty$:
\begin{align}
\Lambda_{21}^{(8)}(\lambda,z)=\frac{1}{2\lambda} +\frac{1}{2\lambda^2}\bigl(-\Xi_{21}^{(8)}(z)-{2\ii z}{}\Xi_{21}^{(8)}(z)^2-\Xi_{11}^{(8)}(z)+\Xi_{22}^{(8)}(z)\bigr)
+\mathcal{O}\big(\lambda^{-3}\big).
 \label{eq:Lambda21-higher-order}
\end{align}
Likewise, using \eqref{eq:rhozero-derivs} and \eqref{eq:Omega-expand-zero-early} gives that as $\lambda \to 0$
\begin{gather}
 \mathbf{\Lambda}^{(8)}(\lambda,z)=\mathbf{\Delta}^{(8)}(z)\Bigg(\frac{1}{\lambda^2}\begin{bmatrix}0&{-\ii z}{} \\0&0\end{bmatrix}\nonumber \\ \phantom{\mathbf{\Lambda}^{(8)}(\lambda,z)=}{}-\frac{1}{4\lambda}\begin{bmatrix}1-{4\ii z}\Pi_{21}(z)&{4\ii z}{}(\Pi_{11}(z)-\Pi_{22}(z))\\2&-1+{4\ii z}\Pi_{21}(z)\end{bmatrix}\Bigg)\mathbf{\Delta}^{(8)}(z)^{-1}
+ \mathcal{O}(1),\nonumber
 \\
 \mathbf{Z}(\lambda,z)=\frac{1}{\lambda}\mathbf{\Delta}^{(8)}(z)\begin{bmatrix}0&{\ii}{} \\0&0\end{bmatrix}\mathbf{\Delta}^{(8)}(z)^{-1} + \mathcal{O}(1).
\label{eq:Lax-coefficients-zero}
\end{gather}
Applying Liouville's theorem yields the exact expressions
\begin{align}
 & \mathbf{\Lambda}^{(8)}(\lambda,z)= \begin{bmatrix}0&{\ii z}{} \\0&0\end{bmatrix} +\frac{1}{4\lambda}\begin{bmatrix}1-{4\ii z}\Xi^{(8)}_{21}(z) & {4\ii z}{}\big(\Xi^{(8)}_{11}(z)-\Xi^{(8)}_{22}(z)\big) \\ 2 &-1 +{4\ii z}{}\Xi^{(8)}_{21}(z)\end{bmatrix}\nonumber
 \\
 &\hphantom{\mathbf{\Lambda}^{(8)}(\lambda,z)= }{}
 +\frac{\ii z}{\lambda^2}\begin{bmatrix}\Delta^{(8)}_{11}(z)\Delta^{(8)}_{21}(z) & -\Delta^{(8)}_{11}(z)^2 \\
 \Delta^{(8)}_{21}(z)^2 & -\Delta^{(8)}_{11}(z)\Delta^{(8)}_{21}(z)\end{bmatrix},
\label{eq:Lambda-exact}
\end{align}
and
\begin{align*}
 & \mathbf{Z}(\lambda,z)= \lambda\begin{bmatrix}0 & {\ii}{} \\0&0\end{bmatrix}+\frac{1}{4z}\begin{bmatrix}
 1-{4\ii z}{}\Xi^{(8)}_{21}(z) &{4\ii z}{}\big(\Xi^{(8)}_{11}(z)-\Xi^{(8)}_{22}(z)\big) \\2 & -1+{4\ii z}{}\Xi^{(8)}_{21}(z)
 \end{bmatrix} \\
 &\hphantom{\mathbf{Z}(\lambda,z)=}{} -
 \frac{\ii }{ \lambda}\begin{bmatrix}\Delta^{(8)}_{11}(z)\Delta^{(8)}_{21}(z) & -\Delta^{(8)}_{11}(z)^2 \\
 \Delta^{(8)}_{21}(z)^2 & -\Delta^{(8)}_{11}(z)\Delta^{(8)}_{21}(z)\end{bmatrix}.
\end{align*}
Using the notation \eqref{eq:un-U-early} and noting the structure of the coefficients of the different powers of $\lambda$, it is convenient to reparametrize the coefficients as follows:
\begin{equation*}
 \mathbf{\Lambda}^{(8)}(\lambda,z)=\begin{bmatrix}0&{\ii z}{} \\0&0\end{bmatrix} +\frac{1}{4\lambda}\begin{bmatrix}V(z) & W(z)\\2 & -V(z)\end{bmatrix} +\frac{1}{\lambda^2}\begin{bmatrix} X(z) & -2 \ii X(z)^2U(z)\\-\ii/(2U(z)) & -X(z)\end{bmatrix}
 % \label{eq:Lambda-exact-reparametrize}\nonumber
\end{equation*}
and
\begin{equation*} \mathbf{Z}(\lambda,z)=\lambda\begin{bmatrix}0&{\ii}{} \\0&0\end{bmatrix} + \frac{1}{4z}\begin{bmatrix}V(z) & W(z)\\2 & -V(z)\end{bmatrix}-\frac{1}{z\lambda}\begin{bmatrix}X(z) & -2\ii X(z)^2U(z)\\-\ii/(2U(z)) & -X(z)\end{bmatrix}.
\end{equation*}
The quantities $U(z)$, $V(z)$, $W(z)$, and $X(z)$ are not independent; comparing the $21$-element of the coefficient of $\lambda^{-1}$ in the expansion of $\mathbf{\Delta}^{(8)}(z)^{-1}\mathbf{\Lambda}^{(8)}(\lambda,z)\mathbf{\Delta}^{(8)}(z)$ computed using \eqref{eq:Lax-coefficients-infty} and~\eqref{eq:Lax-coefficients-zero} gives the identity \eqref{eq:UVWXidentity1}.
At the same time from formula \eqref{eq:Lambda21-higher-order} we get identity \eqref{eq:UVWidentity1}.


Since \eqref{eq:Lambda-Z-def} holds for the same matrix function $\mathbf{\Omega}(\lambda,z)$, the latter satisfies the equations of a compatible Lax system
\begin{equation}
 \frac{\partial\mathbf{\Omega}}{\partial\lambda}(\lambda,z)=\mathbf{\Lambda}^{(8)}(\lambda,z)\mathbf{\Omega}(\lambda,z)\qquad\text{and}\qquad\frac{\partial\mathbf{\Omega}}{\partial z}(\lambda,z)=\mathbf{Z}(\lambda,z)\mathbf{\Omega}(\lambda,z),
\label{eq:Lax-system}
\end{equation}
which coincides with the system \eqref{eq:lax_pair_D8}--\eqref{eq:lax_pair_D8_2}.

\subsection{Monodromy manifold}
\label{sec:monodromy-rep-$D_8$-manifold}
Introducing notation for the connection matrix elements
\begin{equation*}
 \mathbf{C}_{0\infty}=\begin{bmatrix}
 n_1&n_2\\n_3&n_4
 \end{bmatrix},\qquad \det(\mathbf{C}_{0\infty})=1,
\end{equation*}
we have the cyclic relation around the unique nonsingular point of self-intersection of $L^{(8)}$
 \begin{equation*}%\label{eq:cyclic}\nonumber
\mathbf{S}^\infty_{{1}}\ii\sigma_2= \mathbf{C}_{0\infty}\mathbf{S}^0_{{0}}(-\ii\sigma_2)(\mathbf{C}_{0\infty})^{-1},
\end{equation*}
which implies
\[
n_1=-n_4,\qquad t_1^{\infty}=t_{0}^0,\qquad n_3=n_2-n_4t_1^{\infty}.
\]
Denoting
\[
y_1=n_3,\qquad y_2=n_4,\qquad y_3=t_1^{\infty},
\]
the condition $\det(\mathbf{C}_{0\infty})=1$ implies that the coordinates $(y_1,y_2,y_3)$ are related by the cubic equation \eqref{eq:cubic-D8}.

\begin{Remark}
 If the solution $\mathbf{\Omega}(\lambda,z)$ is multiplied by the scalar $-1$ for $|\lambda|<1$ and left unchanged for $|\lambda|>1$, then the elements of the connection matrix $\mathbf{C}_{0\infty}$ change sign while the Stokes multiplier $t_1^\infty$ is invariant. Therefore, this transformation changes $(y_1,y_2,y_3)$ to $(-y_1,-y_2,y_3)$, yielding a different point on the cubic \eqref{eq:cubic-D8}. The matrix coefficient $\mathbf{\Delta}^{(8)}(z)$ also changes sign, however \smash{$\Delta_{21}^{(8)}(z)^2$} is invariant, so the solution $U(z)$ of the Painlev\'e-III($D_8$) equation \eqref{eq:PIII-$D_8$} is the same for both points.
 \label{rem:minus-y1-y2}
\end{Remark}

\section{Schlesinger transformation and proof of Proposition~\ref{prop:schlesinger}}\label{sec:schlesinger}

Fix generic monodromy parameters $(e_1, e_2)$. In view of the parametrization of the Stokes multipliers in \eqref{eq:stokes-parameters} and the eigenvector matrices in \eqref{eq:E-infinity-formula}, \eqref{eq:E-zero-formula}, this data determines from Riemann--Hilbert Problem~\ref{rhp:initial} a matrix $\mathbf{\Psi}(\lambda, x)$ which is meromorphic in $x$ and satisfies asymptotic conditions \eqref{eq:Psi-asymptotic-infinity} and \eqref{eq:Psi-asymptotic-zero}, which we write in the form\footnote{The coefficients $\mathbf{\Psi}_j^{\infty}$, $\mathbf{\Psi}_j^{0}$ should not be confused with the fundamental solutions discussed in the previous sections. The reader can rest assured that this notation will only appear in this section.}
\begin{align*}
&\mathbf{\Psi}(\lambda, x) \lambda_{\lw}^{\Theta_\infty \sigma_3 /2} \ee^{-\ii x \lambda \sigma_3 /2} = \mathbb{I} + \mathbf{\Psi}_1^{\infty}(x) \lambda^{-1} + \mathcal{O}\big(\lambda^{-2}\big) , \qquad \lambda \to \infty, \\
&\mathbf{\Psi}(\lambda, x) \lambda_{\lw}^{-\Theta_0 \sigma_3 /2} \ee^{\ii x \lambda^{-1} \sigma_3 /2} = \mathbf{\Psi}_0^{0}(x) + \mathbf{\Psi}_1^{0}(x) \lambda + \mathcal{O}\big(\lambda^{2}\big) , \qquad \lambda \to 0.
\end{align*}
Define the matrices
\begin{equation*}%\label{eq:sigma+-}
\sigma_+ = \begin{bmatrix} 1 & 0 \\ 0 & 0 \end{bmatrix} \qquad \text{and} \qquad \sigma_- = \begin{bmatrix} 0 & 0 \\ 0 & 1 \end{bmatrix}.
\end{equation*}
Following \cite{BMS18}, assuming the $(1, 1)$ entry of $\mathbf{\Psi}_0^{0}(x)$, denoted ${\Psi}_{0, 11}^{0}(x)$, is not identically zero, we consider the Schlesinger transformation
\begin{equation*}
%\label{eq:Schlesinger-up}
\hat{\mathbf{\Psi}}(\lambda, x) := \big( \sigma_+ \lambda_{\lw}^{1/2} + \hat{\mathbf{S}}(x) \lambda_{\lw}^{-1/2}\big) \mathbf{\Psi}(\lambda, x),
\end{equation*}
where
\begin{equation*}
\hat{\mathbf{S}} (x):= \begin{bmatrix}
\Psi_{0, 21}^0(x) \Psi_{1, 12}^\infty(x)/\Psi^0_{0,11}(x) & -\Psi_{1, 12}^\infty(x) \vspace{1mm}\\ -\Psi_{0, 21}^0(x)/\Psi_{0, 11}^0(x) & 1
\end{bmatrix}.
\end{equation*}
Since $\lambda_{\lw}^{\pm 1/2}$ has its branch cut along part of the curve $L^{(6)}$, we see that $\hat{\mathbf{\Psi}}(\lambda, x)$ is analytic in $\C \setminus L^{(6)}$ and, by direct calculation, has the jumps on \smash{$L^{(6)}$} summarized by Figure \ref{fig:1}, with the exception of the sign changes
\[
\mathbf{S}_2^0 e_0^{-2\sigma_3} \mapsto -\mathbf{S}_2^0 e_0^{-2\sigma_3} \qquad \text{and} \qquad \mathbf{S}_2^\infty e_\infty^{2\sigma_3} \mapsto -\mathbf{S}_2^\infty e_\infty^{2\sigma_3}.
\]
Furthermore, one can verify using the definition of $\hat{\mathbf{\Psi}}(\lambda, x)$ that
\begin{equation}
\label{eq:psi-hat-asymptotics-infty}
\hat{\mathbf{\Psi}}(\lambda, x) \lambda_{\lw}^{(\Theta_\infty - 1)\sigma_3/2} \ee^{-\ii x \lambda \sigma_3/2}= \mathbb{I} + \hat{\mathbf{\Psi}}^\infty_1(x) \lambda^{-1} + \mathcal{O}\big(\lambda^{-2}\big), \qquad \lambda \to \infty,
\end{equation}
where
\[
\hat{\mathbf{\Psi}}_1^\infty(x) := \sigma_+ \mathbf{\Psi}^\infty_1(x) \sigma_+ + \sigma_+ \mathbf{\Psi}_2^\infty(x) \sigma_- + \hat{\mathbf{S}}(x) \sigma_+ + \hat{\mathbf{S}}(x) \mathbf{\Psi}_1^\infty (x) \sigma_-.
\]
Similarly, one can check that
\begin{equation}
\label{eq:psi-hat-asymptotics-zero}
\hat{\mathbf{\Psi}}(\lambda, x) \lambda_{\lw}^{-(\Theta_0 + 1)\sigma_3/2} \ee^{\ii x \lambda^{-1} \sigma_3/2}= \hat{\mathbf{\Psi}}^0_0(x) + \hat{\mathbf{\Psi}}^0_1(x) \lambda + \mathcal{O}\big(\lambda^{2}\big), \qquad \lambda \to 0,
\end{equation}
where
\[
\hat{\mathbf{\Psi}}^0_0(x) = \hat{\mathbf{S}}(x) \mathbf{\Psi}^0_0(x) \sigma_- + \hat{\mathbf{S}}(x) \mathbf{\Psi}^0_1(x) \sigma_+ + \sigma_+ \mathbf{\Psi}^0_0(x) \sigma_+.
\]
The transformation $\mathbf{\Psi} \mapsto \hat{\mathbf{\Psi}}$ is invertible so long as $\hat{\Psi}_{0, 22}^0(x)$ does not identically vanish, and its inverse is given by
\[
\mathbf{\Psi}(x, \lambda) \mapsto \check{\mathbf{\Psi}}(x, \lambda) := \big( \sigma_- \lambda^{1/2}_{\lw} + \check{\mathbf{S}}(x) \lambda^{-1/2}_{\lw} \big) \ \mathbf{\Psi}(x, \lambda),
\]
where
\[
\check{\mathbf{S}}(x) := \begin{bmatrix}
1 & - \Psi_{0, 12}^0(x) /\Psi_{0, 22}^0(x) \vspace{1mm}\\ - \Psi_{1, 21}^\infty(x) & \Psi_{0, 12}^0(x) \Psi_{1, 21}^\infty(x) /\Psi_{0, 22}^0(x)
\end{bmatrix}.
\]
It follows that $\check{\mathbf{\Psi}}$ satisfies conditions similar to \eqref{eq:psi-hat-asymptotics-infty} and \eqref{eq:psi-hat-asymptotics-zero} as $\lambda$ approaches $\infty, 0$, respectively. That these operations are inverses of one another is the content of \cite[Lemma~1]{BMS18}.

In this way, starting with $\mathbf{\Psi}$ and iterating the map $\mathbf{\Psi \mapsto \hat{\mathbf{\Psi}}}$ (assuming $\Psi^0_{0, 11}(x)$, $\Psi^0_{0, 22}(x)$ do not identically vanish after each step), we may define the $n$th iterate of this Schlesinger transformation, which we denote $\mathbf{\Psi}_n$. This matrix, if it exists, satisfies the following Riemann--Hilbert problem.
\begin{rhp} \label{rhp:initial-with-n}
Fix generic monodromy parameters $( e_1, e_2 )$, $n \in \Z$, and $x>0$. We seek a $2 \times 2$ matrix function $\lambda\mapsto\mathbf \Psi(\lambda, x)$ satisfying:
\begin{itemize}\itemsep=0pt
\item  Analyticity: $\mathbf \Psi_n (\lambda, x)$ is analytic in $\C \setminus L^{(6)}$, where $L^{(6)} = \{ | \lambda| = 1\} \cup \ii \R$ is the jump contour shown in Figure {\rm \ref{fig:1}}.
\item  Jump condition: $\mathbf \Psi_n (\lambda, x)$ has continuous boundary values on $L^{(6)} \setminus \{0\}$ from each component of $\mathbb{C}\setminus L^{(6)}$, which satisfy
\[\mathbf \Psi_{n, +} (\lambda, x) = \mathbf \Psi_{n, -} (\lambda, x) \mathbf J_{\mathbf \Psi_n}(\lambda),\] where $\mathbf J_{\mathbf \Psi_n}(\lambda)$ is as shown in Figure {\rm \ref{fig:1}} but with the modification
\[
\mathbf{S}_2^0 e_0^{-2\sigma_3} \mapsto (-1)^n\mathbf{S}_2^0 e_0^{-2\sigma_3} \qquad \text{and} \qquad \mathbf{S}_2^\infty e_\infty^{2\sigma_3} \mapsto (-1)^n \mathbf{S}_2^\infty e_\infty^{2\sigma_3}.
\]
\item  Normalization: $\mathbf \Psi_n(\lambda,x)$ satisfies the asymptotic conditions
\begin{equation}
\label{eq:Psi-n-asymptotic-infinity}
\boldsymbol \Psi_n(\lambda, x) = \big( \mathbb{I}+\mathbf{\Xi}^{(6)}_n(x) + \mathcal{O}\big(\lambda^{-2}\big) \big) \ee^{\ii x \lambda \sigma_3 /2} \lambda_{\lw}^{(n - \Theta_\infty) \sigma_3/2} \qquad \text{as} \quad \lambda \to \infty,
\end{equation}
and
\begin{equation}
\label{eq:Psi-n-asymptotic-zero}
\boldsymbol \Psi_n(\lambda, x) = \big( \mathbf{\Delta}^{(6)}_n(x) + \mathcal{O}(\lambda) \big) \ee^{-\ii x \lambda^{-1} \sigma_3 /2} \lambda_{\lw}^{(\Theta_0 + n) \sigma_3/2} \qquad \text{as} \quad \lambda \to 0,
\end{equation}
where $\mathbf{\Delta}^{(6)}_n(x)$ is a matrix determined from $\mathbf{\Psi}_n(\lambda,x)$ having unit determinant.
\end{itemize}
\end{rhp}
That $\mathbf{\Psi}_n$ solves the above Riemann--Hilbert problem implies the existence of the limit
\begin{equation}
\label{eq:T-n-definition}
 \mathbf{\Xi}^{(6)}_n(x):=\lim_{\lambda\to\infty}\lambda\big[\mathbf{\Psi}_n(\lambda,x)\ee^{-\ii x\lambda\sigma_3/2}\lambda_{\lw}^{\Theta_\infty\sigma_3/2}-\mathbb{I}\big].
\end{equation}
It follows that the function
\begin{equation}
\label{eq:u-n-recover}
u_n(x) = \dfrac{-\ii \Xi^{(6)}_{n, 12}(x)}{\Delta^{(6)}_{n, 11}(x) \Delta^{(6)}_{n, 12}(x)}
\end{equation}
satisfies PIII($D_6$) in the form
\begin{equation*}\label{eq:PIII-D6-n}
u_n''=\dfrac{(u_n')^2}{u_n}-\dfrac{u_n'}{x}+\dfrac{4(n+\Theta_0)u_n^2}{x}+\dfrac{4(1 + n-\Theta_\infty)}{x}+4u_n^3-\frac{4}{u_n}.
\end{equation*}
It was shown in \cite[Lemma 2]{BMS18} that if for some $n\in\mathbb{Z}$ the inverse monodromy problem is solvable for a given $x \in D$, where $D$ is a domain in $\C \setminus \{0\}$, then $\mathbf{\Psi}_n$ satisfies the Lax pair
\begin{align*}
\dpd{\mathbf{\Psi}_n}{\lambda}(\lambda, x) &= \Bigg( \frac{\ii x}{2}\sigma_3 + \frac{1}{2\lambda}\begin{bmatrix} n-\Theta_\infty & 2y \\ 2v & \Theta_\infty-n \end{bmatrix} + \frac{1}{2\lambda^2} \begin{bmatrix} \ii x - 2\ii st & 2\ii s \\ 2\ii t(x - st) & 2\ii st-\ii x\end{bmatrix} \Bigg) \mathbf{\Psi}_n(\lambda, x), \\
\dpd{\mathbf{\Psi}_n}{x}(\lambda, x) &= \Bigg(\dfrac{\ii \lambda}{2} \sigma_3 + \dfrac{1}{x} \begin{bmatrix} 0 & y \\ v & 0\end{bmatrix} - \dfrac{1}{2\lambda x} \begin{bmatrix} \ii x - 2\ii st & 2\ii s \\ 2\ii t(x -st ) & 2\ii st-\ii x\end{bmatrix}\Bigg)\mathbf{\Psi}_n(\lambda, x),
\end{align*}
where potentials $s$, $t$, $u$, $v$, $y$ all depend on $x$ and $n$. Furthermore, in this domain, the functions $\Psi_{0, 11}^{0}(x)$, $\Psi_{0, 22}^{0}(x)$ extracted from $\mathbf{\Psi}_n(\lambda,x)$ are not identically zero.\footnote{Lemma~2 in~\cite{BMS18} was stated for parameters corresponding to rational solutions of Painlev\'e-III, but the proof is almost exactly the same in this case.}

One can check that if a solution to Riemann--Hilbert Problem~\ref{rhp:initial-with-n} exists, it must be unique, and we attempt to identify this solution as a solution of Riemann--Hilbert Problem~\ref{rhp:initial} with possibly different monodromy data. The diagonal elements of $\mathbf{S}_2^0 e_0^{-2\sigma_3}$, $\mathbf{S}_2^\infty e_\infty^{2\sigma_3}$ alternate signs which implies the change
\[
e_0^2 \mapsto (-1)^n e_0^2, \qquad e_\infty^2 \mapsto (-1)^n e_\infty^2.
\]
Furthermore, in view of \eqref{eq:stokes-parameters}, we can write
\[
(-1)^n s_2^\infty e_\infty^2 = (-1)^n\big(1 - e_1^2 e_\infty^2\big) e_\infty^2 = \big(1 - (-1)^ne_1^2(-1)^ne_\infty^2\big)(-1)^ne_\infty^2,
\]
and
\[
(-1)^n s_2^0 e_0^{-2} = (-1)^n\big(e_1^2 e_0^2 - 1\big) e_0^{-2} = \big((-1)^ne_1^2(-1)^ne_0^2 -1\big)(-1)^ne_0^{-2}.
\]
Combining the above with the fact that $\mathbf{C}_{0 \infty}^\pm$ remain invariant under the iterated Schlesinger transformations implies the change in monodromy data
\begin{gather} \label{eq:monodromy-n-dependence}
e_1^2 \mapsto (-1)^ne_1^2 \qquad \text{and} \qquad e_2 \mapsto e_2.
\end{gather}
Since $e_1$, $e_2$ are assumed to be nonvanishing, we may write them in the form \eqref{eq:e1-e2-mu-eta}
for some ${\mu, \eta \in \C}$ with $-1 < \re(\mu), \re(\eta) \leq 1$. Moreover, since the transformations $e_1\mapsto-e_1$, ${e_2\mapsto -e_2}$ preserve the monodromy data, we can assume $-\frac{1}{2}<\re(\eta)\le\frac{1}{2}$ and $-\frac{1}{2}<\re(\mu) \leq \frac{1}{2}$. Equation~\eqref{eq:monodromy-n-dependence} implies in turn that $\eta$ does not depend on $n \in \Z$, while $\mu$ is replaced with
\begin{equation}
\label{eq:mu-n}
\mu \mapsto \mu_n := \begin{cases} \mu, & n \in 2\Z, \\ \mu-\frac{1}{2}, & n +1 \in 2\Z, \end{cases}
\end{equation}
This proves Proposition \ref{prop:schlesinger}. We end this section with two important remarks.

\begin{Remark}
\label{remark:e2-mu-minus-mu-change}
It was noted in the introduction that one could restrict $0 < \re(\mu_n) \leq 1/2$, in which case, the above iterations interchange the roles of $e_1^2$, $e_1^{-2}$ and we have to perform the transformation $\mu\to -\mu$, which corresponds to the replacements
\begin{gather*}
\mathbf{E}^\infty\to \left(\sqrt{\frac{e_\infty^2-e_1^2}{e_1^2\big(e_1^2e_\infty^2-1\big)}}\right)^{\sigma_3}\mathbf{E}^\infty\left(\sqrt{\frac{e_\infty^2-e_1^2}{e_1^2\big(e_1^2e_\infty^2-1\big)}}\right)^{-\sigma_3}\sigma_1\left(\frac{e_1^2 e_\infty^2 \big(1-e_1^2 e_\infty^2\big)}{ \big(e_1^4-1\big)}\right)^{\sigma_3},\\
\mathbf{E}^0\to \left(\sqrt{\frac{e_0^2-e_1^2}{e_1^2\big(e_1^2e_0^2-1\big)}}\right)^{\sigma_3}\mathbf{E}^0\left(\sqrt{\frac{e_0^2-e_1^2}{e_1^2\big(e_1^2e_0^2-1\big)}}\right)^{-\sigma_3}\sigma_1\left(\frac{e_1^2 \big(e_0^2 e_1^2-1\big)}{ e_0^2 \big(e_1^4-1\big)}\right)^{\sigma_3}.
\end{gather*}
This gauge transformation then allows us to identify the monodromy parameter pairs
\begin{gather}
(e_1,e_2)\sim \left(\frac{1}{e_1},\frac{1}{e_2e_0^2e_\infty^2}\sqrt{\frac{\big(e_1^2-e_0^2\big)\big(1-e_0^2e_1^2\big)}{\big(e_1^2-e_\infty^2\big)\big(1-e_1^2e_\infty^2\big)}}\right).
\label{eq:e1-to-1/e1}
\end{gather}
Therefore, we alternatively can write the monodromy data for the Schlesinger transformation as
\begin{gather*}
\mu_n = \begin{cases} \mu, & n \in 2\Z, \\ \frac{1}{2}-\mu, & n +1 \in 2\Z, \end{cases} \\
 e_{2,n}=\begin{cases} e_2, & n \in 2\Z, \\ \dfrac{1}{e_2e_0^2e_\infty^2}\sqrt{\dfrac{\big(e_1^2-e_0^2\big)\big(1-e_0^2e_1^2\big)}{\big(e_1^2-e_\infty^2\big)\big(1-e_1^2e_\infty^2\big)}}, & n +1 \in 2\Z. \end{cases}
\end{gather*}
Furthermore, one can check that $(x_1, x_2, x_3)$ in \eqref{eq:x1}--\eqref{eq:x3} remain invariant under the map described in \eqref{eq:e1-to-1/e1}, whereas $(y_1, y_2, y_3) \mapsto (\pm y_1, \pm y_2, y_3)$ where the sign depends on the choice of the square root in \eqref{eq:e1-to-1/e1} and \eqref{eq:V-root} below. In both cases, the corresponding solution of \eqref{eq:PIII-$D_8$} remains invariant, see Remark \ref{rem:minus-y1-y2}.
\end{Remark}

\begin{Remark}
 \label{rem:n-dependent-notation}
 Moving forward, we will slightly abuse notation by suppressing the $n$-de\-pen\-den\-ce in the parameters
\begin{equation}
\label{eq:e-definitions}
e_\infty = \ee^{\pi \ii (\Theta_\infty - n)/2}, \qquad e_0 = \ee^{\pi \ii (\Theta_0 + n)/2}, \qquad e_1 = \ee^{\pi \ii \mu_n}, \qquad e_2 = \ee^{\pi \ii \eta }.
\end{equation}
\end{Remark}

\section[Asymptotics for large n and small x and proof of Theorem 1.4]{Asymptotics for large $\boldsymbol{n}$ and small $\boldsymbol{x}$\\ and proof of Theorem \ref{thm:general}}
\label{sec:proof}
Let $(e_1, e_2)$ be generic monodromy parameters, see Definition \ref{def:generic}. At this point, we can see more clearly the meaning of the genericity conditions formulated there:
\begin{enumerate}\itemsep=0pt
\item[(i)] $e_1^4 \neq 1$; this is to guarantee diagonalizability in \eqref{eq:Stokes-products-eigenvectors},
\item[(ii)] $e_1 e_2 \neq 0$; this is to guarantee the unit-determinant condition in \eqref{eq:Stokes-products-eigenvectors} and \eqref{eq:C-zero-infty-diagonalization},
\item[(iii)] $e_1^2 \neq e_\infty^{\pm 2}$ and $e_1^2 \neq e_0^{\pm 2}$; this, in particular, implies that the Stokes multipliers \eqref{eq:stokes-parameters} are nonvanishing.
\end{enumerate}

\subsection{Opening the lenses} First, we define a new unknown matrix by $\mathbf \Phi_n(\lambda,x):=\boldsymbol \Psi_n (\lambda,x)\mathbf{L}$ where $\mathbf{L}$ is the piece-wise constant matrix shown in the left-hand panel of Figure~\ref{2}. It follows from \eqref{eq:C-zero-infty-diagonalization}, \eqref{eq:C-zero-infty-plus-diagonalization} that the resulting jump conditions satisfied by $\mathbf \Phi_n(\lambda,x)$ are as shown in the right-hand panel of Figure~\ref{2}.

 \begin{figure}[t]
 \centering
 \raisebox{7mm}{\includegraphics[width=0.43\linewidth]{Figures/2b2.pdf}}\quad
 \includegraphics[width=0.53\linewidth]{Figures/3b2.pdf}
\caption{Left panel: the definition of the matrix $\mathbf L$; the circles are centered at the origin and have radii $\frac{1}{2}$, $1$, and $2$. Right panel: the jump contour $\Gamma$ and jump conditions for $\mathbf{\Phi}_n(\lambda,x)$.}\label{2}
\end{figure}

\begin{Remark}\label{rmk:rotated-lenses}
In the general case $|{\arg}(x)|<\pi$, the lenses shown in Figure \ref{2} must be rotated in the manner shown in the left panel of Figure \ref{fig:2-rotated}. The resulting jumps follow from the identities~\eqref{eq:general-connection-factorization-1}--\eqref{eq:general-connection-factorization-2} and are shown in the right panel of Figure \ref{fig:2-rotated}.
\begin{figure}[t]
 \centering
\raisebox{-0.5\height}{\includegraphics[width=0.435\linewidth]{Figures/2b2-rotated.pdf}}\hfill%
\raisebox{-0.5\height}{\includegraphics[width=0.535\linewidth]{Figures/3b2-rotated.pdf}}
\caption{Analogue of Figure \ref{2} when $\arg(x) \neq 0$. The thick line represents the branch cut for the argument chosen as in Remark \ref{rmk:general-arg-x}.}
\label{fig:2-rotated}
\end{figure}
\end{Remark}

 \begin{figure}[h!]
 \centering
\includegraphics[width=0.5\columnwidth]{Figures/4b2.pdf}
\caption{The contour $\Gamma^{(\infty)}$ includes the circle $|\lambda|=2$.}\label{fig:6}
\end{figure}

\subsection[Parametrix for Phi\_n(lambda,x) near lambda=infty]{Parametrix for $\boldsymbol{{\Phi}_n(\lambda,x)}$ near $\boldsymbol{\lambda=\infty}$}
\label{sec:Parametrix-infinity}
By definition, the parametrix
%\para{\mathbf{\Phi}}^{(\infty)}(\lambda)=
$\para{\boldsymbol \Phi}_n^{(\infty)}(\lambda,x)$
satisfies the following Riemann--Hilbert problem.
\begin{rhp} \label{rhp:infinity_parametrix}
Fix generic monodromy parameters $( e_1, e_2 )$ determining the Stokes and connection matrices, $n\in\mathbb{Z}$, and $x>0$. We seek a $2 \times 2$ matrix function $\lambda\mapsto\para{\boldsymbol \Phi}_n^{(\infty)}(\lambda,x)$ satisfying:
\begin{itemize}\itemsep=0pt
\item Analyticity: $\para{\boldsymbol \Phi}_n^{(\infty)}(\lambda,x)$ is analytic in $\C \setminus \Gamma^{(\infty)}$, where
\[\Gamma^{(\infty)} = \{ | \lambda| = 2\} \cup (\ii \R\cap \{ |{\im}\,\lambda-1| > 1\})\]
 is the jump contour shown in Figure {\rm\ref{fig:6}}.
\item  Jump condition: $\para{\boldsymbol \Phi}_n^{(\infty)}(\lambda,x)$ has continuous boundary values on $\Gamma^{(\infty)} \setminus \{0\}$ from each component of $\mathbb{C}\setminus \Gamma^{(\infty)}$, which satisfy
\[
\para{\boldsymbol \Phi}_{n,+}^{(\infty)}(\lambda,x) = \para{\boldsymbol \Phi}_{n,-}^{(\infty)}(\lambda,x)\mathbf J_{\para{\boldsymbol \Phi}_n^{(\infty)}}(\lambda),
\]
 where \smash{$\mathbf J_{\para{\boldsymbol \Phi}_n^{(\infty)}}(\lambda)$} is as shown in Figure {\rm\ref{fig:6}} and where the $+$ $($resp., $-)$ subscript denotes a~boundary value taken from the left $($resp., right$)$ of an arc of $\Gamma^{(\infty)}$.
\item Normalization: \smash{$\para{\boldsymbol \Phi}_n^{(\infty)}(\lambda,x)$} satisfies the asymptotic conditions{\samepage
\begin{gather}
\label{eq:Phi-asymptotic-infinity}
\para{\boldsymbol \Phi}_n^{(\infty)}(\lambda,x) = \bigg(\mathbb I + \dfrac{\mathbf{A}_{n}(x)}{\lambda} + \mathcal{O}\left( \dfrac{1}{\lambda^2}\right) \bigg) \ee^{\ii x \lambda \sigma_3/2} \lambda_{\lw}^{(n - \Theta_\infty )\sigma_3/2} \qquad \text{as} \quad \lambda \to \infty,
\\
\label{eq:Phi-asymptotic-zero}
\para{\boldsymbol \Phi}_n^{(\infty)}(\lambda,x) = ( \mathbf{B}_{n}(x) + \mathcal{O}(\lambda) ) \lambda_{\lw}^{\mu_n \sigma_3} \qquad \text{as} \quad \lambda \to 0,
\end{gather}
where $\mathbf{A}_n(x)$ has zero trace and $ \mathbf{B}_{n}(x)$ has unit determinant.}
\end{itemize}
\end{rhp}



It is easy to see that \smash{$\para{\mathbf{\Phi}}_n^{(\infty)}(\lambda,x)$} necessarily has unit determinant. Furthermore, note that the jump matrix being \smash{$e_1^{-2\sigma_3}$} across the arc terminating at the origin implies
$e_1 = \ee^{\pi \ii \mu_n}$, which is consistent with \eqref{eq:e-definitions}.

\subsubsection[Dependence on lambda]{Dependence on $\boldsymbol{\lambda}$}
It follows from assuming differentiability of the asymptotics in \eqref{eq:Phi-asymptotic-infinity}--\eqref{eq:Phi-asymptotic-zero} that
\begin{gather}
 \frac{\partial\para{\boldsymbol \Phi}_{n}^{(\infty)}}{\partial\lambda}(\lambda,x) \para{\boldsymbol \Phi}^{(\infty)}_n(\lambda,x)^{-1}\nonumber\\
\qquad= \left(\mathbb I + \dfrac{\mathbf{A}_{n}(x)}{\lambda} + \mathcal{O}\big( \lambda^{-2}\big) \right)\left( \dfrac{\ii x}{2} + \dfrac{n - \Theta_\infty}{2 \lambda} \right) \sigma_3 \left(\mathbb I + \dfrac{\mathbf{A}_{n}(x)}{\lambda} + \mathcal{O}\big( \lambda^{-2}\big) \right)^{-1} + \mathcal{O}\big(\lambda^{-2}\big)\nonumber\\
 \qquad= \dfrac{\ii x}{2} \sigma_3 + \left( \dfrac{\ii x}{2} [\mathbf{A}_{n}(x), \sigma_3] + \dfrac{n - \Theta_\infty}{2} \sigma_3 \right) \dfrac{1}{\lambda} + \mathcal{O}\big( \lambda^{-2} \big) \qquad \text{as} \quad \lambda \to \infty,
\label{psi-infty-expansion}
\end{gather}
and
\begin{align}
 \frac{\partial\para{\boldsymbol \Phi}_{n}^{(\infty)}}{\partial\lambda}(\lambda,x)\para{\boldsymbol \Phi}^{(\infty)}_n(\lambda,x)^{-1} &= \left( \mathbf{B}_{n}(x) + \mathcal{O}(\lambda) \right) \left( \dfrac{\mu_n}{\lambda} \sigma_3 \right) \left( \mathbf{B}_{n}(x) + \mathcal{O}(\lambda) \right)^{-1} + \mathcal{O}(1)
 \nonumber\\
 &= \dfrac{\mu_n}{\lambda} \mathbf{B}_{n}(x) \sigma_3 \mathbf{B}_{n}(x)^{-1} + \mathcal{O}(1) \qquad \text{as} \quad \lambda \to 0.
\label{eq:Lambda-expansion-zero}
\end{align}
Since the quantity on the left-hand side of \eqref{psi-infty-expansion} and \eqref{eq:Lambda-expansion-zero} is otherwise an analytic function of~$\lambda$, it follows from Liouville's Theorem that
\begin{gather}
 \frac{\partial\para{\boldsymbol \Phi}_{n}^{(\infty)}}{\partial\lambda}(\lambda,x) \para{\boldsymbol \Phi}_n^{(\infty)}(\lambda,x)^{-1} = \dfrac{\ii x}{2} \sigma_3 + \dfrac{\mu_n}{\lambda} \mathbf{B}_{n}(x) \sigma_3 \mathbf{B}_{n}(x)^{-1} \implies\nonumber\\ \frac{\partial\para{\boldsymbol \Phi}_{n}^{(\infty)}}{\partial\lambda}(\lambda,x)= \left( \dfrac{\ii x}{2} \sigma_3 + \dfrac{\mu_n}{\lambda} \mathbf{B}_{n}(x) \sigma_3 \mathbf{B}_{n}(x)^{-1} \right) \para{\boldsymbol \Phi}_n^{(\infty)} (\lambda,x).
 \label{Lax-system-lambda}
\end{gather}
Noting that $\Tr\big(\mathbf{B}_{n}(x)\sigma_3\mathbf{B}_{n}(x)^{-1}\big)=0$ and $\det \big(\mathbf{B}_{n}(x) \sigma_3 \mathbf{B}_{n}(x)^{-1}\big) = -1$, we may write
\begin{equation}
\mathbf{B}_{n}(x) \sigma_3 \mathbf{B}_{n}(x)^{-1} = \begin{bmatrix} a_n(x) & b_n(x) \\ c_n(x) & -a_n(x) \end{bmatrix} \qquad \text{subject to}\quad a_n(x)^2 + b_n(x)c_n(x) = 1
\label{eq:zero-coefficient-matrix}
\end{equation}
and use this form in \eqref{Lax-system-lambda} to write a coupled scalar system of differential equations satisfied by the elements $\phi_1(\lambda,x)$ and $\phi_2(\lambda,x)$ of the first and second rows, respectively, of any column of~$\para{\boldsymbol \Phi}_n^{(\infty)}(\lambda,x)$:
\begin{align}
\label{psi1-psi2-sys1} \frac{\partial\phi_{1}}{\partial\lambda}(\lambda,x) &= \left(\dfrac{\ii x}{2} + \dfrac{\mu_n a_n(x)}{\lambda}\right) \phi_1(\lambda,x) + \dfrac{\mu_n b_n(x)}{\lambda} \phi_2(\lambda,x), \\
\label{psi1-psi2-sys2} \frac{\partial\phi_{2}}{\partial\lambda}(\lambda,x) &= \dfrac{\mu_n c_n(x)}{\lambda} \phi_1(\lambda,x)-\left(\dfrac{\ii x}{2}+ \dfrac{\mu_n a_n(x)}{\lambda}\right) \phi_2(\lambda,x).
\end{align}
Before beginning to solve this system, observe that equating the coefficients of $\lambda^{-1}$ in \eqref{psi-infty-expansion} and~\eqref{eq:Lambda-expansion-zero} yields the identity
\begin{equation}
\label{C-A relation}
\mu_n \mathbf{B}_{n}(x) \sigma_3 \mathbf{B}_{n}(x)^{-1} = \dfrac{\ii x}{2} [\mathbf{A}_{n}(x), \sigma_3] + \dfrac{n - \Theta_\infty}{2} \sigma_3.
\end{equation}
Since $[\mathbf{A}_{n}(x), \sigma_3]$ is off-diagonal, we arrive at
\begin{equation}
\mu_n a_n(x) = \frac{n - \Theta_\infty}{2}.
\label{eq:little-an}
\end{equation}
Since $\mu_n$ and $n$ are constants, this equation implies that $a_n(x)$ is independent of $x$, so we will simply write $a_n$ going forward.
Now, solving for $\phi_1(\lambda,x)$ in \eqref{psi1-psi2-sys2} and eliminating it from \eqref{psi1-psi2-sys1} yields (assuming $c_n(x) \neq 0$ and using $b_n(x)c_n(x) = 1 - a_n^2$)
\begin{align*}
 \lambda \frac{\partial^2\phi_{2}}{\partial\lambda^2}(\lambda,x) + \frac{\partial\phi_{2}}{\partial\lambda}(\lambda,x) + \left[ \dfrac{\ii x}{2} + \dfrac{x^2}{4} \lambda - \ii x\mu_n a_n - \dfrac{\mu_n^2}{\lambda} \right]\phi_2(\lambda,x) = 0.
\end{align*}
It is easy to see that the first-order derivative term is removed by the substitution $\phi_2(\lambda,x)=\lambda^{-1/2}w(\lambda,x)$.
Indeed, $w(\lambda,x)$ satisfies
\begin{align*}
 \frac{\partial^2w}{\partial\lambda^2}(\lambda,x) + \bigg[ \dfrac{x^2}{4} + \ii x\left( \dfrac{1}{2} - \mu_n a_n \right) \dfrac{1}{\lambda} + \left( \dfrac{1}{4} - \mu_n^2 \right) \dfrac{1}{\lambda^2} \bigg]w(\lambda,x) &= 0.
\end{align*}
Finally, the explicit $x$-dependence in the coefficients can be removed by setting $Z:= \ii x \lambda$ and writing $w(\lambda,x)=W(Z)$. Note that the notation $W(Z)$ here is not related to $W(z)$ appearing in Section~\ref{sec:monodromy-rep-$D_8$}. In this case, $W(Z)$ satisfies the ordinary differential equation
\begin{equation}
 W''(Z) + \bigg[ - \dfrac{1}{4} + \left( \dfrac{1}{2} - \mu_n a_n \right) \dfrac{1}{Z} + \left( \dfrac{1}{4} - \mu_n^2 \right) \dfrac{1}{Z^2} \bigg] W(Z) = 0,
 \label{Whittaker}
\end{equation}
which is Whittaker's equation (see \cite[Chapter~13]{DLMF}) with parameter
\begin{equation}\label{kappa-mu}
 \kappa = \kappa_n:=\frac{1}{2} - \mu_n a_n = \dfrac{1+\Theta_\infty - n }{2}.
\end{equation}
Given $\phi_2(\lambda,x)=\lambda^{-1/2}W(Z)$ for $Z=\ii x\lambda$, and a solution $W(Z)$ of \eqref{Whittaker}, it follows from \eqref{psi1-psi2-sys2} that the corresponding first-row entry is
\begin{equation}
 \phi_1(\lambda,x)=\frac{\ii x\lambda^{1/2}}{\mu_n c_n(x)}\bigg(W'(Z)+\left(\frac{1}{2}-\frac{\kappa_n}{Z}\right)W(Z)\bigg).
 \label{first-row}
\end{equation}
A fundamental pair of solutions of \eqref{Whittaker} is given by $W(Z)=W_{\pm \kappa_n, \mu_n} (\pm Z)$, $\mathrm{arg}(\pm Z) \in (-\pi, \pi) $.
If we take the particular solution $\phi_2(\lambda,x)=\lambda^{-1/2}W_{\kappa_n,\mu_n}(Z)$, then
using the identity
\begin{equation*}
 %\label{w-prime-backwards} \nonumber
 W'_{\kappa,\mu_n}(Z)=\left(\dfrac{\kappa}{Z}-\dfrac{1}{2}\right)W_{\kappa,\mu_n}(Z)+\bigg(\left(\dfrac{1}{2}-\kappa\right)^2-\mu_n^2\bigg)\frac{1}{Z}W_{\kappa-1,\mu_n}(Z)
\end{equation*}
(see \cite[equation~(13.15.23)]{DLMF}) in \eqref{first-row} gives
\begin{gather*}
 \phi_2(\lambda,x)=\lambda^{-1/2}W_{\kappa_n,\mu_n}(Z)\\
\implies \phi_1(\lambda,x)=\frac{\ii x\lambda^{1/2}}{\mu_n c_n(x)}\bigg(\left(\frac{1}{2}-\kappa_n\right)^2-\mu_n^2\bigg)Z^{-1}W_{\kappa_n-1,\mu_n}(Z).
\end{gather*}
Likewise, if we take the particular solution $\phi_2(\lambda,x)=\lambda^{-1/2}W_{-\kappa_n,\mu_n}(-Z)$, then using the identity
\begin{equation*}
%\label{w-prime}\nonumber
 W'_{\kappa, \mu_n}(Z) = \left( \dfrac{1}{2} - \dfrac{\kappa}{Z} \right)W_{\kappa, \mu_n}(Z) - \dfrac{1}{Z} W_{\kappa+1, \mu_n}(Z)
\end{equation*}
(see \cite[equation~(13.15.26)]{DLMF}) in \eqref{first-row} yields
\begin{equation*}
 \phi_2(\lambda,x)=\lambda^{-1/2}W_{-\kappa_n,\mu_n}(-Z)\implies \phi_1(\lambda,x)=-\frac{\ii x\lambda^{1/2}}{\mu_n c_n(x)}Z^{-1}W_{1-\kappa_n,\mu_n}(-Z).
\end{equation*}
Taking linear combinations with coefficients depending generally on the parameter $x$, the general solution matrix for the system \eqref{Lax-system-lambda} can be written in the form
\begin{equation}
 \para{\boldsymbol\Phi}_n^{(\infty)}(\lambda,x)=\widetilde{\boldsymbol\Phi}^{(\infty)}_n(\lambda,x)\mathbf{K}(x),\qquad\widetilde{\boldsymbol\Phi}_n^{(\infty)}(\lambda,x):=\mathbf{H}(\lambda,x)\mathbf{W}(\ii\lambda x;\kappa_n,\mu_n),
 \label{Psi0-define}
\end{equation}
where $\widetilde{\mathbf{\Phi}}_n^{(\infty)}(\lambda,x)$ is a specific fundamental solution matrix of \eqref{Lax-system-lambda} constructed from
\begin{equation}
 \mathbf{H}(\lambda,x):=\lambda^{\sigma_3/2}\begin{bmatrix}\displaystyle\frac{\ii x}{\mu_n c_n(x)} & 0\\0 & 1\end{bmatrix}
 \label{H-define}
\end{equation}
and
\begin{gather}
 \mathbf{W}(Z;\kappa,\mu_n):=
 \begin{bmatrix}
 \alpha_{\kappa,\mu_n}Z^{-1}W_{\kappa-1,\mu_n}(Z) & -Z^{-1}W_{1-\kappa,\mu_n}(-Z)\\
 W_{\kappa,\mu_n}(Z) & W_{-\kappa,\mu_n}(-Z)
 \end{bmatrix},\nonumber\\ \alpha_{\kappa,\mu_n}:=\left(\frac{1}{2}-\kappa\right)^2-\mu_n^2,
 \label{M-alpha-define}
\end{gather}
in which $\kappa=\kappa_n$ and $\mu_n$ are given by \eqref{kappa-mu} and $\mathbf{K}(x)$ is a matrix of free coefficients.

\subsubsection[Dependence on x]{Dependence on $\boldsymbol{x}$}
Going back to \eqref{eq:Phi-asymptotic-infinity}--\eqref{eq:Phi-asymptotic-zero} and now assuming that the asymptotics are differentiable with respect to $x$,
\begin{gather*}
 \frac{\partial\para{\boldsymbol \Phi}^{(\infty)}_n}{\partial x}(\lambda,x) \para{\boldsymbol \Phi}_n^{(\infty)} (\lambda,x)^{-1}\\
 \qquad= \left(\mathbb I + \dfrac{\mathbf{A}_{n}(x)}{\lambda} + \mathcal{O}\big( \lambda^{-2}\big) \right) \left( \dfrac{\ii \lambda}{2}\sigma_3 \right)\left(\mathbb I + \dfrac{\mathbf{A}_{n}(x)}{\lambda} + \mathcal{O}\big( \lambda^{-2}\big)\right)^{-1} + \mathcal{O}\big( \lambda^{-1} \big)\\
 \qquad = \dfrac{\ii \lambda}{2} \sigma_3 + \dfrac{\ii}{2} [\mathbf{A}_{n}(x), \sigma_3] + \mathcal{O}\big( \lambda^{-1} \big) \qquad \text{as} \quad \lambda \to \infty,
\end{gather*}
and
\begin{align*}
 \frac{\partial\para{\boldsymbol \Phi}^{(\infty)}_n}{\partial x}(\lambda,x) \para{\boldsymbol \Phi}_n^{(\infty)}(\lambda,x)^{-1} &= \mathbf{B}_n'(x) \mathbf{B}_n(x)^{-1} + \mathcal{O}(\lambda) \qquad \text{as} \quad \lambda \to 0.
\end{align*}
So, applying Liouville's theorem yields
\begin{equation}
 \frac{\partial\para{\boldsymbol \Phi}^{(\infty)}_{n}}{\partial x} (\lambda,x) = \left( \dfrac{\ii \lambda}{2} \sigma_3 + \dfrac{\ii}{2} [\mathbf{A}_n(x), \sigma_3] \right) \para{\boldsymbol \Phi}_n^{(\infty)}(\lambda,x),
 \label{Lax-system-x}
\end{equation}
and it follows from \eqref{C-A relation} that
\begin{align}
\dfrac{\ii}{2}[\mathbf{A}_n(x), \sigma_3] &= \dfrac{1}{x} \left( \mu_n\mathbf{B}_n(x) \sigma_3 \mathbf{B}_n(x)^{-1} +\left(\kappa_n-\frac{1}{2}\right) \sigma_3 \right)\nonumber\\
& = \frac{1}{x}\begin{bmatrix} 0 & \dfrac{\mu_n\big(1 - a_n^2\big)}{c_n(x)} \\ \mu_n c_n(x) & 0 \end{bmatrix},
\label{eq:Ainfty-identity}
\end{align}
where $\mu_n$, $a_n$ are independent of $\lambda$, $x$. To determine the $x$-dependence of $c_n(x)$, we use \eqref{eq:Ainfty-identity} to assemble \eqref{Lax-system-lambda} (using also \eqref{C-A relation} and \eqref{kappa-mu}) and \eqref{Lax-system-x} to give the Lax system
\begin{gather}
 \frac{\partial\para{\boldsymbol \Phi}_{n}^{(\infty)}}{\partial\lambda}(\lambda,x) =\para{\boldsymbol\Lambda}(\lambda,x) \para{\boldsymbol \Phi}_n^{(\infty)}(\lambda,x),\nonumber\\
 \para{\mathbf{\Lambda}}(\lambda,x):=\dfrac{\ii x}{2} \sigma_3 + \dfrac{1}{\lambda}\begin{bmatrix} -\frac{1}{2}(\Theta_\infty -n) & \displaystyle \frac{\mu_n\big(1-a_n^2\big)}{c_n(x)} \\ \mu_n c_n(x) & \frac{1}{2}(\Theta_\infty-n) \end{bmatrix}, \label{lax-lambda} \\
 \frac{\partial\para{\boldsymbol \Phi}_{n}^{(\infty)}}{\partial x}(\lambda,x) =\para{\mathbf X}(\lambda,x) \para{\boldsymbol \Phi}_n^{(\infty)}(\lambda,x),\nonumber\\
 \para{\mathbf{X}}(\lambda,x):= \dfrac{\ii\lambda}{2} \sigma_3 + \dfrac{1}{x}\begin{bmatrix} 0 & \displaystyle \frac{\mu_n\big(1-a_n^2\big)}{c_n(x)} \\ \mu_n c_n(x) & 0 \end{bmatrix}. \label{lax-x}
\end{gather}
Since $\para{\boldsymbol\Phi}_n^{(\infty)}(\lambda,x)$ is a simultaneous fundamental solution matrix for these equations, the Lax system is compatible. The compatibility condition reads
\begin{equation*}
 \para{\boldsymbol\Lambda}_{x}(\lambda,x) - \para{\mathbf X}_{\lambda}(\lambda,x) + \big[\para{\boldsymbol{\Lambda}}(\lambda,x), \para{\mathbf X}(\lambda,x)\big] = \mathbf{0},
\end{equation*}
which is equivalent to
\begin{equation}
 xc_n'(x) = (1-2\kappa_n) c_n(x) \implies c_n(x) = \gamma_n x^{1-2\kappa_n},
 \label{c-of-x}
\end{equation}
for some constant $\gamma_n\neq 0$. Thus, the coefficient $c_n(x)$ is determined up to the choice of the constant $\gamma_n$. Note also that the coefficient matrices $\para{\boldsymbol{\Lambda}}(\lambda,x)$ and $\para{\mathbf{X}}(\lambda,x)$ are obviously related by the simple identity
\begin{equation}
 \para{\mathbf{X}}(\lambda,x)-\frac{\lambda}{x}\para{\boldsymbol\Lambda}(\lambda,x)=\frac{1}{x}\left(\kappa_n-\frac{1}{2}\right)\sigma_3.
 \label{XLambda-relation}
\end{equation}
\sloppy Since the fundamental matrix $\widetilde{\boldsymbol \Phi}_n^{(\infty)}(\lambda,x)$ defined by \eqref{Psi0-define} satisfies \eqref{lax-lambda}, then so does \[\para{\boldsymbol \Phi}_n^{(\infty)}(\lambda,x)=\widetilde{\boldsymbol\Phi}_n^{(\infty)}(\lambda,x) \mathbf K(x),\] and $\mathbf K(x)$ must now be chosen so that \eqref{lax-x} is satisfied.
Substituting into \eqref{lax-x}, we obtain an ordinary differential equation on $\mathbf{K}(x)$:
\begin{equation}
 \mathbf{K}'(x)=\bigg(\widetilde{\boldsymbol\Phi}_n^{(\infty)}(\lambda,x)^{-1}\mathbf{X}(\lambda,x)\widetilde{\boldsymbol\Phi}_n^{(\infty)}(\lambda,x)-
 \widetilde{\boldsymbol\Phi}_n^{(\infty)}(\lambda,x)^{-1}\frac{\partial\widetilde{\boldsymbol\Phi}_{n}^{(\infty)}}{\partial x}(\lambda,x)\bigg)\mathbf{K}(x).
 \label{K-ODE}
\end{equation}
Now, from the form of $\widetilde{\boldsymbol\Phi}_n^{(\infty)}(\lambda,x)$ written in \eqref{Psi0-define}, we have both
\begin{gather*}
 \widetilde{\boldsymbol{\Phi}}_n^{(\infty)}(\lambda,x)^{-1}\frac{\partial\widetilde{\boldsymbol\Phi}_{n}^{(\infty)}}{\partial x}(\lambda,x)=\widetilde{\boldsymbol{\Phi}}_n^{(\infty)}(\lambda,x)^{-1}\frac{\partial\mathbf{H}}{\partial x}(\lambda,x)\mathbf{H}(\lambda,x)^{-1}\widetilde{\boldsymbol\Phi}_n^{(\infty)}(\lambda,x) \\
\phantom{ \widetilde{\boldsymbol{\Phi}}_n^{(\infty)}(\lambda,x)^{-1}\frac{\partial\widetilde{\boldsymbol\Phi}_{n}^{(\infty)}}{\partial x}(\lambda,x)=}{}+ \ii\lambda\widetilde{\boldsymbol\Phi}_n^{(\infty)}(\lambda,x)^{-1}\mathbf{H}(\lambda,x)\mathbf{W}'(Z;\kappa_n,\mu_n),\\
 \widetilde{\boldsymbol{\Phi}}_n^{(\infty)}(\lambda,x)^{-1}\frac{\partial\widetilde{\boldsymbol\Phi}_{n}^{(\infty)}}{\partial\lambda}(\lambda,x)=\widetilde{\boldsymbol{\Phi}}_n^{(\infty)}(\lambda,x)^{-1}\frac{\partial\mathbf{H}}{\partial\lambda}(\lambda,x)\mathbf{H}(\lambda,x)^{-1}\widetilde{\boldsymbol\Phi}_n^{(\infty)}(\lambda,x) \\
\phantom{ \widetilde{\boldsymbol{\Phi}}_n^{(\infty)}(\lambda,x)^{-1}\frac{\partial\widetilde{\boldsymbol\Phi}_{n}^{(\infty)}}{\partial\lambda}(\lambda,x)=}{}+ \ii x\widetilde{\boldsymbol\Phi}_n^{(\infty)}(\lambda,x)^{-1}\mathbf{H}(\lambda,x)\mathbf{W}'(Z;\kappa_n,\mu_n),
\end{gather*}
so it follows that
\begin{gather*}
 \widetilde{\boldsymbol{\Phi}}_n^{(\infty)}(\lambda,x)^{-1}\frac{\partial\widetilde{\boldsymbol\Phi}_{n}^{(\infty)}}{\partial x}(\lambda,x)\\
\quad=\frac{\lambda}{x}\widetilde{\boldsymbol{\Phi}}_n^{(\infty)}(\lambda,x)^{-1}\frac{\partial\widetilde{\boldsymbol{\Phi}}_{n}^{(\infty)}}{\partial\lambda}(\lambda,x)\\
 \quad\phantom{=}{}+ \widetilde{\boldsymbol{\Phi}}_n^{(\infty)}(\lambda,x)^{-1}\left[\frac{\partial\mathbf{H}}{\partial x}(\lambda,x)\mathbf{H}(\lambda,x)^{-1}-\frac{\lambda}{x}\frac{\partial\mathbf{H}}{\partial\lambda}(\lambda,x)\mathbf{H}(\lambda,x)^{-1}\right]\widetilde{\boldsymbol{\Phi}}_n^{(\infty)}(\lambda,x)\\
 \quad=\widetilde{\boldsymbol{\Phi}}_n^{(\infty)}(\lambda,x)^{-1}\left[\frac{\lambda}{x}\breve{\boldsymbol\Lambda}(\lambda,x) + \frac{\partial\mathbf{H}}{\partial x}(\lambda,x)\mathbf{H}(\lambda,x)^{-1}-\frac{\lambda}{x}\frac{\partial\mathbf{H}}{\partial\lambda}(\lambda,x)\mathbf{H}(\lambda,x)^{-1}\right]\widetilde{\boldsymbol{\Phi}}_n^{(\infty)}(\lambda,x),
\end{gather*}
where we also used \eqref{lax-lambda}.
Using this in \eqref{K-ODE} along with the explicit definition \eqref{H-define} of $\mathbf{H}(\lambda,x)$ and the identities \eqref{c-of-x} and \eqref{XLambda-relation} gives
\begin{gather*}
 \mathbf{K}'(x)\mathbf{K}(x)^{-1}\\
\quad =\widetilde{\boldsymbol{\Phi}}_n^{(\infty)}(\lambda,x)^{-1}\left[\breve{\mathbf{X}}(\lambda,x)-\frac{\lambda}{x}\breve{\boldsymbol{\Lambda}}(\lambda,x)
\right.\\
\left.
\hphantom{\quad =\widetilde{\boldsymbol{\Phi}}_n^{(\infty)}(\lambda,x)^{-1}}{}
+\frac{\lambda}{x}\frac{\partial\mathbf{H}}{\partial\lambda}(\lambda,x)\mathbf{H}(\lambda,x)^{-1}-\frac{\partial\mathbf{H}}{\partial x}(\lambda,x)\mathbf{H}(\lambda,x)^{-1}\right]\widetilde{\boldsymbol{\Phi}}_n^{(\infty)}(\lambda,x)\\
 \quad =\widetilde{\boldsymbol\Phi}_n^{(\infty)}(\lambda,x)^{-1}\Bigg[\frac{1}{x}\left(\kappa_n-\frac{1}{2}\right)\sigma_3 +\frac{1}{2x}\sigma_3+\frac{\dd}{\dd x}\log\left(\frac{c_n(x)}{x}\right)\begin{bmatrix}1&0\\0&0\end{bmatrix}\Bigg]\widetilde{\boldsymbol{\Phi}}_n^{(\infty)}(\lambda,x)\\
 \quad=\frac{\kappa_n}{x}\widetilde{\boldsymbol{\Phi}}_n^{(\infty)}(\lambda,x)^{-1}\Bigg[\sigma_3-2
 \begin{bmatrix}1&0\\0&0\end{bmatrix}\Bigg]\widetilde{\boldsymbol{\Phi}}_n^{(\infty)}(\lambda,x)
=-\frac{\kappa_n}{x}\mathbb{I}.
\end{gather*}
Therefore, the $x$-dependence of the matrix $\mathbf{K}(x)$ is explicitly given by
\begin{equation*}
 \mathbf{K}(x)=x^{-\kappa_n}\mathbf{K},
\end{equation*}
where $\mathbf{K}$ is now independent of both $\lambda$ and $x$. However, as the domain of analyticity of $\para{\mathbf{\Phi}}_n^{(\infty)}(\lambda,x)$ in the $\lambda$-plane consists of three disjoint regions, we expect to have to specify a different matrix~$\mathbf{K}$ for each. Note also that the constant $\gamma_n$ remains to be determined.

\subsubsection[The parametrix Phi\_n\^infty(lambda,x) on the two regions with |lambda|>2]{The parametrix $\boldsymbol{\para{\mathbf{\Phi}}_n^{(\infty)}(\lambda,x)}$ on the two regions with $\boldsymbol{|\lambda|>2}$}
To fully specify the parametrix $\para{\boldsymbol{\Phi}}_n^{(\infty)}(\lambda,x)$ for $|\lambda|>2$, we concretely take the jump contours for $|\lambda|>2$ to lie along the real axis in the $Z$-plane, corresponding to $\R_+$ and $\R_-$, respectively. Thus, the part of the domain of analyticity of $\para{\boldsymbol\Phi}_n^{(\infty)}(\lambda,x)$ with $|\lambda|>2$ has two components, corresponding to the upper and lower half $Z$-planes. To properly define $\para{\boldsymbol{\Phi}}_n^{(\infty)}(\lambda,x)$ in these two exterior domains, we firstly take the matrix factor $\mathbf{H}(\lambda,x)$ defined in \eqref{H-define} in the precise form
\begin{equation}
 \mathbf{H}(\lambda,x)=\lambda_{\lw}^{\sigma_3/2}\begin{bmatrix}\displaystyle \frac{\ii x^{2\kappa_n}}{\mu_n \gamma_n} & 0\\0 & 1
 \end{bmatrix} = x^{\kappa_n}\lambda_{\lw}^{\sigma_3/2}x^{\kappa_n\sigma_3}\mathbf{D}_n,\qquad\mathbf{D}_n:=\begin{bmatrix}\displaystyle\frac{\ii}{\mu_n\gamma_n} & 0\\0 & 1\end{bmatrix}.
 \label{H-define-2}
\end{equation}
Then, we assume different constant matrices $\mathbf{K}=\mathbf{K}_n^\pm$ in the two domains by writing the parametrix for $|\lambda|>2$ and $\pm\im(Z)> 0$ as
\begin{align}
 &\para{\boldsymbol{\Phi}}_n^{(\infty)}(\lambda,x)=\para{\boldsymbol{\Phi}}_n^{(\infty)\pm}(\lambda,x)=x^{-\kappa_n}\mathbf{H}(\lambda,x)\mathbf{W}(\ii x\lambda;\kappa_n,\mu_n)\mathbf{K}_n^\pm \nonumber\\&\phantom{\para{\boldsymbol{\Phi}}_n^{(\infty)}(\lambda,x)}{}= \lambda_{\lw}^{\sigma_3/2}x^{\kappa_n\sigma_3}\mathbf{D}_n\mathbf{W}(\ii x\lambda;\kappa_n,\mu_n)\mathbf{K}_n^\pm.
 \label{infinity-parametrix-form}
\end{align}
We now express the matrices $\mathbf{K}_n^\pm$ in terms of the remaining constants $\mu_n$ and $\gamma_n$ by enforcing the asymptotic condition \eqref{eq:Phi-asymptotic-infinity} in each of the two sectors with $|\lambda|>2$. According to \cite[equation~(13.19.3)]{DLMF},
\begin{equation*}
 W_{\kappa,\mu_n}(Z)=\ee^{-Z/2}Z^\kappa\big(1+\mathcal{O}\big(Z^{-1}\big)\big),\qquad Z\to\infty,\qquad |\mathrm{arg}(Z)|\le\frac{3\pi}{2}-\delta
\end{equation*}
holds for each $\delta>0$. Hence also
\begin{gather*} % \label{F-matrix-expand}
\begin{split}
& \mathbf{W}(Z;\kappa,\mu_n)=\begin{bmatrix}
 \alpha_{\kappa,\mu_n}Z^{\kappa-2}\big(1+\mathcal{O}\big(Z^{-1}\big)\big) & (-Z)^{-\kappa}\big(1+\mathcal{O}\big(Z^{-1}\big)\big)\vspace{1mm}\\
 Z^\kappa\big(1+\mathcal{O}\big(Z^{-1}\big)\big) & (-Z)^{-\kappa}\big(1+\mathcal{O}\big(Z^{-1}\big)\big)
 \end{bmatrix}\ee^{-Z\sigma_3/2},\\
 & Z\to\infty,\qquad |\mathrm{arg}(Z)|\le\frac{3\pi}{2}-\delta.
\end{split}
\end{gather*}
Under the condition given on $\mathrm{arg}(Z)$, we have
\begin{equation*}
 (-Z)^{-\kappa}=Z^{-\kappa}\begin{cases} \ee^{\ii\pi\kappa},& \text{$\im(Z)>0$ \ (i.e., $0<\arg(Z)<\pi$),}\\
 \ee^{-\ii\pi\kappa},& \text{$\im(Z)<0$ \ (i.e., $-\pi<\arg(Z)<0$).}
 \end{cases}
 %\label{z-power-upper-lower-half-planes}
\end{equation*}
To calculate $Z^{\pm \kappa}$, we recall $Z=\ii x\lambda$ and use \cite[equation~(49)]{BMS18}:
\[
-\dfrac{\pi}{2} - \arg(x) < \mathrm{arg}_{\lw} (\lambda) < \dfrac{3\pi}{2} - \arg(x), \qquad |\lambda| \to \infty.
\]
Next,
\begin{align*}
\im(Z)>0 &\implies 0<\arg(Z) <\pi\!\!\!\! &&\implies\! -\frac{\pi}{2} - \arg(x)< \mathrm{arg}_{\lw}(\lambda) < \frac{\pi}{2} - \arg(x),\!& \\
\im(Z)<0 &\implies -\pi<\arg(Z) < 0\!\!\!\!\! &&\implies\! -\frac{3\pi}{2} - \arg(x) < \mathrm{arg}_{\lw}(\lambda) - 2\pi <-\dfrac{\pi}{2} - \arg(x)\!&
\end{align*}
and hence, for any $x \in \C \setminus \{0\}$ such that $|{\arg}(x)| < \pi$,
\[
Z^{\pm \kappa} = x^{\pm \kappa} \lambda_{\lw}^{\pm \kappa}\begin{cases} \ee^{\pm \ii \pi \kappa/2}, &\im(Z) > 0, \\ \ee^{\mp 3\ii \pi \kappa/2}, & \im(Z) < 0.
\end{cases}
\]
Therefore, for $\lambda$ large such that $\im(Z)>0$,
\begin{gather*}
 \lambda_{\lw}^{\sigma_3/2}x^{\kappa_n\sigma_3}\mathbf{D}_n\mathbf{W}(\ii x\lambda;\kappa_n,\mu_n)\\
 \qquad=
 \mathbf{D}_n
 \begin{bmatrix}
 \mathcal{O}(\lambda^{-1}) & \ee^{\ii\pi\kappa_n/2}\big(1+\mathcal{O}\big(\lambda^{-1}\big)\big)\\
 \ee^{\ii\pi\kappa_n/2}\big(1+\mathcal{O}\big(\lambda^{-1}\big)\big) & \mathcal{O}\big(\lambda^{-1}\big)
 \end{bmatrix}\lambda_{\lw}^{(\kappa_n -\frac{1}{2})\sigma_3}\ee^{-\ii x\lambda\sigma_3/2},
\end{gather*}
so choosing $\mathbf{K}_n^+$ so that \eqref{infinity-parametrix-form} is consistent with \eqref{eq:Phi-asymptotic-infinity} in the sector $\im(Z)>0$ requires that $\mathbf{K}_n^+$ is an off-diagonal matrix, namely,
\begin{equation*}
 \mathbf{K}_n^+:=\begin{bmatrix}
 0 & \ee^{-\ii\pi\kappa_n/2}\\
 -\ii\ee^{-\ii\pi\kappa_n/2}\mu_n\gamma_n & 0
 \end{bmatrix}.
 %\label{K-infinity-plus}\nonumber
\end{equation*}
Similarly, for $\lambda$ large such that $\im(Z)<0$,
\begin{gather*}
 \lambda_{\lw}^{\sigma_3/2}x^{\kappa_n\sigma_3}\mathbf{D}_n\mathbf{W}(\ii x\lambda;\kappa_n,\mu_n) \\
 \qquad= \mathbf{D}_n\begin{bmatrix}
 \mathcal{O}(\lambda^{-1}) & \ee^{\ii\pi\kappa_n/2}\big(1+\mathcal{O}\big(\lambda^{-1}\big)\big)\\
 \ee^{-3\ii\pi\kappa_n/2}\big(1+\mathcal{O}\big(\lambda^{-1}\big)\big) & \mathcal{O}\big(\lambda^{-1}\big)
 \end{bmatrix}\lambda_{\lw}^{(\kappa_n-\frac{1}{2})\sigma_3}\ee^{-\ii x\lambda\sigma_3/2},
\end{gather*}
so consistency of \eqref{infinity-parametrix-form} with \eqref{eq:Phi-asymptotic-infinity} in the sector $\im(Z)<0$ requires
\begin{equation*}
 \mathbf{K}_n^-:=\begin{bmatrix}
 0 & \ee^{3\ii\pi\kappa_n/2}\\
 -\ii\ee^{-\ii\pi\kappa_n/2}\mu_n\gamma_n & 0
 \end{bmatrix}.
% \label{K-infinity-minus}
\end{equation*}

Some additional useful information can be gleaned by enforcing on
$\para{\boldsymbol{\Phi}}_n^{(\infty)}(\lambda,x)$ the jump conditions for $|\lambda|>2$. The jump rays are illustrated in the $Z$-plane with their orientations in Figure~\ref{fig1}.
\begin{figure}
 \centering
 \begin{tikzpicture}
 \draw (-8, 0) -- (-2, 0);
 \node at (-5,0){\textbullet};
 \node[below] at (-5,0){$\infty$};

 \node at (-5, 1){$\im(Z) < 0$};
 \node at (-5, -1){$\im(Z) > 0$};

 \node at (-4, 0){$<$};
 \node at (-6, 0){$>$};
 \node[right] at (-2, 0){$Z \in \R_+$};
 \node[left] at (-8, 0){$Z \in \R_-$};
 \end{tikzpicture}
 \caption{Jump contour for $\para{\boldsymbol{\Phi}}_n^{(\infty)}(\lambda,x)$ near $\lambda=\infty$.}
\label{fig1}
\end{figure}
The Whittaker function $W_{\kappa,\mu_n}(Z)$ can be viewed as an analytic function on the cut plane $|{\arg}(Z)|<\pi$, and it follows from the connection formula \cite[equation~(13.14.13)]{DLMF} that the boundary values on the negative real axis are related by
\begin{align}
&W_{\kappa,\mu_n}(-Z+\ii 0)= \ee^{2\pi \ii\kappa}W_{\kappa,\mu_n}(-Z-\ii 0)\nonumber \\
&\hphantom{W_{\kappa,\mu_n}(-Z+\ii 0)=}{} +
 \frac{2\pi \ii\ee^{\ii\pi\kappa}}{\Gamma\big(\frac{1}{2}+\mu_n-\kappa\big)\Gamma\big(\frac{1}{2}-\mu_n-\kappa\big)}W_{-\kappa,\mu_n}(Z),\qquad Z>0.
\label{Whittaker-jump}
\end{align}
Note that the denominators in the second term on the right-hand side of \eqref{Whittaker-jump} are finite due to condition (iii) in the definition of generic data; see the beginning of Section \ref{sec:proof}. Indeed, it follows from~\eqref{eq:e-definitions} and \eqref{kappa-mu} that
\[
e_1^{\pm 2} = e_\infty^2 \ \Leftrightarrow \ \frac{1}{2} - \kappa_n \pm \mu_n \in \Z.
\]

On $Z\in\R_-$ the left ($+$) and right ($-$) boundary values correspond to limits from ${\im(Z)<0}$ and $\im(Z)>0$, respectively. Therefore, the second column of $\mathbf{W}(Z;\kappa,\mu_n)$ is continuous across~$\mathbb{R}_-$, and from \eqref{Whittaker-jump} (replacing $Z$ with $-Z$),
\begin{gather*}%\label{F-jump-1}
 \mathbf{W}_{-}(Z;\kappa,\mu_n)
 =
 \begin{bmatrix}
 \alpha_{\kappa,\mu_n}Z^{-1}W_{\kappa-1,\mu_n}(Z+\ii 0) & -Z^{-1}W_{1-\kappa,\mu_n}(-Z)\vspace{1mm}\\
 W_{\kappa,\mu_n}(Z+\ii 0) & W_{-\kappa,\mu_n}(-Z)
 \end{bmatrix}\\
 =\begin{bmatrix}
 \displaystyle \ee^{2\pi \ii\kappa}\alpha_{\kappa,\mu_n}Z^{-1}W_{\kappa-1,\mu_n}(Z-\ii 0) - \frac{2\pi \ii\ee^{\ii\pi \kappa}Z^{-1}W_{1-\kappa,\mu_n}(-Z)}{\Gamma\big(\frac{1}{2}+\mu_n-\kappa\big)\Gamma\big(\frac{1}{2}-\mu_n-\kappa\big)}& -Z^{-1}W_{1-\kappa,\mu_n}(-Z)\vspace{1mm}\\
 \displaystyle \ee^{2\pi \ii\kappa}W_{\kappa,\mu_n}(Z-\ii 0)+\frac{2\pi \ii\ee^{\ii\pi\kappa}W_{-\kappa,\mu_n}(-Z)}{\Gamma\big(\frac{1}{2}+\mu_n-\kappa\big)\Gamma\big(\frac{1}{2}-\mu_n-\kappa\big)} & W_{-\kappa,\mu_n}(-Z)
 \end{bmatrix}\\
 =\mathbf{W}_{+}(Z;\kappa,\mu_n)\begin{bmatrix}
 \ee^{2\pi \ii\kappa} & 0\vspace{1mm}\\
 \displaystyle \frac{2\pi \ii\ee^{\ii\pi\kappa}}{\Gamma\big(\frac{1}{2}+\mu_n-\kappa\big)\Gamma\big(\frac{1}{2}-\mu_n-\kappa\big)} & 1
 \end{bmatrix},\qquad Z<0.
\end{gather*}
Here, on the {third} line we used the definition \eqref{M-alpha-define} of $\alpha_{\kappa,\mu_n}$ and the factorial identity $\Gamma(\diamond+1)=\diamond\Gamma(\diamond)$. Since $\mathbf{H}(\lambda,x)$ is analytic across $\ii\mathbb{R}_+$, it follows that \smash{$\para{\mathbf{\Phi}}^{(\infty)}(\lambda)=\para{\mathbf{\Phi}}_n^{(\infty)}(\lambda,x)$} satisfies the jump condition
\begin{align*}
 \para{\boldsymbol{\Phi}}_+^{(\infty)}(\lambda)&=\para{\boldsymbol{\Phi}}_-^{(\infty)}(\lambda)\big[\mathbf{W}_-(\ii x\lambda;\kappa_n,\mu_n)\mathbf{K}_n^{+}\big]^{-1}\mathbf{W}_+(\ii x\lambda;\kappa_n,\mu_n)\mathbf{K}_n^-\\
 &=\para{\boldsymbol{\Phi}}_-^{(\infty)}(\lambda)
 \big(\mathbf{K}^+_n\big)^{-1}
 \begin{bmatrix}
 \ee^{2\pi \ii\kappa_n} & 0\vspace{1mm}\\
 \displaystyle \frac{2\pi \ii\ee^{\ii\pi\kappa_n}}{\Gamma\big(\frac{1}{2}+\mu_n-\kappa_n\big)\Gamma\big(\frac{1}{2}-\mu_n-\kappa_n\big)} & 1
 \end{bmatrix}^{-1}
 \mathbf{K}^-_n\\
 &=\para{\boldsymbol{\Phi}}_-^{(\infty)}(\lambda)\begin{bmatrix}
 1 & \displaystyle\frac{2\pi \ee^{\ii\pi\kappa_n}}{\mu_n\gamma_n\Gamma\big(\frac{1}{2}+\mu_n-\kappa_n\big)\Gamma\big(\frac{1}{2}-\mu_n-\kappa_n\big)}\\
 0 & 1
 \end{bmatrix},\qquad \lambda\in \ii\mathbb{R}_+.
\end{align*}
Requiring that this matches with the corresponding jump condition in Figure \ref{fig:6} gives the condition
\begin{equation*}% \label{infinity-match-jump-up}
 \frac{2\pi \ee^{\ii\pi\kappa_n}}{\mu_n\gamma_n\Gamma\big(\frac{1}{2}+\mu_n-\kappa_n\big)\Gamma\big(\frac{1}{2}-\mu_n-\kappa_n\big)} = \frac{e_\infty^2-e_1^2}{e_1^2e_\infty^4}.
\end{equation*}
For $Z\in\R_+$ the $\pm$ boundary values correspond to the limit from $\im(Z)\gtrless 0$. Therefore, now the first column of $\mathbf{W}(Z;\kappa,\mu_n)$ is continuous across $\mathbb{R}_+$, and from \eqref{Whittaker-jump},
\begin{align*} % \label{F-jump-2}
 &\mathbf{W}_{-}(Z;\kappa,\mu_n) =
 \begin{bmatrix}
 \alpha_{\kappa,\mu_n}Z^{-1}W_{\kappa-1,\mu_n}(Z) & -Z^{-1}W_{1-\kappa,\mu_n}(-Z+\ii 0)\vspace{1mm}\\
 W_{\kappa,\mu_n}(Z) & W_{-\kappa,\mu_n}(-Z+\ii 0)
 \end{bmatrix}\\
 & =\begin{bmatrix}
 \! \alpha_{\kappa,\mu_n}Z^{-1}W_{\kappa-1,\mu_n}(Z) \!&\displaystyle\!\! -\ee^{-2\pi \ii\kappa}Z^{-1} W_{1-\kappa,\mu_n}(-Z\!-\ii 0)+\frac{2\pi \ii\ee^{-\ii\pi\kappa}\alpha_{\kappa,\mu_n}Z^{-1}W_{\kappa-1,\mu_n}(Z)}{\Gamma\big(\frac{1}{2} +\mu_n+\kappa\big)\Gamma\big(\frac{1}{2} -\mu_n+\kappa\big)} \! \vspace{1mm}\\
 W_{\kappa,\mu_n}(Z)\!& \!\displaystyle \ee^{-2\pi \ii\kappa}W_{-\kappa,\mu_n}(-Z\!-\ii 0) +
 \frac{2\pi \ii \ee^{-\ii\pi\kappa}W_{\kappa,\mu_n}(Z)}{\Gamma\big(\frac{1}{2} +\mu_n+\kappa\big)\Gamma\big(\frac{1}{2} -\mu_n+\kappa\big)}
 \end{bmatrix} \\
 & =\mathbf{W}_{+}(Z;\kappa,\mu_n)\begin{bmatrix}
 1 & \displaystyle \frac{2\pi \ii\ee^{-\ii\pi\kappa}}{\Gamma\big(\frac{1}{2}+\mu_n+\kappa\big)\Gamma\big(\frac{1}{2}-\mu_n+\kappa\big)} \vspace{1mm} \\
 0 & \ee^{-2\pi \ii\kappa}
 \end{bmatrix},\qquad Z>0.
\end{align*}
Again here, the finiteness of the denominators is guaranteed by condition (iii) at the beginning of Section \ref{sec:proof}. Since $\mathbf{H}(\lambda,x)$ changes sign across $\ii\mathbb{R}_-$, we get that \smash{$\para{\mathbf{\Phi}}^{(\infty)}(\lambda)=\para{\mathbf{\Phi}}^{(\infty)}_n(\lambda,x)$} satisfies the jump condition
\begin{align*}
 \para{\boldsymbol{\Phi}}_+^{(\infty)}(\lambda)&=-\para{\boldsymbol{\Phi}}_-^{(\infty)}(\lambda)[\mathbf{W}_{-}(\ii x\lambda;\kappa_n,\mu_n)\mathbf{K}_n^-]^{-1}\mathbf{W}_{+}(\ii x\lambda;\kappa_n,\mu_n)\mathbf{K}_n^+\\
 &=\para{\boldsymbol{\Phi}}_-^{(\infty)}(\lambda)(\mathbf{K}^-_n)^{-1}\begin{bmatrix}
 -1 & \displaystyle -\frac{2\pi \ii\ee^{-\ii\pi\kappa_n}}{\Gamma\big(\frac{1}{2}+\mu_n+\kappa_n\big)\Gamma\big(\frac{1}{2}-\mu_n+\kappa_n\big)}\vspace{1mm}\\
 0 & -\ee^{-2\pi \ii\kappa_n}
 \end{bmatrix}^{-1}\mathbf{K}_n^+\\
 &=\para{\boldsymbol{\Phi}}_-^{(\infty)}(\lambda)\begin{bmatrix}
 -\ee^{2\pi \ii\kappa_n} & 0\vspace{1mm}\\
 \displaystyle \frac{2\pi \ee^{-\ii\pi\kappa_n}\mu_n\gamma_n }{\Gamma\big(\frac{1}{2}+\mu_n+\kappa_n\big)\Gamma\big(\frac{1}{2}-\mu_n+\kappa_n\big)} &-\ee^{-2\pi \ii\kappa_n}
 \end{bmatrix} \\
 &= \para{\boldsymbol{\Phi}}_-^{(\infty)}(\lambda) \begin{bmatrix}e_{\infty}^{2}&0 \vspace{1mm}\\ e_\infty^2s_2^\infty&e_{\infty}^{-2}\end{bmatrix} = \para{\boldsymbol{\Phi}}_-^{(\infty)}(\lambda) \mathbf{S}_2^\infty e_\infty^{2\sigma_3}, \qquad \lambda\in\mathrm{i}\mathbb{R}_-.
\end{align*}
The last two equalities follow by a direct calculation using the definitions of $s_2^\infty$, $\kappa_n$, $e_1$, and $e_\infty$ in \eqref{eq:stokes-parameters}, \eqref{kappa-mu}, and \eqref{eq:e-definitions}, respectively, along with the expression
\begin{equation} \label{Qinfty-gammainfty}
 \mu_n \gamma_n = \dfrac{2\pi \ee^{\pi \ii \kappa_n}}{\Gamma\big( \frac{1}{2} + \mu_n - \kappa_n\big)\Gamma\big( \frac{1}{2} - \mu_n - \kappa_n\big)} \cdot \frac{e_1^2 e_\infty^4}{e_\infty^2 - e_1^2},
\end{equation}
and the classical identity
\begin{equation}
\label{eq:Gamma-reflection}
\Gamma \left(\frac{1}{2} - z \right) \Gamma \left(\frac{1}{2} + z \right) = \dfrac{\pi }{\cos(\pi z)}.
\end{equation}

\subsubsection[The parametrix Phi\^(infty)\_n(lambda,x) in the region lambda less than 2 in absolute value]{The parametrix $\boldsymbol{\para{\boldsymbol{\Phi}}^{(\infty)}_n(\lambda,x)}$ in the region $\boldsymbol{|\lambda|<2}$}

We use the identity \cite[equation~(13.14.33)]{DLMF} to express the elements of $\mathbf{W}(Z;\kappa,\mu)$ in terms of the alternative basis of solutions $M_{-\kappa,\pm\mu}(-Z)$ of Whittaker's equation with parameters $(\kappa,\mu)$ that form a numerically satisfactory pair in a neighborhood of the origin and that are analytic for $\arg(-Z)\in (-\pi,\pi)$. Moreover, these functions are the Maclaurin series associated with the regular singular point at $Z=0$, so they have the property that
\begin{equation}
M_{\kappa,\mu}(-Z)(-Z)^{-\frac{1}{2}-\mu} = 1+\mathcal{O}(Z) \qquad \text{as} \quad Z\to 0,
\label{eq:WhittakerM-near-origin}
\end{equation}
where the power function denotes the principal branch and where the error term represents an analytic function of $Z$ vanishing at the origin. To deal with the first column of $\mathbf{W}(Z;\kappa,\mu)$ we also use the corresponding identity $M_{\kappa,\mu}(Z)=\ee^{\pm \ii\pi(\frac{1}{2}+\mu)}M_{-\kappa,\mu}(-Z)$ which holds for $\pm\im(Z)>0$ (see also \cite[equation~(13.14.10)]{DLMF}).
Using the above identities, and under the condition $2\mu \not \in \Z$ (which follows from the condition~(i) at the beginning of Section~\ref{sec:proof} in our case), we can write the elements of $\mathbf{W}(Z;\kappa,\mu)$ in the form
\begin{gather*}\nonumber
 W_{11}(Z;\kappa,\mu)= -Z^{-1}\frac{\Gamma(-2\mu)\Gamma\big(\frac{1}{2}-\mu+\kappa\big)}{\Gamma\big(\frac{1}{2}-\mu-\kappa\big)\Gamma\bigl(-\frac{1}{2}-\mu+\kappa\bigr)}\ee^{\pm \ii\pi(\frac{1}{2}+\mu)}M_{1-\kappa,\mu}(-Z)\\
\hphantom{W_{11}(Z;\kappa,\mu)=}{}
-Z^{-1}\frac{\Gamma(2\mu)\Gamma\big(\frac{1}{2}+\mu+\kappa\big)}{\Gamma\big(\frac{1}{2}+\mu-\kappa\big)\Gamma\bigl(-\frac{1}{2}+\mu+\kappa\bigr)}
\ee^{\pm \ii\pi(\frac{1}{2}-\mu)}M_{1-\kappa,-\mu}(-Z),\quad\! \pm\im(Z)>0,\\
 W_{12}(Z;\kappa,\mu)= -Z^{-1}\frac{\Gamma(-2\mu)}{\Gamma\bigl(-\frac{1}{2}-\mu+\kappa\bigr)}M_{1-\kappa,\mu}(-Z)-Z^{-1}\frac{\Gamma(2\mu)}{\Gamma\bigl({-}\frac{1}{2}+\mu+\kappa\bigr)}M_{1-\kappa,-\mu}(-Z),\\
 W_{21}(Z;\kappa,\mu)= \frac{\Gamma(-2\mu)}{\Gamma\big(\frac{1}{2}-\mu-\kappa\big)}\ee^{\pm \ii\pi (\frac{1}{2}+\mu)}M_{-\kappa,\mu}(-Z)\\
 \hphantom{W_{21}(Z;\kappa,\mu)= }{}
 +
 \frac{\Gamma(2\mu)}{\Gamma\big(\frac{1}{2}+\mu-\kappa\big)}\ee^{\pm \ii\pi(\frac{1}{2}-\mu)}M_{-\kappa,-\mu}(-Z),\qquad\pm\im(Z)>0,
\end{gather*}%
and
\begin{equation*}
 W_{22}(Z;\kappa,\mu)=\frac{\Gamma(-2\mu)}{\Gamma\big(\frac{1}{2}-\mu+\kappa\big)}M_{-\kappa,\mu}(-Z) +\frac{\Gamma(2\mu)}{\Gamma\big(\frac{1}{2}+\mu+\kappa\big)}M_{-\kappa,-\mu}(-Z).
\end{equation*}
These expressions can be usefully combined into a matrix identity:
\begin{equation}
 \mathbf{W}(Z;\kappa,\mu)=\mathbf{M}(Z;\kappa,\mu)\mathbf{G}^\pm_{\kappa,\mu},\qquad \pm\im(Z)>0,
 \label{F-to-N}
\end{equation}
where
\begin{gather}
 \mathbf{M}(Z;\kappa,\mu):=
 \begin{bmatrix}
 \big(\frac{1}{2}-\kappa+\mu\big)Z^{-1}M_{1-\kappa,\mu}(-Z) & \big(\frac{1}{2}-\kappa-\mu\big)Z^{-1}M_{1-\kappa,-\mu}(-Z)\vspace{1mm}\\
 M_{-\kappa,\mu}(-Z) & M_{-\kappa,-\mu}(-Z)
 \end{bmatrix},\nonumber\\
 \arg(-Z)\in (-\pi,\pi),
 \label{N-define}
\end{gather}
and
\begin{equation*} % \label{coefficients-C}
 \mathbf{G}^\pm_{\kappa,\mu}:=\begin{bmatrix}
 \displaystyle \frac{\Gamma(-2\mu)\ee^{\pm \ii\pi(\frac{1}{2}+\mu)}}{\Gamma\big(\frac{1}{2}-\mu-\kappa\big)} &
 \displaystyle \frac{\Gamma(-2\mu)}{\Gamma\big(\frac{1}{2}-\mu+\kappa\big)}\vspace{1mm}\\
 \displaystyle\frac{\Gamma(2\mu)\ee^{\pm \ii\pi(\frac{1}{2}-\mu)}}{\Gamma\big(\frac{1}{2}+\mu-\kappa\big)} &
 \displaystyle \frac{\Gamma(2\mu)}{\Gamma\big(\frac{1}{2}+\mu+\kappa\big)}
 \end{bmatrix}.
\end{equation*}

To define the parametrix $\para{\boldsymbol{\Phi}}_n^{(\infty)}(\lambda,x)$ for $|\lambda|<2$, we first introduce a constant matrix by
\begin{equation}
 \mathbf{J}_n:=\mathbf{G}_{\kappa_n,\mu_n}^+ \mathbf{K}_n^+ \mathbf{S}_1^{\infty} \mathbf{E}^{\infty}=\mathbf{G}_{\kappa_n,\mu_n}^- \mathbf{K}_n^- \mathbf{E}^{\infty}.
\label{eq:M-infty}
\end{equation}
The equality of these two expressions can be seen as follows. First, combining \eqref{infinity-parametrix-form} and \eqref{F-to-N}, and using the fact that the matrix $\mathbf{M}(\ii x\lambda;\kappa,\mu_n)$ is analytic in a neighborhood of $\lambda=2\ii$, the jump condition for \smash{$\para{\mathbf{\Phi}}_n^{(\infty)}(\lambda,x)$} across the positive imaginary axis for $|\lambda|>2$ shown in Figure~\ref{fig:6} implies the identity $\mathbf{G}_{\kappa_n,\mu_n}^+\mathbf{K}_n^+ \mathbf{S}_1^{\infty}=\mathbf{G}_{\kappa_n,\mu_n}^-\mathbf{K}_n^-$, which yields the desired equality.

Then, we set
\begin{equation}
 \para{\boldsymbol{\Phi}}_n^{(\infty)}(\lambda,x):= \lambda_{\lw}^{\sigma_3/2}x^{\kappa_n\sigma_3}\mathbf{D}_n\mathbf{M}(\ii x\lambda;\kappa_n,\mu_n)\mathbf{J}_n,\qquad\text{for} \quad |\lambda|<2.
 \label{infinity-parametrix-form-origin}
\end{equation}
It is straightforward to then check that, regardless of the choice of $\mu_n$,
the matrix $\mathbf{J}_n$ defined by~\eqref{eq:M-infty} is diagonal. Comparing \eqref{infinity-parametrix-form} and \eqref{infinity-parametrix-form-origin} shows that the jump conditions for~\smash{$\para{\mathbf{\Phi}}_n^{(\infty)}(\lambda,x)$} across the arcs of the circle $|\lambda|=2$ shown in Figure~\ref{fig:6} are satisfied. Using~\eqref{eq:WhittakerM-near-origin} then proves that~\smash{$\para{\mathbf{\Phi}}_n^{(\infty)}(\lambda,x)$} satisfies the simple jump condition across the negative imaginary axis with~${|\lambda|<2}$ shown in Figure~\ref{fig:6} and that an expansion of the form shown in \eqref{eq:Phi-asymptotic-zero} holds.
To check that the matrix $\mathbf{J}_n$ is diagonal and arrive at its final form below, we use the identity~\eqref{eq:Gamma-reflection} to get
\begin{gather}
\Gamma\left( \dfrac{1}{2} + \kappa_n - \mu_n \right) \Gamma\left( \dfrac{1}{2} - \kappa_n + \mu_n \right) = \dfrac{\pi}{\cos \left( \pi (\kappa_n - \mu_n)\right)} = \dfrac{2 \pi \ii e_1e_\infty}{e_1^2 - e_\infty^2}, \nonumber\\
\Gamma\left( \dfrac{1}{2} + \kappa_n + \mu_n \right) \Gamma\left( \dfrac{1}{2} - \kappa_n - \mu_n \right) = \dfrac{\pi}{\cos \left( \pi (\kappa_n + \mu_n)\right)} = \dfrac{2\pi \ii e_1e_\infty}{1 - e_1^2 e_\infty^2}.\label{eq:cosine-identity-1}
\end{gather}
The result is that the diagonal matrix $\mathbf{J}_n$ from \eqref{eq:M-infty} is given by
\begin{equation}
\mathbf{J}_n =
\begin{bmatrix}\dfrac{\ee^{{\ii\pi/4}} e_1 e_\infty^{7/2} \Gamma(-2 \mu_n )}{\Gamma \big(\frac{1}{2}-\kappa_n -\mu_n \big)}&0 \\
0&\dfrac{\ee^{{\ii\pi/4}} \big(e_1^4-1\big) e_\infty^{3/2} \Gamma(2 \mu_n )}{e_1 \big(e_1^2-e_\infty^2\big) \Gamma \big(\frac{1}{2}-\kappa_n +\mu_n \big)}
\label{eq:J-n}
\end{bmatrix}.
\end{equation}

\subsection[Parametrix for Phi\_n(lambda,x) near lambda=0]{Parametrix for $\boldsymbol{{\Phi}_n(\lambda,x)}$ near $\boldsymbol{\lambda=0}$}
By definition, the parametrix $\para{\boldsymbol \Phi}_n^{(0)}(\lambda,x)$ satisfies the following Riemann--Hilbert problem.
\begin{rhp} \label{rhp:zero_parametrix}
Fix generic monodromy parameters $( e_1, e_2 )$ determining the Stokes and connection matrices, $n\in\mathbb{Z}$, and $x>0$. We seek a $2 \times 2$ matrix function $\lambda\mapsto\para{\boldsymbol \Phi}_n^{(0)}(\lambda,x)$ satisfying:
\begin{itemize}\itemsep=0pt
\item Analyticity: $\para{\boldsymbol \Phi}_n^{(0)}(\lambda,x)$ is analytic in $\C \setminus \Gamma^{(0)}$, where $\Gamma^{(0)} = \big\{ | \lambda| = \frac{1}{2}\big\} \cup \big(\ii \R\cap \big\{ \im \lambda < \frac{1}{2}\big\}\big)$ is the jump contour shown in Figure {\rm\ref{fig:8}}.
\item  Jump condition: $\para{\boldsymbol \Phi}_n^{(0)}(\lambda,x)$ has continuous boundary values on $\Gamma^{(0)} \setminus \{0\}$ from each component of $\mathbb{C}\setminus \Gamma^{(0)}$, which satisfy
\[
\para{\boldsymbol \Phi}_{n,+}^{(0)}(\lambda,x) = \para{\boldsymbol \Phi}_{n,-}^{(0)}(\lambda,x)\mathbf J_{\para{\boldsymbol \Phi}_n^{(0)}}(\lambda),
\]
 where \smash{$\mathbf J_{\para{\boldsymbol \Phi}_n^{(0)}}(\lambda)$} is as shown in Figure {\rm\ref{fig:8}} and where the $+$ $($resp., $-)$ subscript denotes a boundary value taken from the left $($resp., right$)$ of an arc of $\Gamma^{(0)}$.
\item  Normalization: $\para{\boldsymbol \Phi}_n^{(0)}(\lambda,x)$ satisfies the asymptotic conditions
\begin{gather}
\label{eq:Phi_z-asymptotic-infinity}
\para{\boldsymbol \Phi}_n^{(0)}(\lambda,x) = \mathcal{O}(1) \lambda_{\lw}^{ {\mu_n \sigma_3}{}} \qquad \text{as} \quad \lambda \to \infty,
\end{gather}
where $\mathcal{O}(1)$ refers to a function analytic and bounded in a neighborhood of $\lambda=\infty$ and
\begin{equation}
\label{eq:Phi_z-asymptotic-zero}
\para{\boldsymbol \Phi}_n^{(0)}(\lambda,x) = ( \mathbb I + \mathcal{O}(\lambda) )\ee^{-\ii x\lambda^{-1}\sigma_3/2} \lambda_{\lw}^{ (n+ \Theta_0 )\sigma_3/2 } \qquad \text{as} \quad \lambda \to 0.
\end{equation}
\end{itemize}
\end{rhp}

 \begin{figure}
 \centering
\includegraphics[width=0.4\columnwidth]{Figures/5b2.pdf}
\caption{The contour $\Gamma^{(0)}$, which includes the circle $|\lambda|=\frac{1}{2}$.}\label{fig:8}
\end{figure}
We can write down the unique solution $\para{\mathbf{\Phi}}_n^{(0)}(\lambda,x)$ explicitly in terms of the parametrix \smash{$\para{\mathbf{\Phi}}_n^{(\infty)}(\lambda,x)$} obtained in Section~\ref{sec:Parametrix-infinity}, but taken with the index $1-n$ instead of $n$ and $\Theta_\infty$ replaced by $\Theta_0$. If we indicate the dependence of \smash{$\para{\mathbf{\Phi}}_n^{(\infty)}(\lambda,x)$} and \smash{$\para{\mathbf{\Phi}}_n^{(0)}(\lambda,x)$} on $\Theta_\infty$ and $\Theta_0$ respectively with the notation \smash{$\para{\mathbf{\Phi}}_n^{(\infty)}(\lambda,x)=\para{\mathbf{\Phi}}_n^{(\infty)}(\lambda,x,\Theta_\infty)$} and \smash{$\para{\mathbf{\Phi}}_n^{(0)}(\lambda,x)=\para{\mathbf{\Phi}}_n^{(0)}(\lambda,x,\Theta_0)$},
then we have the following.
\begin{prop}\label{prop:symmetry}
Fix $\Theta_\infty, \Theta_0 \in \C$ and generic monodromy parameters $(e_1, e_2)$. Then
\begin{equation}
 \para{\mathbf{\Phi}}^{(0)}_n(\lambda,x, \Theta_0)=\begin{cases}
 e_0^{3\sigma_3}\para{\mathbf{\Phi}}^{(\infty)}_{-n}\bigl(-\lambda^{-1},x, \Theta_0\bigr) e_0^{-2\sigma_3}\quad ,& |\lambda|<\frac{1}{2}, \vspace{1mm}\\
 e_0^{3\sigma_3}\para{\mathbf{\Phi}}^{(\infty)}_{-n}\bigl(-\lambda^{-1},x, \Theta_0\bigr)\begin{bmatrix}0 & \beta_n\\-\beta_n^{-1} & 0\end{bmatrix},& |\lambda|>\frac{1}{2},
 \end{cases}
 \label{eq:Phi0-Phiinfty}
\end{equation}
where
\begin{equation}
 \beta_n:= \dfrac{1-e_1^4 }{e_1^2\big(1 - e_0^2 e_1^2\big)}.
 \label{eq:a-def}
\end{equation}
\end{prop}
\begin{proof}
The mapping $\lambda\mapsto -\lambda^{-1}$ takes the contour $\Gamma^{(0)}$ onto the contour $\Gamma^{(\infty)}$ up to the reversal of orientation of certain arcs, and swaps the circles centered at the origin of radii $\frac{1}{2}$ and $2$. Therefore, the domain of analyticity of \smash{$\para{\mathbf{\Phi}}^{(0)}_n(\lambda,x, \Theta_0)$} is as desired.

Under the map $n \mapsto -n$, $\Theta_\infty \mapsto \Theta_0$, the exponentials defined in \eqref{eq:e-definitions} satisfy $e_{0}^2 \mapsto e_{\infty}^2$, whereas $\mu_n = \mu_{-n}$ since this quantity depends only on the parity of $n$. This implies that the Stokes matrices defined in \eqref{eq:Stokes-infty}--\eqref{eq:Stokes-zero} satisfy the corresponding identities
\begin{equation*}
 \mathbf{S}^\infty_1 \mapsto e_0^{-2\sigma_3}\big(\mathbf{S}^0_{1}\big)^{-1}e_0^{2\sigma_3} \qquad\text{and}\qquad
 \mathbf{S}^\infty_{{2}}e_\infty^{2\sigma_3} \mapsto e_0^{-2\sigma_3}\big(\mathbf{S}^0_{2} e_0^{-2\sigma_3}\big)^{-1} e_0^{2\sigma_3}.
\end{equation*}
Comparing Figures~\ref{fig:6} and \ref{fig:8} then shows that the function defined by \eqref{eq:Phi0-Phiinfty} satisfied the required jump conditions across the imaginary axis for $|\lambda|<\frac{1}{2}$. Likewise, the jump condition on the negative imaginary axis for $|\lambda|>\frac{1}{2}$ is easily verified due to the identity valid for any $a\neq 0$:
\begin{equation*}
 e_1^{-2 \sigma_3}\begin{bmatrix}0&a\\-a^{-1} & 0\end{bmatrix} +\begin{bmatrix}0 & a\\-a^{-1} & 0\end{bmatrix}e_1^{2 \sigma_3}=\mathbf{0}.
\end{equation*}
Finally, the fact that $\para{\mathbf{\Phi}}^{(0)}_n(\lambda,x, \Theta_0)$ defined by \eqref{eq:Phi0-Phiinfty} satisfies the required jump conditions across the circle $|\lambda|=\frac{1}{2}$ follows from the corresponding jump conditions for \smash{$\para{\mathbf{\Phi}}^{(\infty)}_{1-n}(\lambda,x, \Theta_0)$} for $|\lambda|=2$ and the identities
\begin{equation*}
 (\mathbf{E}^{\infty})^{-1} \mapsto \begin{bmatrix}0 & \beta_n\\-\beta_n^{-1} & 0\end{bmatrix} \big(\mathbf{S}_1^0 \mathbf{E}^0\big)^{-1} e_0^{2\sigma_3}, \qquad(\mathbf{S}_1^\infty \mathbf{E}^{\infty})^{-1} \mapsto \begin{bmatrix}0 & \beta_n\\-\beta_n^{-1} & 0\end{bmatrix} \big( \mathbf{E}^0\big)^{-1} e_0^{2\sigma_3}
\end{equation*}
among the matrices $\mathbf{E}^\infty$, $\mathbf{E}^0$ defined in \eqref{eq:E-infinity-formula}--\eqref{eq:E-zero-formula}, which hold for the value of $\beta_n$ indicated in~\eqref{eq:a-def}.

It only remains to verify the asymptotics in \eqref{eq:Phi_z-asymptotic-infinity}--\eqref{eq:Phi_z-asymptotic-zero}. However, these follow from the corresponding formul\ae\ in \eqref{eq:Phi-asymptotic-infinity}--\eqref{eq:Phi-asymptotic-zero} with the help of the identity
\begin{equation*}
 \lambda^p_{\lw}=\ee^{\ii \pi p}\zeta^{-p}_{\lw},\qquad\zeta:=-\lambda^{-1},
\end{equation*}
which holds for all $\lambda$ not on the negative imaginary axis.

Since the matrix function $\para{\mathbf{\Phi}}_n^{(0)}(\lambda,x, \Theta_0)$ defined by \eqref{eq:Phi0-Phiinfty}--\eqref{eq:a-def} satisfies all the required Riemann--Hilbert conditions, and there is at most one solution of those conditions, as is easily confirmed by a Liouville argument, the proof is finished.
\end{proof}

\subsection{An equivalent Riemann--Hilbert problem on the unit circle}
The parametrix for $\mathbf{\Phi}_n(\lambda,x)$ is by definition the following matrix function:
\begin{equation*} % \label{eq:Parametrix}
 \para{\mathbf{\Phi}}_n(\lambda,x):=\begin{cases}
 \para{\mathbf{\Phi}}^{(\infty)}_n(\lambda,x, \Theta_\infty),& |\lambda|>1,\\
 \para{\mathbf{\Phi}}_n^{(0)}(\lambda,x, \Theta_0),& |\lambda|<1.
 \end{cases}
\end{equation*}
This matrix function satisfies exactly the same jump conditions in the domains $|\lambda|>1$ and~${|\lambda|<1}$ as does $\mathbf{\Phi}_n(\lambda,x)$ itself, and it is also consistent with the asymptotics given in \eqref{eq:Psi-n-asymptotic-infinity}--\eqref{eq:Psi-n-asymptotic-zero} (note that $\mathbf{\Psi}_n(\lambda,x)=\mathbf{\Phi}_n(\lambda,x)$ for $|\lambda|$ sufficiently large or small). The parametrix has unit determinant, so the matrix quotient
\begin{equation*}
 \mathbf{Q}_n(\lambda,x):=\mathbf{\Phi}_n(\lambda,x)\para{\mathbf{\Phi}}_n(\lambda,x)^{-1}
\end{equation*}
is an analytic function of $\lambda$ except possibly on the jump contour $\Gamma$ shown in Figure~\ref{2} and on the unit circle, where there is a discontinuity in the definition of $\para{\mathbf{\Phi}}_n(\lambda,x)$. However, since the jumps of $\para{\mathbf{\Phi}}_n(\lambda,x)$ and $\mathbf{\Phi}_n(\lambda,x)$ agree on $\Gamma$, a Morera argument shows that $\mathbf{Q}_n(\lambda,x)$ is actually analytic both for $|\lambda|>1$ and for $0<|\lambda|<1$. The asymptotic behavior of the factors in $\mathbf{Q}_n(\lambda,x)$ as $\lambda\to 0$ then shows that any singularity of $\mathbf{Q}_n(\lambda,x)$ at the origin $\lambda=0$ is removable, and the asymptotic behavior of the same factors in the limit $\lambda\to\infty$ shows that $\mathbf{Q}_n(\lambda,x)\to\mathbb{I}$ as $\lambda\to\infty$.

$\mathbf{Q}_n(\lambda,x)$ is therefore characterized by its jump condition across the unit circle $|\lambda|=1$. Taking counterclockwise orientation for the circle, the jump condition for $\mathbf{Q}_n(\lambda,x)$ reads
\begin{equation*}
 \mathbf{Q}_{n,+}(\lambda,x)=\mathbf{Q}_{n,-}(\lambda,x)\para{\mathbf{\Phi}}_n^{(\infty)}(\lambda,x, \Theta_\infty)e_2^{\sigma_3}\para{\mathbf{\Phi}}_n^{(0)}(\lambda,x, \Theta_0)^{-1},\qquad |\lambda|=1.
\end{equation*}
Using Proposition~\ref{prop:symmetry}, the jump matrix can be written as
\begin{gather}
 \para{\mathbf{\Phi}}_n^{(\infty)}(\lambda,x, \Theta_\infty)e_2^{\sigma_3}\para{\mathbf{\Phi}}_n^{(0)}(\lambda,x, \Theta_0)^{-1}\nonumber\\
 \qquad =\para{\mathbf{\Phi}}^{(\infty)}_n(\lambda,x, \Theta_\infty) e_2^{\sigma_3}
 \begin{bmatrix}0 & -\beta_n\\ \beta_n^{-1} & 0\end{bmatrix}
 \para{\mathbf{\Phi}}^{(\infty)}_{-n}(-\lambda^{-1},x, \Theta_0)^{-1}e_0^{-3\sigma_3},\qquad |\lambda|=1.\label{eq:VQ-rewrite}
\end{gather}
We summarize by writing the Riemann--Hilbert problem for $\mathbf{Q}_n(\lambda,x)$.
\begin{rhp}\label{rhp:Q}
Fix generic monodromy parameters $( e_1, e_2)$, $n\in\mathbb{Z}$, and $x\in\mathbb{C}$. Seek a $2\times 2$ matrix function $\lambda\mapsto\mathbf{Q}_n(\lambda,x)$ with the following properties:
\begin{itemize}\itemsep=0pt\samepage
 \item Analyticity: $\mathbf{Q}_n(\lambda,x)$ is an analytic function of $\lambda$ for $|\lambda|\neq 1$.
 \item Jump condition:
 $\mathbf{Q}_n(\lambda,x)$ takes analytic boundary values on the unit circle from the interior and exterior, denoted $\mathbf{Q}_{n,+}(\lambda,x)$ and $\mathbf{Q}_{n,-}(\lambda,x)$ for $|\lambda|=1$ respectively, and they are related by
 \begin{equation*}
 %\label{Q-jump}\nonumber \mathbf{Q}_{n,+}(\lambda,x)=\mathbf{Q}_{n,-}(\lambda,x)\para{\mathbf{\Phi}}^{(\infty)}_n(\lambda,x, \Theta_\infty) e_2^{\sigma_3}
 \mathbf{Q}_{n,+}(\lambda,x)=\mathbf{Q}_{n,-}(\lambda,x)\breve{\mathbf{\Phi}}_n^{(\infty)}(\lambda,x,\Theta_\infty)e_2^{\sigma_3}
 \begin{bmatrix}0 & -\beta_n\\\beta_n^{-1} & 0\end{bmatrix}
 \para{\mathbf{\Phi}}^{(\infty)}_{-n}\bigl(-\lambda^{-1},x, \Theta_0\bigr)^{-1}e_0^{-3\sigma_3}.
 \end{equation*}
 \item Normalization:
 $\mathbf{Q}_n(\lambda,x)\to\mathbb{I}$ as $\lambda\to\infty$.
\end{itemize}
\end{rhp}
Henceforth, to avoid the notation becoming unwieldy, we understand that all quantities appearing with subscript $n$ are evaluated at parameter $\Theta_\infty$ while quantities appearing with subscript $-n$ are evaluated at parameter $\Theta_0$.

\subsection[The limit n to infinity]{The limit $\boldsymbol{n\to+\infty}$}
Having succeeded in removing the problematic jump conditions along rays emanating from~$0$,~$\infty$ in the $\lambda$ plane by defining $\mathbf{Q}_n(\lambda,x)$, we would next like to consider the limiting behavior of this problem as $n\to+\infty$ with $x=z/n$ and $z$ fixed. It is convenient to first renormalize~$\mathbf{Q}_n(\lambda,x)$, essentially by a transformation that diagonalizes the coefficient $\mu_n \mathbf{B}_n(x)\sigma_3\mathbf{B}_n(x)^{-1}$ of $\lambda^{-1}$ in the matrix of the Lax equation \eqref{Lax-system-lambda}.
In other words, in the jump condition for $\mathbf{Q}_n(\lambda,x)$ we prefer to replace \smash{$\para{\mathbf{\Phi}}_n^{(\infty)}(\lambda,x)$} with a suitable left-diagonal multiple of~\smash{$\mathbf{B}_n(x)^{-1}\para{\mathbf{\Phi}}_n^{(\infty)}(\lambda,x)$}. Observe that the coefficient $\mathbf{B}_n(x)$ is determined up to right-multiplication by a diagonal matrix by
\eqref{eq:zero-coefficient-matrix}, in which the second row of the matrix on the right-hand side is $(c_n(x),-a_n)=\big(\gamma_n x^{1-2\kappa_n},\mu_n ^{-1}\big(\kappa_n-\frac{1}{2}\big)\big)$, where we used \eqref{kappa-mu} and \eqref{c-of-x}. Indeed, the first column \smash{$\mathbf{b}_n^{(1)}(x)$} satisfies \smash{$\big(\gamma_n x^{1-2\kappa_n},\mu_n ^{-1}\big(\kappa_n-\frac{1}{2}\big)-1\big)\mathbf{b}_n^{(1)}(x)=0$} while the second column \smash{$\mathbf{b}_n^{(2)}(x)$} satisfies \smash{$\big(\gamma_n x^{1-2\kappa_n},\mu_n ^{-1}\big(\kappa_n-\frac{1}{2}\big)+1\big)\mathbf{b}_n^{(2)}(x)=0$}. By selecting specific constant factors for each column, we obtain a matrix $\mathbf{P}_n(x)$ differing from $\mathbf{B}_n(x)$ by right-multiplication by a diagonal matrix, and given explicitly by
\begin{equation*}
 \mathbf{P}_n(x):=\begin{bmatrix}\frac{1}{2}-\kappa_n+\mu_n & \frac{1}{2}-\kappa_n-\mu_n \\
 \mu_n \gamma_n x^{1-2\kappa_n} & \mu_n \gamma_n x^{1-2\kappa_n}\end{bmatrix},
\end{equation*}
in which the dependence on the index $n$ enters via \eqref{kappa-mu} and \eqref{Qinfty-gammainfty}. Then to get the desired modification of the jump matrix we set
\begin{equation*}
 \mathbf{R}_n(\lambda,x):=\begin{cases}
 \mathbf{P}_n(x)^{-1}\mathbf{Q}_n(\lambda,x)\mathbf{P}_n(x),& |\lambda|>1,\\
 \mathbf{P}_n(x)^{-1}\mathbf{Q}_n(\lambda,x)e_0^{3\sigma_3}\mathbf{P}_{-n}(x)d_n(x),& |\lambda|<1,
 \end{cases}
\end{equation*}
where $d_n(x)$ is a scalar satisfying
\begin{equation}
 d_n(x)^2 = \frac{\det(\mathbf{P}_{n}(x))}{\det(\mathbf{P}_{-n}(x))}.
 \label{eq:dn-squared}
\end{equation}
Then, $\mathbf{R}_n(\lambda,x)$ solves the following
Riemann--Hilbert problem.

\begin{rhp}
Fix generic monodromy parameters $( e_1, e_2)$, $n\in\mathbb{Z}$, and $x\in\mathbb{C}$. Seek a $2\times 2$ matrix function $\lambda\mapsto\mathbf{R}_n(\lambda,x)$ with the following properties:
\begin{itemize}\itemsep=0pt
 \item Analyticity: $\mathbf{R}_n(\lambda,x)$ is an analytic function of $\lambda$ for $|\lambda|\neq 1$.
 \item Jump condition:
 $\mathbf{R}_n(\lambda,x)$ takes analytic boundary values on the unit circle from the interior and exterior, denoted $\mathbf{R}_{n,+}(\lambda,x)$ and $\mathbf{R}_{n,-}(\lambda,x)$ for $|\lambda|=1$ respectively, and they are related by
 \begin{align}
 & \mathbf{R}_{n,+}(\lambda,x)= \mathbf{R}_{n,-}(\lambda,x)\mathbf{P}_n(x)^{-1}\para{\mathbf{\Phi}}_n^{(\infty)}(\lambda,x)e_2^{\sigma_3}\begin{bmatrix}0 & -\beta_n\\\beta_n^{-1} & 0\end{bmatrix}\nonumber\\
 &\hphantom{\mathbf{R}_{n,+}(\lambda,x)=}{} \times
 \para{\mathbf{\Phi}}_{-n}^{(\infty)}\bigl(-\lambda^{-1},x\bigr)^{-1}\mathbf{P}_{-n}(x)d_n(x)\label{eq:R-jump}
 \end{align}
 \item Normalization:
 $\mathbf{R}_n(\lambda,x)\to\mathbb{I}$ as $\lambda\to\infty$.
\end{itemize}
%\label{rhp:R}
\end{rhp}
The matrices $\mathbf{\Xi}_n^{(6)}(x)$ and $\mathbf{\Delta}_n^{(6)}(x)$ defined in \eqref{eq:T-n-definition} and \eqref{eq:Psi-n-asymptotic-zero}, respectively, can be expressed in terms of $\mathbf{R}_n(\lambda,x)$ as follows:
\begin{gather}
 \mathbf{\Xi}_n^{(6)}(x)=\mathbf{A}_n(x)+\mathbf{P}_n(x)\big[\lim_{\lambda\to\infty}\lambda(\mathbf{R}_n(\lambda,x)-\mathbb{I})\big]\mathbf{P}_n(x)^{-1}-\frac{1}{2}\ii x\sigma_3,\nonumber\\
 \mathbf{\Delta}_n^{(6)}(x)=\mathbf{P}_n(x)\mathbf{R}_n(0,x)\mathbf{P}_{-n}(x)^{-1}d_n(x)^{-1}e_0^{-3\sigma_3}.\label{eq:TU-R}
\end{gather}
Here, $\mathbf{A}_n(x)$ is the matrix coefficient defined in \eqref{eq:Phi-asymptotic-infinity}.

We now show that the jump matrix in \eqref{eq:R-jump} has explicit limits as $n\to +\infty$ along even or odd subsequences, with the convergence being uniform for $|\lambda|=1$ and bounded $z$ where $x=z/n$. To this end, we compute the asymptotic behavior of \smash{$\mathbf{P}_n\big(n^{-1}z\big)^{-1}\para{\mathbf{\Phi}}_n^{(\infty)}\big(\lambda,n^{-1}z\big)$} assuming that $|\lambda|=1$. The relevant formula for \smash{$\para{\mathbf{\Phi}}_n^{(\infty)}(\lambda,x)$} in this setting is \eqref{infinity-parametrix-form-origin}.
When $|\lambda|=1$,
\begin{gather*}
 \mathbf{P}_n(x)^{-1}\para{\mathbf{\Phi}}_n^{(\infty)}(\lambda,x)\\
 \qquad=\frac{1}{2\mu_n^2\gamma_nx^{1-2\kappa_n}}\begin{bmatrix}\mu_n\gamma_nx^{1-2\kappa_n} & -\frac{1}{2}+\mu_n+\kappa_n \vspace{1mm}\\-\mu_n\gamma_nx^{1-2\kappa_n} & \frac{1}{2}+\mu_n-\kappa_n\end{bmatrix}\mathbf{D}_nx^{\kappa_n\sigma_3}\lambda_{\lw}^{\sigma_3/2}\mathbf{M}(\ii x\lambda;\kappa_n,\mu_n)\mathbf{J}_n.
\end{gather*}
Using \eqref{H-define-2}, we see that
\begin{equation*}
 \mathbf{P}_n(x)^{-1}\para{\mathbf{\Phi}}_n^{(\infty)}(\lambda,x)=\frac{1}{2\mu_n^2 \gamma_n}x^{\kappa_n-\frac{1}{2}}\begin{bmatrix}
 \ii & -\frac{1}{2} + \mu_n+\kappa_n \vspace{1mm}\\-\ii & \frac{1}{2} + \mu_n-\kappa_n\end{bmatrix}x^{\sigma_3/2}\lambda_{\lw}^{\sigma_3/2}\mathbf{M}(\ii x\lambda;\kappa_n,\mu_n)\mathbf{J}_n.
\end{equation*}
Now, for $x>0$, the principal branch power $(-\ii x\lambda)^{\sigma_3/2}$ has the same domain of analyticity as~\smash{$\lambda_{\lw}^{\sigma_3/2}$}, and these two analytic functions are related by the identity \[\lambda_{\lw}^{\sigma_3/2}=x^{-\sigma_3/2}\ee^{\ii\pi\sigma_3/4}(-\ii x\lambda)^{\sigma_3/2}.\] Therefore,
\begin{gather}
 \mathbf{P}_n(x)^{-1}\para{\mathbf{\Phi}}_n^{(\infty)}(\lambda,x)\nonumber\\
 \qquad =\frac{1}{2\mu_n^2 \gamma_n}x^{\kappa_n-\frac{1}{2}}\begin{bmatrix}
 \ii & -\frac{1}{2} + \mu_n+\kappa_n \\-\ii & \frac{1}{2} + \mu_n-\kappa_n\end{bmatrix}\ee^{\ii\pi\sigma_3/4}(-\ii x\lambda)^{\sigma_3/2}\mathbf{M}(\ii x\lambda;\kappa_n,\mu_n)\mathbf{J}_n\nonumber\\
 \qquad=\frac{\ee^{-\ii\pi/4}}{2 \mu_n^2 \gamma_n}x^{\kappa_n-\frac{1}{2}}\begin{bmatrix}-1 & -\frac{1}{2} + \mu_n +\kappa_n \\1 & \frac{1}{2} + \mu_n-\kappa_n\end{bmatrix}(-Z)^{\sigma_3/2}\mathbf{M}(Z;\kappa_n,\mu_n)\mathbf{J}_n,\qquad Z=\ii x\lambda.
\label{eq:Psi-infty-hat}
\end{gather}
Now using \eqref{N-define}, we have
\begin{gather*}
 (-Z)^{\sigma_3/2}\mathbf{M}(Z;\kappa_n,\mu_n) \\
 \quad=
 \begin{bmatrix}
 \bigl(-\frac{1}{2} - \mu_n+\kappa_n\bigr)(-Z)^{-\frac{1}{2}}M_{1-\kappa_n,\mu_n}(-Z) & \bigl(-\frac{1}{2}+\mu_n+\kappa_n\bigr)(-Z)^{-\frac{1}{2}}M_{1-\kappa_n,-\mu_n}(-Z)\vspace{1mm}\\
 (-Z)^{-\frac{1}{2}}M_{-\kappa_n,\mu_n}(-Z) & (-Z)^{-\frac{1}{2}}M_{-\kappa_n,-\mu_n}(-Z)
 \end{bmatrix}.
\end{gather*}
The diagonal elements in \eqref{eq:Psi-infty-hat} can be simplified using the identity (see \cite[equation~(13.15.3)]{DLMF})
\begin{equation*}
 \big(\kappa-\mu-\tfrac{1}{2}\big)M_{\kappa-\frac{1}{2},\mu+\frac{1}{2}}(\diamond)+(1+2\mu)\diamond^{\frac{1}{2}}M_{\kappa,\mu}(\diamond) -\big(\kappa+\mu+\tfrac{1}{2}\big)M_{\kappa+\frac{1}{2},\mu+\frac{1}{2}}(\diamond)=0
\end{equation*}
replacing $\kappa\to\frac{1}{2}-\kappa_n$ and $\mu\to\mu_n-\frac{1}{2}$ for the $(1,1)$ entry and $\mu\to-\frac{1}{2}-\mu_n$ for the $(2,2)$ entry,
and the off-diagonal elements can be simplified using the identity (see \cite[equation~(13.15.4)]{DLMF})
\begin{equation*}
 2\mu M_{\kappa-\frac{1}{2},\mu-\frac{1}{2}}(\diamond)-2\mu M_{\kappa+\frac{1}{2},\mu-\frac{1}{2}}(\diamond)-\diamond^{\frac{1}{2}}M_{\kappa,\mu}(\diamond)=0
\end{equation*}
replacing $\kappa \to \frac{1}{2}-\kappa_n$ and $\mu\to\frac{1}{2}-\mu_n$ for the $(1,2)$ entry and $\mu\to\frac{1}{2}+\mu_n$ for the $(2,1)$ entry. The result is that
\begin{gather}
 \mathbf{P}_n(x)^{-1}\para{\mathbf{\Phi}}_n^{(\infty)}(\lambda,x)
 =\frac{\ee^{-\ii\pi/4}}{2 \mu_n^2 \gamma_n}x^{\kappa_n-\frac{1}{2}}\nonumber\\
 \qquad {}\times\begin{bmatrix}2\mu_n M_{\frac{1}{2}-\kappa_n,\mu_n-\frac{1}{2}}(-Z) & \left(\dfrac{\kappa_n + \mu_n - \frac{1}{2}}{1 - 2\mu_n}\right) M_{\frac{1}{2}-\kappa_n,\frac{1}{2}-\mu_n}(-Z) \\
 \left(\dfrac{\frac{1}{2} - \kappa_n + \mu_n}{1 + 2\mu_n}\right)M_{\frac{1}{2}-\kappa_n,\frac{1}{2}+\mu_n}(-Z) & 2\mu_n M_{\frac{1}{2}-\kappa_n,-\frac{1}{2}-\mu_n}(-Z)\end{bmatrix}\mathbf{J}_n,\nonumber\\ Z=\ii x\lambda.
 \label{eq:Psi-infty-hat-1}
\end{gather}
We will need the following result for the large $n$ limit of Whittaker functions appearing here, cf.~\cite[equation~(13.21.1)]{DLMF}.

\begin{Lemma}\label{lem:M-expansion}
Assume that $\mu$ is fixed with $2\mu\neq -1,-2,-3,\dots$, and let $f(\zeta;\mu)$ denote the entire function
\begin{equation}
 f(\zeta;\mu):=\sum_{s=0}^\infty \frac{(-\zeta)^s}{\Gamma(1+2\mu+s)s!}.
\label{eq:Bessel-f-define}
\end{equation}
Then the asymptotic formula
\begin{equation*}%\label{eq:M-asymptotic-formula}
 M_{\kappa,\mu}\left(\frac{\zeta}{\kappa}\right) = \Gamma(1+2\mu)\left(\frac{\zeta}{\kappa}\right)^{\mu+1/2}\big[f(\zeta;\mu) + \mathcal{O}\big(\kappa^{-2}\big)\big]
\end{equation*}
holds uniformly in the limit $\kappa\to\infty$ in any $($possibly complex$)$ direction under the assumption~${\zeta=\mathcal{O}(1)}$.
\end{Lemma}
\begin{proof}
We start from the formula \cite[equation~(13.14.6)]{DLMF} which holds under the indicated condition on $\mu$:
\begin{equation*}
 M_{\kappa,\mu}\left(\frac{\zeta}{\kappa}\right)=
 \Gamma(1+2\mu)\left(\frac{\zeta}{\kappa}\right)^{\mu+1/2}\ee^{-\zeta/(2\kappa)}\sum_{s=0}^\infty\frac{(-\zeta)^s}{\Gamma(1+2\mu+s)s!}\prod_{j=0}^{s-1}\left(1-\frac{\mu+1/2+j}{\kappa}\right).
 %\label{eq:M-expansion} \nonumber
\end{equation*}
Clearly, $\ee^{-\zeta/(2\kappa)}=1-\zeta/(2\kappa)+\mathcal{O}\big(\kappa^{-2}\big)$ as $\kappa\to\infty$ for $\zeta=\mathcal{O}(1)$, and the product in the summand has the expansion
\begin{gather}
 \prod_{j=0}^{s-1}\left(1-\frac{\mu+1/2+j}{\kappa}\right) = 1 -\frac{1}{\kappa}\bigg[\left(\mu+\frac{1}{2}\right)s +\frac{1}{2}s(s-1)\bigg] + \mathcal{O}\big(\kappa^{-2}s^s\big),\qquad\kappa\to\infty\!\!\!
\label{eq:uniform-estimate}
\end{gather}
uniformly for all indices $s$.
This follows from the Fredholm expansion formula
\begin{equation*}
 \prod_{j=0}^{s-1}(1+r_j) = 1+\sum_{j=0}^{s-1}r_j + \sum_{k=2}^{s}
 \mathop{\sum_{S\subset \mathbb{Z}_s}}_{|S|=k}\prod_{l\in S}r_l
\end{equation*}
and the estimate
\begin{equation*}
 \Biggl| \sum_{\substack{S\subset \mathbb{Z}_s \\ |S|=k}}\prod_{l\in S}r_l\Biggr|\le \binom{s}{k}R_s^k,\qquad R_s:=\max_{0\le l\le s-1}\{|r_l|\}.
\end{equation*}
Indeed, with $r_j=-\kappa^{-1}(\mu+\frac{1}{2}+j)$, we have
\begin{equation*}
 \sum_{j=0}^{s-1}r_j=-\frac{1}{\kappa}\sum_{j=0}^{s-1}\left(\mu+\frac{1}{2}+j\right) = -\frac{1}{\kappa}\left[\left(\mu+\frac{1}{2}\right)s + \frac{1}{2}s(s-1)\right],
\end{equation*}
and $R_s\le |\kappa|^{-1}\big(\big|\mu+\frac{1}{2}\big| + (s-1)\big)$. Therefore, $R_s^k\le |\kappa|^{-2}\big(\big|\mu+\frac{1}{2}\big|+(s-1)\big)^k$ holds for all $k\ge 2$ whenever $|\kappa|\ge 1$. Consequently,
\begin{gather*}
 \left|\prod_{j=0}^{s-1}\left(1-\frac{\mu+1/2+j}{\kappa}\right)-1+\right. \left.\frac{1}{\kappa}\left[\left(\mu+\frac{1}{2}\right)s+ \frac{1}{2}s(s-1)\right]\right|\\
 \qquad=\Biggl|\sum_{k=2}^s\mathop{\sum_{S\subset\mathbb{Z}_s}}_{|S|=k}\prod_{l\in S}r_l\Biggr|\le \frac{1}{|\kappa|^2}\sum_{k=2}^s\binom{s}{k}\left(\left|\mu+\frac{1}{2}\right|+(s-1)\right)^k1^{s-k}\\
 \qquad\le \frac{1}{|\kappa|^2}\sum_{k=0}^s\binom{s}{k}\left(\left|\mu+\frac{1}{2}\right|+(s-1)\right)^k1^{s-k}=\frac{\big(\big|\mu+\frac{1}{2}\big|+s\big)^s}{|\kappa|^2},
\end{gather*}
which proves \eqref{eq:uniform-estimate}.
Since the series
\begin{equation*}%\nonumber
 \sum_{s=0}^\infty \frac{(-\zeta)^ss^s}{\Gamma(1+2\mu_n+s)s!}
\end{equation*}
converges uniformly for $|\zeta|$ bounded, it follows that
\begin{equation*}%\label{eq:M-asymptotic-formula-intermediate}
 M_{\kappa,\mu}\left(\frac{\zeta}{\kappa}\right) = \Gamma(1+2\mu)\left(\frac{\zeta}{\kappa}\right)^{\mu+1/2}\left[f(\zeta;\mu) -\frac{1}{\kappa}g(\zeta;\mu) + \mathcal{O}\big(\kappa^{-2}\big)\right]
\end{equation*}
holds as $\kappa\to\infty$ in $\mathbb{C}$ with $\zeta=\mathcal{O}(1)$, where
\begin{equation*}%\nonumber
 g(\zeta;\mu):=\frac{1}{2}\zeta f(\zeta;\mu)+\left(\mu+\frac{1}{2}\right)\zeta f'(\zeta;\mu) +\frac{1}{2}\zeta^2f''(\zeta;\mu).
\end{equation*}
Now using the series \eqref{eq:Bessel-f-define} one checks that for all indicated values of $\mu$, $g(\zeta;\mu)$ vanishes identically, so the proof is complete.
\end{proof}

The series in \eqref{eq:Bessel-f-define} defines an entire function of $\zeta$ related to Bessel functions (see \cite[Chapter~10]{DLMF}) in the following way:
\begin{equation}
\label{f-bessel-j}
 f(\zeta,\mu)=\zeta^{-{\mu}{}}J_{2\mu}\big(2\sqrt{\zeta}\big).
\end{equation}
We apply Lemma~\ref{lem:M-expansion} to \eqref{eq:Psi-infty-hat-1} by taking $\zeta=-\big(\frac{1}{2}-\kappa_n\big)Z=-\big(\frac{1}{2}-\kappa_n\big)\ii x\lambda$. If $x=z/n$ and $z=\mathcal{O}(1)$, then using~\eqref{kappa-mu}, we see that $\zeta=-\frac{1}{2}\ii z\lambda + \mathcal{O}\big(n^{-1}\big)$ holds for $|\lambda|=1$. So, \eqref{eq:Psi-infty-hat-1}~becomes the statement that
\begin{gather*}
 \mathbf{P}_n(x)^{-1}\para{\mathbf{\Phi}}_n^{(\infty)}(\lambda,x)
\\
\quad=\frac{\ee^{-\ii\pi/4}}{2\mu_n^2\gamma_n}x^{\kappa_n-\frac{1}{2}}\Bigg(\frac{\rho_\infty}{2}\begin{bmatrix}
 \Gamma(2\mu_n+1) J_{2\mu_n-1}(\rho_\infty) & -\Gamma(1-2\mu_n) J_{1-2\mu_n}(\rho_\infty) \vspace{1mm}\\
 \Gamma(2\mu_n+1) J_{2\mu_n+1}(\rho_\infty) & -\Gamma(1-2\mu_n) J_{-1-2\mu_n}(\rho_\infty)
 \end{bmatrix}+\mathcal{O}\big(n^{-1}\big)\Bigg) \\
 \phantom{\quad=}{}\times\left(\dfrac{n}{2} \right)^{-\mu_n \sigma_3}
 \mathbf{J}_n
\end{gather*}
with
\begin{equation}
 \rho_\infty=\rho_\infty(\lambda,z):=(-2\ii z\lambda)^{1/2}\qquad\text{(principal branch)}
 \label{eq:rho-infty-branch}
\end{equation}
holds in the limit $n\to\infty$ with $x=n^{-1}z$ uniformly for $z=\mathcal{O}(1)$ and $|\lambda|=1$. Similarly, to study the parametrix near 0, it will be convenient to rewrite formula \eqref{f-bessel-j} in terms of modified Bessel functions:
\begin{equation*}
 f(\zeta, \mu) = (-\zeta)^{-\mu} I_{2\mu}\big(2 \sqrt{-\zeta} \big).
\end{equation*}

Replacing $n$ with $-n$, $\Theta_\infty$ with $\Theta_0$, and $\lambda$ with $-\lambda^{-1}$ and recalling that $\mu_{-n} = \mu_n$, gives in the same limit,
\begin{gather*}
 \mathbf{P}_{-n}(x)^{-1}\para{\mathbf{\Phi}}_{-n}^{(\infty)}\bigl(-\lambda^{-1},x\bigr)
\\
 \quad =\frac{\ee^{-\ii\pi/4}}{2\mu_n^2\gamma_{-n}}x^{\kappa_{-n}-\frac{1}{2}}\Bigg(\frac{\rho_0}{2}\begin{bmatrix}{\Gamma(2\mu_n+1)}{} I_{2\mu_n-1}(\rho_0) & {\Gamma(1-2\mu_n)}{} I_{1-2\mu_n}(\rho_0) \vspace{1mm}\\
 -{\Gamma(2\mu_n+1)}{}I_{2\mu_n+1}(\rho_0)&-{\Gamma(1-2\mu_n)}{}I_{-2\mu_n-1}(\rho_0)
 \end{bmatrix} +\mathcal{O}\big(n^{-1}\big)\Bigg)\\
\phantom{ \quad =}{} \times \left(\frac{n}{2}\right)^{-\mu_n \sigma_3}
 \mathbf{J}_{-n}
\end{gather*}
with
\begin{equation}
\rho_0=\rho_0(\lambda,z):=\big(2\ii z\lambda^{-1}\big)^{1/2}\qquad\text{(principal branch)}.
\label{eq:rho-zero-branch}
\end{equation}
The jump matrix in \eqref{eq:R-jump} therefore reads
\begin{gather}
 \mathbf{R}_{n,-}\big(\lambda,n^{-1}z\big)^{-1}
 \mathbf{R}_{n,+}\big(\lambda,n^{-1}z\big) \nonumber\\
 = \frac{\rho_\infty}{2}\begin{bmatrix}
 \Gamma(2\mu_n+1) J_{2\mu_n-1}(\rho_\infty)+ \mathcal{O}\big(n^{-1}\big) & -\Gamma(1-2\mu_n) J_{1-2\mu_n}(\rho_\infty)+ \mathcal{O}\big(n^{-1}\big) \vspace{1mm}\\
 \Gamma(2\mu_n+1) J_{2\mu_n+1}(\rho_\infty)+ \mathcal{O}\big(n^{-1}\big) & -\Gamma(1-2\mu_n) J_{-1-2\mu_n}(\rho_\infty)+ \mathcal{O}\big(n^{-1}\big)
 \end{bmatrix}\nonumber\\
\phantom{ =}{} \times{} \frac{\gamma_{-n}}{\gamma_n}x^{\kappa_n-\kappa_{-n}}d_n(x)\left(\frac{n}{2}\right)^{-\mu_n\sigma_3}\mathbf{J}_ne_2^{\sigma_3}\begin{bmatrix}0 & -\beta_n\\\beta_n^{-1} & 0\end{bmatrix}\mathbf{J}_{-n}^{-1}\left(\frac{n}{2}\right)^{\mu_n\sigma_3}\frac{2}{\rho_0}\nonumber
 \\
\phantom{ =}{}\times \begin{bmatrix}{\Gamma(2\mu_n+1)}{} I_{2\mu_n-1}(\rho_0)+ \mathcal{O}\big(n^{-1}\big) & {\Gamma(1-2\mu_n)}{} I_{1-2\mu_n}(\rho_0)+ \mathcal{O}\big(n^{-1}\big) \vspace{1mm}\\
 -{\Gamma(2\mu_n+1)}{}I_{2\mu_n+1}(\rho_0)+ \mathcal{O}\big(n^{-1}\big)&-{\Gamma(1-2\mu_n)}{}I_{-2\mu_n-1}(\rho_0)+ \mathcal{O}\big(n^{-1}\big)
 \end{bmatrix}^{-1}.
 \label{eq:E-hat-jump-1}
\end{gather}
Expanding \eqref{eq:dn-squared} for large $n>0$ gives
\begin{align*}
& d_n(x)^2= \frac{\gamma_n}{\gamma_{-n}}x^{2\kappa_{-n}-2\kappa_n} = \dfrac{e_\infty^5}{e_0^5} \dfrac{e_0^2 - e_1^2}{e_\infty^2 - e_1^2} \dfrac{\Gamma\big( \frac{1}{2} - \kappa_{-n} -\mu_n\big) \Gamma\big( \frac{1}{2} - \kappa_{-n} + \mu_n\big)}{\Gamma\big( \frac{1}{2} - \kappa_{n} -\mu_n\big) \Gamma\big( \frac{1}{2} - \kappa_{n} +\mu_n\big)} x^{2\kappa_{-n}-2\kappa_n}\\
&\hphantom{d_n(x)^2}{} = \dfrac{e_\infty^5 e_1^2}{e_0^3\big(1 - e_1^2 e_0^2\big)\big(e_\infty^2 - e_1^2\big)} \\
&\hphantom{d_n(x)^2=}{} \times\dfrac{ 4\pi^2x^{2\kappa_{-n}-2\kappa_n}}{\Gamma \big( \frac{1}{2} + \kappa_{-n}- \mu_n \big) \Gamma \big( \frac{1}{2} + \kappa_{-n} + \mu_n \big) \Gamma \big( \frac{1}{2} - \kappa_n + \mu_n\big) \Gamma \big( \frac{1}{2} - \kappa_n - \mu_n \big)} \\
&\hphantom{d_n(x)^2}{}=
 4\dfrac{e_\infty^4}{e_0^4}\dfrac{e_0e_\infty e_1^2}{\big(1 - e_1^2 e_0^2\big)\big(e_\infty^2 - e_1^2\big)} x^{2\kappa_{-n}-2\kappa_n}n^{-2n-\Theta_0+\Theta_\infty}\ee^{2n}2^{2n+\Theta_0-\Theta_\infty-2}\big(1+\mathcal{O}\big(n^{-1}\big)\big), \\
& n \to +\infty.
\end{align*}
We now properly define $d_n(x)$ for large $n$ by selecting a definite value for the square root of
\begin{equation}
\label{eq:V-root}
%\dfrac{e_1 e_\infty^{5/2}}{e_0^{3/2}} \dfrac{2}{(1 - e_1^2 e_0^2)}
\sqrt{\dfrac{1 - e_1^2 e_0^2}{e_\infty^2 - e_1^2} }
\end{equation}
after which $d_n(x)$ has the asymptotic expansion
\begin{align*}
&d_n(x) = \dfrac{e_1 e_\infty^{5/2}}{e_0^{3/2}} \dfrac{1}{\big(1 - e_1^2 e_0^2\big)}\sqrt{\dfrac{1 - e_1^2 e_0^2}{e_\infty^2 - e_1^2} } x^{\kappa_{-n}-\kappa_n} n^{-n+\frac{1}{2}(-\Theta_0+\Theta_\infty)}\ee^{n}2^{n + \frac{1}{2}(\Theta_0-\Theta_\infty)}\\
&\hphantom{d_n(x) =}{}
\times\big(1+\mathcal{O}\big(n^{-1}\big)\big), \qquad n \to +\infty.
\end{align*}
Then, by definition, we have
\begin{align*}
\dfrac{\gamma_{-n}}{\gamma_n} d_n(x) x^{\kappa_n - \kappa_{-n}} ={}& \dfrac{e_0^{3/2}}{e_1 e_\infty^{5/2}} {\big(1 - e_1^2 e_0^2\big)}\left(\sqrt{\dfrac{1 - e_1^2 e_0^2}{e_\infty^2 - e_1^2} }\right)^{-1}n^{n-\frac{1}{2}(-\Theta_0+\Theta_\infty)}\ee^{-n}\\&
\times2^{-n - \frac{1}{2}(\Theta_0-\Theta_\infty)}\big(1+\mathcal{O}\big(n^{-1}\big)\big),\qquad n \to +\infty.
\end{align*}
Furthermore, using identities \eqref{eq:Gamma-reflection}, \eqref{eq:cosine-identity-1}, and Stirling's formula yields
\begin{align*}
 & -\left(\frac{2}{n}\right)^{2\mu_n}\beta_n\frac{J_{n,11}}{J_{-n,22}} = -\left(\frac{2}{n}\right)^{2\mu_n} \frac{e_\infty^{7/2} \big(e_1^2-e_0^2\big) \Gamma (-2 \mu_n ) \Gamma \big(\frac{1}{2}-\kappa_{-n}+\mu_n \big)}{e_0^{3/2} \big(e_0^2 e_1^2-1\big) \Gamma (2 \mu_n ) \Gamma \big(\frac{1}{2}-\kappa_{n}-\mu_n \big)} \\
 &\hphantom{-\left(\frac{2}{n}\right)^{2\mu_n}\beta_n\frac{J_{n,11}}{J_{-n,22}}}{}
 = \frac{e_\infty^{7/2}}{e_0^{3/2}} \dfrac{\ii e_0e_1}{\big(1- e_0^2 e_1^2\big) } \dfrac{\Gamma (-2 \mu_n ) }{ \Gamma (2 \mu_n ) } n^{-n+\frac{1}{2}(-\Theta_0+\Theta_\infty)}\ee^n \\
 &\hphantom{-\left(\frac{2}{n}\right)^{2\mu_n}\beta_n\frac{J_{n,11}}{J_{-n,22}}=}{}
 \times 2^{n + \frac{1}{2}(\Theta_0-\Theta_\infty)}\big(1+\mathcal{O}\big(n^{-1}\big)\big),\qquad n\to+\infty,
\end{align*}
and similarly,
\begin{align*}
&\left(\frac{n}{2}\right)^{2\mu_n}\frac{J_{n,22}}{\beta_nJ_{-n,11}} = \left(\frac{n}{2}\right)^{2\mu_n} \frac{e_\infty^{3/2}\big(e_0^2 e_1^2-1\big) \Gamma (2 \mu_n ) \Gamma \big(\frac{1}{2}-\kappa_{-n}-\mu_n \big)}{e_0^{7/2} \big(e_1^2-e_\infty^2\big) \Gamma (-2 \mu_n ) \Gamma \big(\frac{1}{2}-\kappa_n+\mu_n \big)}\\
&\hphantom{\left(\frac{n}{2}\right)^{2\mu_n}\frac{J_{n,22}}{\beta_nJ_{-n,11}}}{}
= \frac{e_\infty^{3/2}}{e_0^{7/2}} \dfrac{\ii e_0e_1}{\big(e_\infty^2 - e_1^2\big)}\dfrac{ \Gamma (2 \mu_n ) }{ \Gamma (-2 \mu_n )}n^{-n+\frac{1}{2}(-\Theta_0+\Theta_\infty)}\ee^n \\
&\hphantom{\left(\frac{n}{2}\right)^{2\mu_n}\frac{J_{n,22}}{\beta_nJ_{-n,11}}=}{}
\times 2^{n + \frac{1}{2}(\Theta_0-\Theta_\infty)}\big(1+\mathcal{O}\big(n^{-1}\big)\big),\qquad n\to+\infty.
\end{align*}
Therefore, the central factor on the right-hand side of \eqref{eq:E-hat-jump-1} satisfies
\begin{gather*}
\frac{\gamma_{-n}}{\gamma_n}x^{\kappa_n-\kappa_{-n}}d_n(x)\left(\frac{n}{2}\right)^{-\mu_n\sigma_3}\mathbf{J}_ne_2^{\sigma_3}\begin{bmatrix}0 & -\beta_n\\\beta_n^{-1} & 0\end{bmatrix}\mathbf{J}_{-n}^{-1}\left(\frac{n}{2}\right)^{\mu_n\sigma_3}\\
\qquad = \begin{bmatrix} 0 & -\dfrac{e_0 e_2 e_\infty}{\ii}\dfrac{ \Gamma (-2 \mu_n )}{ \Gamma (2 \mu_n )} \left(\sqrt{\dfrac{1 - e_0^2 e_1^2}{e_\infty^2-e_1^2 } }\right)^{-1} \\
 \dfrac{ \ii }{e_0 e_2 e_\infty} \dfrac{\Gamma (2 \mu_n )}{ \Gamma (-2 \mu_n )} \sqrt{\dfrac{1 - e_0^2 e_1^2}{e_\infty^2-e_1^2 }} & 0
\end{bmatrix} + \mathcal{O}\big(n^{-1}\big).
\end{gather*}
The leading term is independent of $n\pmod{2}$ and has unit determinant. This proves the following.
\begin{prop}
Define the constant matrix which depends only on the even/odd parity of $n$ via $\mu_n$, $e_1$, $e_0^2$, and $e_\infty^2$:
\begin{equation}
 \mathbf{V}^\mathrm{even/odd}:=\begin{bmatrix} 0 & -\dfrac{e_0 e_2 e_\infty}{\ii} \left(\sqrt{\dfrac{1 - e_0^2 e_1^2}{e_\infty^2-e_1^2 } }\right)^{-1} \\
 \dfrac{ \ii }{e_0 e_2 e_\infty} \sqrt{\dfrac{1 - e_0^2 e_1^2}{e_\infty^2-e_1^2 }} & 0
\end{bmatrix}.
 % \nonumber
 \label{eq:Vm-even}
\end{equation}
Then the following asymptotic formula holds uniformly for $|\lambda|=1$ and $z$ bounded:
\begin{align*}
& \mathbf{R}_{n,-}(\lambda,z/n)^{-1}\mathbf{R}_{n,+}(\lambda,z/n)= {\rho_\infty}\begin{bmatrix}J_{2\mu_n - 1}(\rho_\infty) & -J_{1- 2\mu_n}(\rho_\infty)\\
 J_{2\mu_n +1}(\rho_\infty) & -J_{-1-2\mu_n}(\rho_\infty)\end{bmatrix}\cdot\mathbf{V}^\mathrm{even/odd}\\
&\hphantom{\mathbf{R}_{n,-}(\lambda,z/n)^{-1}\mathbf{R}_{n,+}(\lambda,z/n)=}{}
\times
 \dfrac{1}{\rho_0}\begin{bmatrix}I_{2\mu_n - 1}(\rho_0) & I_{1- 2\mu_n}(\rho_0)\\
 -I_{2\mu_n +1}(\rho_0) & -I_{-1-2\mu_n}(\rho_0)\end{bmatrix}^{-1}+\mathcal{O}\big(n^{-1}\big),
\end{align*}
as $n\to\infty$ along even/odd subsequences.
\label{prop:limiting-jump}
\end{prop}
Proposition~\ref{prop:limiting-jump} suggests defining the following limiting Riemann--Hilbert problem.
\begin{rhp}[limiting problem, even/odd subsequences of $n$]
Fix generic monodromy parameters $( e_1, e_2)$, and $z\in\mathbb{C}$ with $|{\arg}(z)|<\pi$. Seek a $2\times 2$ matrix function $\lambda\mapsto\hat{\mathbf{R}}^\mathrm{even/odd}(\lambda,z)$ with the following properties:
\begin{itemize}\itemsep=0pt
 \item Analyticity: $\hat{\mathbf{R}}^\mathrm{even/odd}(\lambda,z)$ is an analytic function of $\lambda$ for $|\lambda|\neq 1$.
 \item Jump condition:
 $\hat{\mathbf{R}}^\mathrm{even/odd}(\lambda,z)$ takes analytic boundary values on the unit circle from the interior and exterior, denoted $\hat{\mathbf{R}}^\mathrm{even/odd}_+(\lambda,z)$ and $\hat{\mathbf{R}}^\mathrm{even/odd}_-(\lambda,z)$ for $|\lambda|=1$ respectively, and they are related by
 \begin{align}
 &\hat{\mathbf{R}}^\mathrm{even/odd}_+(\lambda,z)=\hat{\mathbf{R}}^\mathrm{even/odd}_-(\lambda,z) {\rho_\infty}\begin{bmatrix}J_{2\mu_n - 1}(\rho_\infty) & -J_{1- 2\mu_n}(\rho_\infty)\\
 J_{2\mu_n +1}(\rho_\infty) & -J_{-1-2\mu_n}(\rho_\infty)\end{bmatrix}\nonumber\\
 & \hphantom{\hat{\mathbf{R}}^\mathrm{even/odd}_+(\lambda,z)=}{}
 \times\mathbf{V}^\mathrm{even/odd}
 \cdot \dfrac{1}{\rho_0}\begin{bmatrix}I_{2\mu_n - 1}(\rho_0) & I_{1- 2\mu_n}(\rho_0)\\
 -I_{2\mu_n +1}(\rho_0) & -I_{-1-2\mu_n}(\rho_0)\end{bmatrix}^{-1}.\label{r-hat-jump}
 \end{align}
 \item Normalization:
 $\hat{\mathbf{R}}^\mathrm{even/odd}(\lambda,z)\to\mathbb{I}$ as $\lambda\to\infty$.
\end{itemize}
\label{rhp:Rhat-even/odd}
\end{rhp}
Note that the Bessel functions $J_\nu(\rho_\infty)$ and $I_\nu(\rho_0)$ appearing in the jump matrix in \eqref{r-hat-jump} are analytic on the unit circle $|\lambda|=1$ except at the point $\lambda=\lambda_\mathrm{c}:=-\ii\ee^{-\ii\arg(z)}$. However, from the identities
\[
\left.J_\nu(\rho_\infty)\right|_{\lambda=\lambda_\mathrm{c}\ee^{-\ii 0}}=\ee^{\ii\pi\nu}\left.J_\nu(\rho_\infty)\right|_{\lambda=\lambda_\mathrm{c}\ee^{\ii 0}} \qquad \text{and} \qquad
\left.I_\nu(\rho_0)\right|_{\lambda=\lambda_\mathrm{c}\ee^{-\ii 0}}=\ee^{-\ii\pi\nu}\left.I_\nu(\rho_0)\right|_{\lambda=\lambda_\mathrm{c}\ee^{\ii 0}}
\]
and the fact that the indices $\nu$ in each column of the Bessel matrix factors in \eqref{r-hat-jump} differ by $2$, combined with the fact that $\mathbf{V}^\mathrm{even/odd}$ is an off-diagonal matrix, one sees easily that
\[
\begin{bmatrix}J_{2\mu_n - 1}(\rho_\infty) & -J_{1- 2\mu_n}(\rho_\infty)\\
 J_{2\mu_n +1}(\rho_\infty) & -J_{-1-2\mu_n}(\rho_\infty)\end{bmatrix}\cdot\mathbf{V}^\mathrm{even/odd}
 \cdot \begin{bmatrix}I_{2\mu_n - 1}(\rho_0) & I_{1- 2\mu_n}(\rho_0)\\
 -I_{2\mu_n +1}(\rho_0) & -I_{-1-2\mu_n}(\rho_0)\end{bmatrix}^{-1}
\]
is continuous at $\lambda=\lambda_\mathrm{c}$ and hence is an analytic function of $\lambda$ on the unit circle. The scalar factor $\rho_\infty/\rho_0$ is also analytic for $|\lambda|=1$, and therefore the jump matrix in \eqref{r-hat-jump} is an analytic function of $\lambda$ when $|\lambda|=1$. At this stage, the existence of a matrix function $\hat{\mathbf{R}}^\mathrm{even/odd}(\lambda, z)$ satisfying Riemann--Hilbert Problem \ref{rhp:Rhat-even/odd} is not clear. However, it turns out that there exists a~discrete set $\Sigma^{\mathrm{even/odd}} \subset \C$ such that for $z \in \C \setminus \Sigma^{\mathrm{even/odd}}$, such a matrix does exist and is in fact a meromorphic function of $z$, see Section \ref{sec:suleimanov-solution-connection} below.

\begin{Lemma}
\label{lemma:limit}
Let $(e_1, e_2)$ be generic monodromy parameters and take $z \in \C \setminus \Sigma^\mathrm{even/odd}$. Then,
\begin{align}
 \mathop{\lim_{n\to\infty}}_{\text{$n$ even/odd}}u_n\big(n^{-1}z;m\big)={}&-\dfrac{8 \mu_n^2}{z}\big[\hat{R}_{11}^\mathrm{even/odd}(0,z)+\hat{R}_{21}^\mathrm{even/odd}(0,z)\nonumber\\
 &-\hat{R}_{12}^\mathrm{even/odd}(0,z)+\hat{R}_{22}^\mathrm{even/odd}(0,z)\big]^{-2}.\label{eq:un-Rhat}
\end{align}
\end{Lemma}
\begin{proof}

Noting that $\hat{\mathbf{R}}^\mathrm{even/odd}(\lambda,z)$ necessarily has unit determinant, we form the matrix quotient
\begin{equation*}
 \mathbf{E}_n(\lambda,z):=\mathbf{R}_n\big(\lambda,n^{-1}z\big)\hat{\mathbf{R}}^\mathrm{even/odd}(\lambda,z)^{-1},\qquad |\lambda|\neq 1.
\end{equation*}
Clearly, $\mathbf{E}_n(\lambda,z)$ is analytic as a function of $\lambda$ in the domain of definition, and for each fixed $n$ it tends to $\mathbb{I}$ as $\lambda\to\infty$ as this is true for both $\mathbf{R}_n\big(\lambda,n^{-1}z\big)$ and $\hat{\mathbf{R}}^\mathrm{even/odd}(\lambda,z)$. Across the unit circle, the boundary values of $\mathbf{E}_n(\lambda,z)$ are related by
\begin{align*}
 &\mathbf{E}_{n,+}(\lambda,z)= \mathbf{E}_{n,-}(\lambda,z)\hat{\mathbf{R}}_-^\mathrm{even/odd}(\lambda,z)\big[\mathbf{R}_{n,-}\big(\lambda,n^{-1}z\big)^{-1}\mathbf{R}_{n,+}\big(\lambda,n^{-1}z\big)\big]\\
 &\hphantom{\mathbf{E}_{n,+}(\lambda,z)=}{}
 \times \big[\hat{\mathbf{R}}^\mathrm{even/odd}_-(\lambda,z)^{-1}\hat{\mathbf{R}}^\mathrm{even/odd}_+(\lambda,z)\big]^{-1}
 \hat{\mathbf{R}}^\mathrm{even/odd}_-(\lambda,z)^{-1},\qquad |\lambda|=1.
\end{align*}
Thus, the jump matrix for $\mathbf{E}_n(\lambda,z)$ is the conjugation, by a unit-determinant matrix function of $\lambda$ independent of $n$, of the matrix ratio of the jump matrices for $\mathbf{R}_n\big(\lambda,n^{-1}z\big)$ and for \smash{$\hat{\mathbf{R}}^\mathrm{even/odd}(\lambda,z)$}. But by Proposition~\ref{prop:limiting-jump}, the latter ratio is $\mathbb{I}+\mathcal{O}\big(n^{-1}\big)$ uniformly on the unit circle as $n\to\infty$ along even or odd subsequences. The conjugating factors exist and are uniformly bounded for $z$ in compact subsets of $\C\setminus \Sigma(y_1, y_2, y_3)$. %{\color{magenta}Here the conjugation amplifies the error if the conjugating factors can grow without bound, which should be the locus $\Sigma(y_1,y_2,y_3)$.}
It follows that in this limit, $\mathbf{E}_{n,+}(\lambda,z)=\mathbf{E}_{n,-}(\lambda,z)\big(\mathbb{I}+\mathcal{O}\big(n^{-1}\big)\big)$ uniformly for $|\lambda|=1$ and $z$ in compact subsets of~$\C\setminus \Sigma(y_1, y_2, y_3)$ as $n\to\infty$. By standard small-norm theory, $\mathbf{E}_n(\lambda,z)$ exists for large enough even or odd $n$, and tends to the identity as $n\to\infty$, in particular in the sense that
\begin{equation*}
 \lim_{\lambda\to\infty}\lambda(\mathbf{E}_n(\lambda,z)-\mathbb{I})\to\mathbf{0}\qquad\text{and}\qquad\mathbf{E}_n(0,z)\to\mathbb{I}
\end{equation*}
as $n\to\infty$ along even/odd subsequences. By the definition of $\mathbf{E}_n(\lambda,z)$ it follows that in the same limit
\begin{gather*}
 \lim_{\lambda\to\infty}\lambda\big(\mathbf{R}_n\big(\lambda,n^{-1}z\big)-\mathbb{I}\big)\to\lim_{\lambda\to\infty}\lambda\big(\hat{\mathbf{R}}^\mathrm{even/odd}(\lambda,z)-\mathbb{I}\big), \qquad \text{and}\\
 \mathbf{R}_n\big(0,n^{-1}z\big)\to\hat{\mathbf{R}}^\mathrm{even/odd}(0,z).
\end{gather*}
Combining \eqref{eq:u-n-recover} with \eqref{eq:TU-R} then shows \eqref{eq:un-Rhat}.
Partly, this works because the dominant term in $\Xi_{n,12}^{(6)}\big(n^{-1}z;m\big)$ is $A_{n,12}\big(n^{-1}z\big)$.
\end{proof}

\subsection{Transformations of the limiting Riemann--Hilbert problem} In this section we transform Riemann--Hilbert Problem \ref{rhp:Rhat-even/odd} to match the form of Riemann--Hilbert Problem \ref{rhp:D8}. To this end, using \cite[equations~(10.4.4) and (10.4.6)]{DLMF} to express the Bessel function~$J_\nu(\diamond)$ in terms of the Hankel functions \smash{$H_\nu^{(1)}(\diamond)$}, \smash{$H_\nu^{(2)}(\diamond)$} and the relations \cite[equation~(10.4.4) and~(10.4.6)]{DLMF},
\begin{align*}
 H_{-\nu}^{(1)}(\diamond) &= \ee^{\pi \ii \nu}H_{\nu}^{(1)}(\diamond) , \qquad
 H_{-\nu}^{(2)}(\diamond) = \ee^{-\pi \ii \nu}H_{\nu}^{(2)}(\diamond) ,
\end{align*}
we arrive at the identity:
\begin{gather}
 \rho_\infty \begin{bmatrix}J_{2\mu_n - 1}(\rho_\infty) & -J_{1- 2\mu_n}(\rho_\infty)\\
 J_{2\mu_n +1}(\rho_\infty) & -J_{-1-2\mu_n}(\rho_\infty)\end{bmatrix}\nonumber \\
 \qquad= \dfrac{\rho_\infty}{2} \begin{bmatrix}H_{2\mu_n - 1}^{(1)}(\rho_\infty) & H_{1 -2\mu_n}^{(2)}(\rho_\infty) \\ H^{(1)}_{2\mu_n + 1}(\rho_\infty) & H_{-1 -2\mu_n}^{(2)}(\rho_\infty) \end{bmatrix} \begin{bmatrix} 1 & \ee^{2\pi \ii \mu_n} \\ -\ee^{2\pi \ii \mu_n} & -1 \end{bmatrix}. \label{J-to-H}
\end{gather}
To obtain appropriate asymptotic formul\ae\ for the matrix on the right-hand side of \eqref{J-to-H}, we first apply the identity \cite[equation~(10.6.1)]{DLMF}
\begin{equation*}
 H^{(k)}_{\nu-1}(\diamond) + H^{(k)}_{\nu + 1}(\diamond) = \dfrac{2\nu}{\diamond} H^{(k)}_\nu(\diamond), \qquad k = 1, 2,
\end{equation*}
which gives
\begin{gather}
 \rho_\infty^{-\sigma_3/2} \begin{bmatrix} 1 & 0 \\ 1 & 1\end{bmatrix} \dfrac{\rho_\infty}{2} \begin{bmatrix}H_{2\mu_n - 1}^{(1)}(\rho_\infty) & H_{1 -2\mu_n}^{(2)}(\rho_\infty) \vspace{1mm}\\ H^{(1)}_{2\mu_n + 1}(\rho_\infty) & H_{-1 -2\mu_n}^{(2)}(\rho_\infty) \end{bmatrix}\nonumber\\
 \qquad= \dfrac{\sqrt{\rho_\infty}}{2} \begin{bmatrix} H_{2\mu_n-1}^{(1)}(\rho_\infty) & H_{1-2\mu_n}^{(2)}(\rho_\infty)\vspace{1mm}\\4\mu_n H_{2\mu_n}^{(1)}(\rho_\infty) & -4\mu_n H_{-2\mu_n}^{(2)}(\rho_\infty)\end{bmatrix}.\label{H-inf-identity}
\end{gather}
The matrix on the right-hand side is amenable to asymptotic analysis as $\rho_\infty \to \infty$; using the asymptotics of Hankel functions \cite[equations~(10.17.5) and (10.17.6)]{DLMF} and \eqref{H-inf-identity} yields
\begin{gather}
 \dfrac{\rho_\infty}{2} \begin{bmatrix}H_{2\mu_n - 1}^{(1)}(\rho_\infty) & H_{1 -2\mu_n}^{(2)}(\rho_\infty) \vspace{1mm}\\ H^{(1)}_{2\mu_n + 1}(\rho_\infty) & H_{-1 -2\mu_n}^{(2)}(\rho_\infty) \end{bmatrix}\nonumber\\
 \qquad= \begin{bmatrix} 1 & \dfrac{1}{2}-\dfrac{\mu_n}{2}-\dfrac{3}{32 \mu_n} \vspace{1mm}\\ -1 & \dfrac{1}{2} + \dfrac{\mu_n}{2}+\dfrac{3}{32 \mu_n}
\end{bmatrix} \Biggl( \mathbb{I}+ \dfrac{1}{128 \rho_\infty^2}  \nonumber\\
 \phantom{\qquad=}{}\times \begin{bmatrix} \big(16\mu_n^2 - 9\big)\big(16\mu_n^2 - 1\big) & \dfrac{\big(16 \mu_n^2 - 13\big)\big(16\mu_n^2 - 9\big)\big(16\mu_n^2 - 1\big)}{48\mu_n} \vspace{1mm}\\ 64\mu_n\big(16\mu_n^2 - 1\big) & - \big(16\mu_n^2 - 9\big)\big(16\mu_n^2 - 1\big)\end{bmatrix} + \mathcal{O}\big(\rho_\infty^{-4}\big)\Biggr) \rho_\infty^{\sigma_3/2} \nonumber\\
\phantom{\qquad=}{}\times (2\sqrt{\mu_n})^{\mathbb{I}-\sigma_3}{\dfrac{\ee^{-\pi \ii \mu_n}}{\sqrt{2\pi}}} {}{} \ee^{\pi \ii\sigma_3/4}\begin{bmatrix} 1 & \ii\\ 1 &-\ii \end{bmatrix} \ee^{\ii \rho_\infty \sigma_3}, \qquad \arg(\rho_\infty) \in (-\pi, \pi).\label{h-infty-asymptotics}
\end{gather}
We turn to analogously treating the final factor of the jump of $\hat{\mathbf{R}}^{\mathrm{even/odd}}(\lambda,z)$; using \cite[equation~(10.27.7)]{DLMF} and the above relations, we have
\begin{gather*}
% \label{I-to-H}
{\rho_0} \begin{bmatrix}I_{2\mu_n - 1}(\rho_0) & I_{1- 2\mu_n}(\rho_0)\vspace{1mm}\\
 -I_{2\mu_n +1}(\rho_0) & -I_{-1-2\mu_n}(\rho_0)\end{bmatrix} \\
 \qquad= \dfrac{\ee^{-\pi \ii/2}\rho_0}{2} \begin{bmatrix} H^{(1)}_{2\mu_n-1}\big(\ee^{-\pi \ii/2} \rho_0\big) & H^{(2)}_{1-2\mu_n}\big(\ee^{-\pi \ii/2} \rho_0\big)\vspace{1mm}\\ H^{(1)}_{1+2\mu_n}\big(\ee^{-\pi \ii/2} \rho_0\big) & H^{(2)}_{-1-2\mu_n}\big(\ee^{-\pi \ii/2} \rho_0\big)\end{bmatrix} \begin{bmatrix} \ee^{\pi \ii \mu_n}& \ee^{\pi \ii \mu_n}\\ - \ee^{3\pi \ii \mu_n} & - \ee^{-\pi \ii \mu_n} \end{bmatrix}.
\end{gather*}
This allows us to find the following large-$\rho_0$ asymptotics:
\begin{gather}
\dfrac{\ee^{-\pi \ii /2}\rho_0}{2} \begin{bmatrix} H^{(1)}_{2\mu_n-1}\big(\ee^{-\pi \ii/2} \rho_0\big) & H^{(2)}_{1 -2\mu_n}\big(\ee^{-\pi \ii/2} \rho_0\big)\vspace{1mm}\\ H^{(1)}_{2\mu_n + 1}\big(\ee^{-\pi \ii/2} \rho_0\big) & H^{(2)}_{-1-2\mu_n}\big(\ee^{-\pi \ii/2} \rho_0\big)\end{bmatrix} \nonumber\\
\qquad= \begin{bmatrix} 1 & \dfrac{1}{2} - \dfrac{\mu_n}{2} - \dfrac{3}{32\mu_n} \vspace{1mm}\\ -1 & \dfrac{1}{2} + \dfrac{\mu_n}{2} + \dfrac{3}{32\mu_n} \end{bmatrix} \Biggl( \mathbb{I} - \dfrac{1}{128 \rho_0^2}  \nonumber\\
\phantom{\qquad=}{}  \times\begin{bmatrix} \big(16\mu_n^2 - 9\big)\big(16\mu_n^2 - 1\big) & \dfrac{\big(16 \mu_n^2 - 13\big)\big(16\mu_n^2 - 9\big)\big(16\mu_n^2 - 1\big)}{48\mu_n} \vspace{1mm}\\ 64\mu_n\big(16\mu_n^2 - 1\big) & - \big(16\mu_n^2 - 9\big)\big(16\mu_n^2 - 1\big)\end{bmatrix} + \mathcal{O}\big(\rho_0^{-4}\big) \Biggr) \rho_0^{\sigma_3/2}\nonumber\\
\phantom{\qquad=}{} \times \big(2\sqrt{\mu_n}\big)^{\mathbb{I}-\sigma_3}{\dfrac{\ee^{-\pi \ii \mu_n}}{\sqrt{2\pi}}} {}{} \begin{bmatrix} 1 & \ii\\ 1 &-\ii \end{bmatrix} \ee^{ \rho_0 \sigma_3}, \qquad \arg(\rho_0) \in \left( -\dfrac{\pi}{2}, \dfrac{3\pi}{2} \right).
\label{eq:H-expand-zero}
\end{gather}
For convenience, we introduce the notation
\begin{equation}
 \mathbf{H}_n(\diamond) := \sqrt{\dfrac{\pi}{4\mu_n}} \ee^{\pi \ii/4 }\ee^{\pi \ii \mu_n} \cdot \dfrac{\diamond}{2} \begin{bmatrix} H^{(1)}_{2\mu_n - 1}(\diamond) & H^{(2)}_{1 - 2\mu_n }(\diamond) \vspace{1mm}\\ H^{(1)}_{1 + 2\mu_n}(\diamond) & H^{(2)}_{-1-2\mu_n}(\diamond) \end{bmatrix},
 \label{eq:Bn-def}
\end{equation}
with a fixed determination of the square root; this choice of prefactor guarantees that we have~$\det (\mathbf{H}_n ) = 1$ identically. Using the identity \cite[equation~(10.11.4)]{DLMF}, we note that $\mathbf{H}_n$ satisfies
\begin{gather}
\label{Bn-jump}
\mathbf{H}_n(\ee^{\pi \ii}\diamond) = \mathbf{H}_n(\diamond) \begin{bmatrix}
0 & -1 \\ 1 & 2\cos(2\pi \mu_n)
\end{bmatrix}.
\end{gather}
We can now rewrite the jump condition \eqref{r-hat-jump} as
\begin{align*}
 &\hat{\mathbf{R}}^\mathrm{even/odd}_+(\lambda,z)= \hat{\mathbf{R}}^\mathrm{even/odd}_-(\lambda,z) \mathbf{H}_n(\rho_\infty) \begin{bmatrix} 1 & \ee^{2\pi \ii \mu_n} \\ -\ee^{2\pi \ii \mu_n} & -1 \end{bmatrix} \mathbf{V}^\mathrm{even/odd}\\
 & \hphantom{\hat{\mathbf{R}}^\mathrm{even/odd}_+(\lambda,z)=}{}
 \times \begin{bmatrix} \ee^{\pi \ii \mu_n}& \ee^{\pi \ii \mu_n}\\ - \ee^{3\pi \ii \mu_n} & - \ee^{-\pi \ii \mu_n} \end{bmatrix}^{-1} \mathbf{H}_n^{-1}\big(\ee^{-\pi \ii/2}\rho_0\big).
\end{align*}
Next, define
\begin{align}
 & \mathbf{\Omega}^\mathrm{even/odd} (\lambda, z) := \begin{bmatrix} 1 & \dfrac{1}{2} - \dfrac{\mu_n}{2} - \dfrac{3}{32\mu_n} \vspace{1.5mm}\\ -1 & \dfrac{1}{2} + \dfrac{\mu_n}{2} + \dfrac{3}{32\mu_n} \end{bmatrix}^{-1} \big(2\sqrt{\mu_n}\big)^{\sigma_3}\nonumber\\
 &\hphantom{\mathbf{\Omega}^\mathrm{even/odd} (\lambda, z) :=}{}
 \times\hat{\mathbf{R}}^\mathrm{even/odd}(\lambda, z) \begin{cases}
 \mathbf{H}_n(\rho_\infty), & |\lambda| >1, \\
 \mathbf{H}_n\big(\ee^{-\pi \ii/2}\rho_0\big), & |\lambda| < 1.
 \end{cases}
 \label{eq:Omega-def}
\end{align}
Then, $\mathbf{\Omega}^\mathrm{even/odd}$ satisfies
\begin{align}
 &\mathbf{\Omega}^\mathrm{even/odd}_+(\lambda, z) = \mathbf{\Omega}_-^\mathrm{even/odd}(\lambda, z) \begin{bmatrix} 1 & \ee^{2\pi \ii \mu_n} \\ -\ee^{2\pi \ii \mu_n} & -1 \end{bmatrix} \cdot \mathbf{V}^\mathrm{even/odd} \nonumber\\
 &\hphantom{\mathbf{\Omega}^\mathrm{even/odd}_+(\lambda, z) =}{}
 \times\begin{bmatrix} \ee^{\pi \ii \mu_n}& \ee^{\pi \ii \mu_n}\\ - \ee^{3\pi \ii \mu_n} & - \ee^{-\pi \ii \mu_n} \end{bmatrix}^{-1}, \qquad |\lambda| = 1,
 \label{eq:Omega-jump-circle}
\end{align}
where the jump depends only on the parity of $n$. Furthermore, since $\rho_\infty$ and $\rho_0$ change signs across the negative imaginary axis, we may use \eqref{Bn-jump} to find
\begin{equation}
 \mathbf{\Omega}_+^\mathrm{even/odd} (\lambda, z) = \mathbf{\Omega}_-^\mathrm{even/odd} (\lambda, z) \begin{bmatrix}
0 & -1 \\ 1 & 2\cos(2\pi \mu_n)
\end{bmatrix},
\end{equation}
for $\lambda$ on the negative imaginary axis with $|\lambda| > 1$, oriented towards the origin and
\begin{equation}
 \mathbf{\Omega}_+^\mathrm{even/odd} (\lambda, z) = \mathbf{\Omega}_-^\mathrm{even/odd} (\lambda, z) \begin{bmatrix}
0 & -1 \\ 1 & 2\cos(2\pi \mu_n)
\end{bmatrix},
\label{eq:Omega-jump-outside-circle}
\end{equation}
for $\lambda$ on the negative imaginary axis with $|\lambda| <1$, oriented away from the origin.

It follows from Riemann--Hilbert Problem \ref{rhp:Rhat-even/odd}, \eqref{eq:Omega-def}, and \eqref{h-infty-asymptotics} that $\mathbf{\Omega}^\mathrm{even/odd}$ has the following asymptotic behavior as $\lambda \to \infty$:
\begin{align}
\mathbf{\Omega}^\mathrm{even/odd}(\lambda, z) = \big( \mathbb{I} + \mathbf{\Xi^{\mathrm{even/odd}}}(z) \lambda^{-1} + \mathcal{O}\big(\lambda^{-2}\big) \big) \rho_\infty^{\sigma_3/2} {\dfrac{1}{\sqrt{2}}} \begin{bmatrix} \ii & -1\\ 1 &-\ii \end{bmatrix} \ee^{\ii \rho_\infty \sigma_3},
 \label{eq:Omega-expand-infty}
\end{align}
where the $\mathcal{O}\big(\lambda^{-2}\big)$ represents an asymptotic series that is differentiable term-by-term with respect to both $\lambda$ and $z$. Analogously, we have
\begin{align}
 &\mathbf{\Omega}^{\mathrm{even/odd}}(\lambda,z)= \mathbf{\Delta}^\mathrm{even/odd}(z)\big(\mathbb{I}+\mathbf{\Pi}^\mathrm{even/odd}(z)\lambda+\mathcal{O}\big(\lambda^2\big)\big)\rho_0^{\sigma_3/2}\ee^{-\pi \ii \sigma_3/4} \nonumber\\
 &\hphantom{\mathbf{\Omega}^{\mathrm{even/odd}}(\lambda,z)=}{}
 \times{\dfrac{1}{\sqrt{2}}} {}{} \begin{bmatrix} \ii & -1 \\ 1 &-\ii \end{bmatrix} \ee^{ \rho_0 \sigma_3}, \qquad \lambda\to 0,
 \label{eq:Omega-expand-zero}
\end{align}
where $\mathcal{O}\big(\lambda^2\big)$ represents an asymptotic series at the origin $\lambda=0$ which is similarly term-by-term differentiable.
Notice that we can now relate the limiting formula from Lemma \ref{lemma:limit} to \smash{$\mathbf{\Omega}^{\mathrm{even/odd}}$} using definitions \eqref{eq:Omega-def} and \eqref{eq:Bn-def} to find that
\begin{gather*}
 \hat{R}^\mathrm{even/odd}_{11}(0,z)+\hat{R}^\mathrm{even/odd}_{21}(0,z)-\hat{R}^\mathrm{even/odd}_{12}(0,z)-\hat{R}^\mathrm{even/odd}_{22}(0,z)\\
 \qquad =\sqrt{\frac{\pi}{4\mu_n}} \ee^{\pi \ii/4} \ee^{\pi \ii \mu_n} \frac{\ee^{-\pi \ii/2}\rho_0}{2} \big[ \big(H^{(2)}_{-1-2\mu_n}\big(\ee^{-\pi \ii/2}\rho_0\big) + H^{(2)}_{1-2\mu_n}\big(\ee^{-\pi \ii/2}\rho_0\big)\big)\Omega^{\mathrm{even/odd}}_{21}(\lambda, z) \\
 \phantom{\qquad =}{}- \big(H^{(1)}_{2\mu_n-1}\big(\ee^{-\pi \ii/2}\rho_0\big) + H^{(1)}_{2\mu_n +1}\big(\ee^{-\pi \ii/2}\rho_0\big)\big)\Omega^{\mathrm{even/odd}}_{22}(\lambda, z) \big]_{\lambda=0}.
\end{gather*}
Then, using \eqref{eq:Omega-expand-zero} and \eqref{eq:H-expand-zero} yields
\begin{gather*}
 \mathop{\lim_{n\to\infty}}_{\text{$n$ even/odd}}u_n\big(n^{-1}z\big)=U^\mathrm{even/odd}(z):=-\frac{1}{2z\Delta^\mathrm{even/odd}_{21}(z)^2}.
 % \label{eq:un-U}
\end{gather*}

To extract the monodromy parameters of $U(z)$ from $\mathbf{\Omega}(\lambda, z)$, we notice that it solves Riemann--Hilbert Problem \ref{rhp:D8} with
\begin{equation}\label{eq:limiting-stokes-data}
t_1^\infty = t_0^0 = -2\cos(2\pi \mu_n)= -\left(e_1^2+\frac{1}{e_1^2}\right),
\end{equation}
and
\begin{equation}
\mathbf{C}_{0\infty} = \begin{bmatrix} 1 & \ee^{2\pi \ii \mu_n} \\ -\ee^{2\pi \ii \mu_n} & -1 \end{bmatrix} \cdot \mathbf{V}^\mathrm{even/odd} \cdot \begin{bmatrix} \ee^{\pi \ii \mu_n}& \ee^{\pi \ii \mu_n}\\ - \ee^{3\pi \ii \mu_n} & - \ee^{-\pi \ii \mu_n} \end{bmatrix}^{-1}.
\label{eq:limiting-connection-matrix}
\end{equation}
Since $\mu_n$ and $\mathbf{V}^\mathrm{even/odd}$ depend only on the parity of $n$ (see \eqref{eq:mu-n}, \eqref{eq:Vm-even}, respectively), and ${e_1 = e_{1, n} = \ee^{\pi \ii \mu_n}}$ (see Remark \ref{rem:n-dependent-notation}), it follows that \eqref{eq:limiting-stokes-data}--\eqref{eq:limiting-connection-matrix} depend only on the parity of $n$. Recalling the formul\ae\ for $y_i$ in Section \ref{sec:monodromy-rep-$D_8$-manifold}, one immediately arrives at formul\ae\ \eqref{eq:y1}--\eqref{eq:y3}.

